\documentclass[11pt, a4paper]{article}

\usepackage[english]{estilo-mate}
\tema{Measure Theory}

\hypersetup{pdftitle={Measure Theory Notes}}

\title{
    \vspace{-2cm}
    \textbf{\Huge Measure Theory Notes}\\
    \vspace{0.5cm}
}
\author{Lucas Allara}
\date{}

\begin{document}

\maketitle
\newpage

\tableofcontents
\newpage

\section{Motivation: from Riemann to Lebesgue}

Let $f: [a, b] \to \mathbb{R}$ be bounded. The Riemann integral is constructed by dividing the domain
$[a, b]$ into finitely many subintervals $[t_{k-1}, t_k]$ and approximating the area under the
curve
with sums
$$ \sum_{k=1}^n f(\xi_k) (t_k - t_{k-1}), \qquad \xi_k \in [t_{k-1}, t_k]. $$
When these sums converge as the partition is refined, $f$ is said to be Riemann-integrable.

This construction has two limitations:
\begin{itemize}
    \item Functions with many discontinuities are not Riemann-integrable, even though
    intuitively they should have an area. The classical example is
    $\mathbf{1}_{\mathbb{Q} \cap [0,1]}$.
    \item The theorems on interchanging limit and integral require strong hypotheses.
\end{itemize}

An earlier attempt (Jordan, 1892) was to measure a bounded set $E$ by approximating it from outside
with \emph{finite} covers by intervals:
$$ \bar c(E) = \inf\left\{ \sum_{k=1}^n \ell(I_k) : E \subseteq I_1 \cup \cdots \cup I_n \right\}, $$
where $\ell$ denotes length. This \emph{Jordan content} fails on dense sets:
every finite cover of $\mathbb{Q} \cap [0,1]$ by open intervals is an open set that
contains a dense subset of $[0,1]$, and its closure (which only adds finitely many endpoints) contains
$[0,1]$; the sum of the lengths is then $\ge 1$, so that
$\bar c(\mathbb{Q} \cap [0,1]) = 1$ and, by the same argument,
$\bar c(\mathbb{Q}^c \cap [0,1]) = 1$: two disjoint sets whose union is $[0,1]$ would have
``measures'' adding up to $2$. The decisive correction was to admit \emph{countable} covers:
with them $\mathbb{Q} \cap [0,1]$ is covered by intervals of arbitrarily small total length
(an interval of length $\epsilon/2^n$ around the $n$-th rational), and its measure
turns out to be $0$.

\textbf{Lebesgue's idea (1904).} Instead of partitioning the domain, the range is partitioned. Take a
partition $a_0 < a_1 < \cdots < a_n$ of the range of $f$, and for each $k$ consider the set
$$ A_k = \{ t \in [a, b] : a_{k-1} < f(t) \le a_k \}. $$
The integral is approximated by
$$ \sum_{k=1}^n a_k \cdot m(A_k), $$
where $m(A_k)$ is a notion of \emph{measure} of the set $A_k$.

The goal is to construct, starting from a reasonable class of subsets called
\emph{measurable}, a function $m: \mathcal{A} \to [0, \infty]$ with properties such as:
\begin{itemize}
    \item $A \subseteq B \implies m(A) \le m(B)$.
    \item $m(A \cup B) = m(A) + m(B)$ when $A \cap B = \emptyset$.
    \item On $\mathbb{R}$, $m([a, b]) = b - a$ and $m$ is translation invariant.
\end{itemize}
It will be seen that \emph{not} every subset of $\mathbb{R}$ admits such a measure; there are subsets
 that are not measurable.

\newpage

\section{\texorpdfstring{$\sigma$}{σ}-algebras}

\begin{definicion}{$\sigma$-algebra}{sigma_algebra}
Let $X$ be a set. A family $\mathcal{A} \subseteq \mathcal{P}(X)$ is a $\sigma$-algebra
on $X$ if:
\begin{enumerate}
    \item $\emptyset \in \mathcal{A}$.
    \item Closed under complements: $A \in \mathcal{A} \implies A^c \in \mathcal{A}$.
    \item Closed under countable unions: $(A_n) \subset \mathcal{A} \implies
    \bigcup_{n=1}^\infty A_n \in \mathcal{A}$.
\end{enumerate}
The elements of $\mathcal{A}$ are called \emph{measurable sets} (with respect to $\mathcal{A}$).
\end{definicion}

\obs By De Morgan, a $\sigma$-algebra is also closed under countable intersections: if
$(A_n) \subset \mathcal{A}$, then
$$ \bigcap_n A_n = \left( \bigcup_n A_n^c \right)^c \in \mathcal{A}. $$
Moreover, $X = \emptyset^c \in \mathcal{A}$.

\subsection{Examples}

\begin{ejemplo}{Trivial $\sigma$-algebra}{}
$\mathcal{A} = \{\emptyset, X\}$ is the smallest possible $\sigma$-algebra on $X$. It satisfies
the three axioms by inspection.
\end{ejemplo}

\begin{ejemplo}{Discrete $\sigma$-algebra}{}
$\mathcal{A} = \mathcal{P}(X)$ is the largest $\sigma$-algebra, where every subset is
measurable. On small sets this is the natural choice.
\end{ejemplo}

\begin{ejemplo}{Cofinite sets: an algebra that is not a $\sigma$-algebra}{}
Let $X$ be infinite and $\mathcal{A} = \{A \subseteq X : A \text{ finite or } A^c \text{ finite}\}$.
This family is an \emph{algebra} (closed under finite unions and intersections and under
complements) but it is \emph{not} a $\sigma$-algebra.
\end{ejemplo}

\begin{ejemplo}{$\sigma$-algebra of coin sequences}{ej_monedas}
Let $X = \{0, 1\}^{\mathbb{N}}$ be the space of infinite sequences of zeros and ones (the
results of tossing a coin infinitely many times, with $1$ = heads). For each $n$ define the event
$$ E_n = \{ x \in X : x_n = 1 \} \qquad (\text{``the } n\text{-th toss is heads''}), $$
and take $\mathcal{A} = \sigma(E_1, E_2, \dots)$. This $\sigma$-algebra contains all the
events described by countably many tosses: for example $\bigcap_n E_n$
(``always heads''), $E_1 \cap E_2^c$ (``heads and then tails''), or the limits of the
next subsection. This space will reappear when studying measures and the
Borel-Cantelli lemmas.
\end{ejemplo}

\begin{proposicion}{Intersection of $\sigma$-algebras}{interseccion_sigmas}
If $(\mathcal{A}_i)_{i \in I}$ is an arbitrary family of
$\sigma$-algebras on a set $X$, then the intersection
$$ \mathcal{B} = \bigcap_{i \in I} \mathcal{A}_i $$
is again a $\sigma$-algebra on $X$.
\end{proposicion}

\begin{proof}\leavevmode\\
The three axioms that define a $\sigma$-algebra are verified.

\textit{It contains the empty set.} For each $i \in I$, $\mathcal{A}_i$ is a $\sigma$-algebra and
therefore $\emptyset \in \mathcal{A}_i$. Since this holds for every $i$, $\emptyset \in \mathcal{B}$.

\textit{Closed under complements.} Let $A \in \mathcal{B}$. For each $i \in I$, $A \in
\mathcal{A}_i$, and since $\mathcal{A}_i$ is closed under complements, $A^c \in \mathcal{A}_i$.
Therefore $A^c$ belongs to all the $\mathcal{A}_i$, that is, $A^c \in \mathcal{B}$.

\textit{Closed under countable unions.} Let $(A_n)_{n \in \mathbb{N}} \subseteq \mathcal{B}$ be a
sequence. Fix $i \in I$. For each $n$, $A_n \in \mathcal{B} \subseteq \mathcal{A}_i$, so
that $(A_n) \subseteq \mathcal{A}_i$. Since $\mathcal{A}_i$ is closed under countable unions,
$\bigcup_{n=1}^{\infty} A_n \in \mathcal{A}_i$. This holds for every $i \in I$, which implies
$\bigcup_n A_n \in \mathcal{B}$.
\end{proof}

\obs This proposition is what allows defining the \emph{$\sigma$-algebra generated} by a
family: given a set $\mathcal{F} \subseteq \mathcal{P}(X)$ there is always at least one
$\sigma$-algebra containing $\mathcal{F}$ (for example $\mathcal{P}(X)$), and the intersection
of all of them is the smallest $\sigma$-algebra with that property.

\subsection{\texorpdfstring{$\sigma$}{σ}-algebra generated}

\begin{definicion}{Generated $\sigma$-algebra}{sigma_generada}
Given a family $\mathcal{F} \subseteq \mathcal{P}(X)$, the $\sigma$-algebra generated by
$\mathcal{F}$ is
$$ \sigma(\mathcal{F}) = \bigcap \{ \mathcal{A} : \mathcal{A} \text{ is a }\sigma\text{-algebra and }
\mathcal{F} \subseteq \mathcal{A} \}. $$
It is the smallest $\sigma$-algebra containing $\mathcal{F}$.
\end{definicion}

\subsection{Structure on countable sets}

On a countable set, every $\sigma$-algebra comes from a partition, whose elements
are called \emph{atoms}: the smallest nonempty measurable sets containing a given point.

\begin{teorema}{$\sigma$-algebras on countable sets}{atomos_sigma}
Let $X$ be countable and $\Sigma$ a $\sigma$-algebra on $X$. For each $x \in X$ define
$$ A_x = \bigcap_{\substack{B \in \Sigma \\ x \in B}} B. $$
The family $\mathcal{A} = \{A_x : x \in X\}$ is a partition of $X$ by elements of $\Sigma$,
and $\Sigma = \sigma(\mathcal{A})$.
\end{teorema}

\begin{proof}
Each $B$ in the intersection contains $x$, so that $x \in A_x$. In particular the $A_x$
cover $X$.

To see that $A_x \in \Sigma$, the intersection is reformulated in terms of the complement, which is
at most countable. For each $y \in A_x^c$, since $y$ does not belong to the intersection,
there is some $B_y \in \Sigma$ with $x \in B_y$ but $y \notin B_y$. Define
$$ A_x' = \bigcap_{y \in A_x^c} B_y, $$
which is an at most countable intersection of elements of $\Sigma$ and therefore belongs to
$\Sigma$. Each $B_y$ is among the sets that define $A_x$, which gives $A_x \subseteq B_y$
for every $y \in A_x^c$, and then $A_x \subseteq A_x'$. Conversely, if $y \in A_x^c$
then $y \notin B_y$ and therefore $y \notin A_x'$, which shows $A_x^c \subseteq (A_x')^c$,
that is $A_x' \subseteq A_x$. It follows that $A_x = A_x' \in \Sigma$.

Now a minimality property is shown: if $B \in \Sigma$ is nonempty and $B \subseteq A_x$,
then $B = A_x$. Suppose first that $x \notin B$. The complement $B^c$ belongs to $\Sigma$
and contains $x$, so that it is among the sets that define $A_x$; then $A_x \subseteq
B^c$, equivalent to $A_x \cap B = \emptyset$. But $B \subseteq A_x$ nonempty forces $A_x \cap B
= B \ne \emptyset$, a contradiction. Thus $x \in B$, and since $B$ is a measurable set containing $x$,
$A_x \subseteq B$. Together with the hypothesis $B \subseteq A_x$, one gets $B = A_x$.

The partition property follows from this. If $A_x \cap A_y$ is nonempty, that intersection
is in $\Sigma$, is nonempty and is contained in $A_x$, so that minimality forces
$A_x \cap A_y = A_x$. Consequently $A_x \subseteq A_y$, and exchanging $x$ with $y$ one
obtains $A_y \subseteq A_x$, that is $A_x = A_y$. Combined with the fact that the $A_x$ cover $X$, this
proves that $\mathcal{A}$ is a partition.

It remains to see $\Sigma = \sigma(\mathcal{A})$. The inclusion $\sigma(\mathcal{A}) \subseteq \Sigma$
is a consequence of each $A_x \in \Sigma$. For the other inclusion, given $B \in \Sigma$
it is claimed that $B = \bigcup_{x \in B} A_x$. The inclusion $\subseteq$ follows from $x \in A_x$ for each $x
\in B$. For the inclusion $\supseteq$, if $x \in B$ then $B$ is one of the sets that
define $A_x$, which gives $A_x \subseteq B$. The union is countable because $B \subseteq X$ is,
and each $A_x \in \sigma(\mathcal{A})$, so that $B \in \sigma(\mathcal{A})$.
\end{proof}

\obs The $\sigma$-algebras on countable sets are then always \emph{discrete}: they
are generated by a partition.

\begin{definicion}{Borel $\sigma$-algebra}{borel}
The Borel $\sigma$-algebra on $\mathbb{R}$ is the one generated by the open sets:
$$ \mathcal{B}(\mathbb{R}) = \sigma(\{U \subseteq \mathbb{R} : U \text{ open}\}). $$
Its elements are called \emph{Borel sets}.
\end{definicion}

\subsection{Borel from intervals}

Although the Borel $\sigma$-algebra is defined from \emph{all} open sets, on
$\mathbb{R}$ a much smaller family suffices: the open intervals $(a, b)$. Something more is proved:
every open subset of $\mathbb{R}$ decomposes as an at most countable \emph{disjoint} union of
open intervals. These intervals are the connected components of the open set.

\begin{proposicion}{Structure of open sets in $\mathbb{R}$}{abiertos_R}
Every open $A \subseteq \mathbb{R}$ is written as an at most countable union of disjoint open
intervals. In particular,
$$ \mathcal{B}(\mathbb{R}) = \sigma(\{(a, b) : a < b\}). $$
\end{proposicion}

\begin{proof}
The strategy is to assign to each point $a \in A$ the largest open interval contained in $A$ that
contains it. These intervals cover $A$ and, by their maximality, are pairwise disjoint. Since each
one contains a rational, the number of distinct intervals turns out to be at most countable.

For each $a \in A$ define
$$ x(a) = \sup\{x \in \mathbb{R} \setminus A : x < a\}, \qquad y(a) = \inf\{y \in \mathbb{R}
\setminus A : y > a\}, $$
with the convention $\sup \emptyset = -\infty$, $\inf \emptyset = +\infty$, and define
$I_a = (x(a), y(a))$.

Since $A$ is open and $a \in A$, there is $\delta > 0$ with
$(a - \delta, a + \delta) \subseteq A$. Since no point of the neighborhood $(a - \delta, a)$ is in
$\mathbb{R} \setminus A$, every point in the set whose supremum is $x(a)$ satisfies
$x \le a - \delta$. Therefore $x(a) \le a - \delta < a$, and likewise $y(a) > a$, so that
$a \in I_a$.

For the inclusion $I_a \subseteq A$, let $t \in I_a$ with $t \ne a$, and suppose by contradiction
$t \notin A$. If $t < a$, then $t$ belongs to the set $\{x \in \mathbb{R} \setminus A :
x < a\}$ whose supremum is $x(a)$, which forces $t \le x(a)$ and contradicts $t > x(a)$. The case
$t > a$ is symmetric with respect to $y(a)$.

As a consequence, $A = \bigcup_{a \in A} I_a$: the inclusion $\supseteq$ is $I_a \subseteq A$, and
the $\subseteq$ one is $a \in I_a$.

Suppose $I_a \cap I_b \ne \emptyset$ and let $z \in I_a \cap I_b$. By definition of
$I_a = (x(a), y(a))$ and $I_b = (x(b), y(b))$, the inequalities
$$ x(a) < z < y(a) \qquad \text{and} \qquad x(b) < z < y(b) $$
hold.

First $y(a) = y(b)$ is proved by ruling out the possibility $y(a) < y(b)$ (the case $y(a) > y(b)$
is symmetric, exchanging the roles of $a$ and $b$). Indeed, suppose $y(a) < y(b)$. By
definition of $y(a)$ as the infimum of $\{y \in \mathbb{R} \setminus A : y > a\}$,
$y(a) \notin A$. By the inequalities $z < y(a) < y(b)$, one has $y(a) \in (z, y(b))$. But
$(z, y(b)) \subseteq I_b \subseteq A$, whence $y(a) \in A$, contradicting $y(a) \notin A$.
It follows that $y(a) = y(b)$.

An analogous argument on the other side proves $x(a) = x(b)$: if $x(a) > x(b)$, by definition
of $x(a)$ as the supremum of $\{x \in \mathbb{R} \setminus A : x < a\}$, $x(a) \notin A$. By the
inequalities $x(b) < x(a) < z$, one has $x(a) \in (x(b), z)$. But $(x(b), z) \subseteq
I_b \subseteq A$, whence $x(a) \in A$, contradicting $x(a) \notin A$. The case $x(a) < x(b)$ is
symmetric. It follows that $x(a) = x(b)$.

With $x(a) = x(b)$ and $y(a) = y(b)$ it follows that $I_a = I_b$.

Each $I_a$ is a nonempty open subset of $\mathbb{R}$, so that by
density of $\mathbb{Q}$ it contains at least one rational. Choosing one rational per distinct
interval yields an injective map from the family of distinct intervals into
$\mathbb{Q}$ (it is injective because two distinct intervals are disjoint). Consequently, the
family injects into $\mathbb{Q}$ and is at most countable.

Let $\mathcal{I} = \{(a, b)
: a < b\}$. For one inclusion, each $(a, b)$ is open, which gives $\mathcal{I} \subseteq
\mathcal{B}(\mathbb{R})$ and, by minimality, $\sigma(\mathcal{I}) \subseteq
\mathcal{B}(\mathbb{R})$.
For the other inclusion, every open set is written as a countable union of elements of
$\mathcal{I}$ by what was proved, and therefore belongs to $\sigma(\mathcal{I})$. Thus,
$\sigma(\mathcal{I})$ is a $\sigma$-algebra containing the open sets, and by minimality of
$\mathcal{B}(\mathbb{R})$, $\mathcal{B}(\mathbb{R}) \subseteq \sigma(\mathcal{I})$.
\end{proof}

\begin{corolario}{Sets in $\mathcal{B}(\mathbb{R})$}{borel_contiene}
$\mathcal{B}(\mathbb{R})$ contains the open sets, the closed sets, the points, the half-open
intervals and the half-lines.
\end{corolario}
 
\subsection{Operations with sequences}

\begin{definicion}{Limit superior and inferior of sets}{limsup_conjuntos}
Given $(A_n) \subset \mathcal{A}$:
\begin{align*}
\varlimsup_n A_n &= \bigcap_{m=1}^\infty \bigcup_{k \ge m} A_k = \{x : x \in A_n \text{ for
infinitely many } n\}, \\
\varliminf_n A_n &= \bigcup_{m=1}^\infty \bigcap_{k \ge m} A_k = \{x : x \in A_n \text{ from
some } n \text{ on}\}.
\end{align*}
\end{definicion}

\obs If $\mathcal{A}$ is a $\sigma$-algebra and $(A_n) \subset \mathcal{A}$, then
$\varlimsup A_n, \varliminf A_n \in \mathcal{A}$.

\begin{proposicion}{Properties of the limit superior and inferior}{props_limsup_liminf}
Let $(A_n)$ be a sequence of subsets of $X$. Then:
\begin{enumerate}
    \item $\varliminf_n A_n \subseteq \varlimsup_n A_n$.
    \item $\left(\varliminf_n A_n\right)^c = \varlimsup_n A_n^c$ and
    $\left(\varlimsup_n A_n\right)^c = \varliminf_n A_n^c$.
    \item If $A \subseteq A_n \subseteq B$ for every $n$, then
    $A \subseteq \varliminf_n A_n \subseteq \varlimsup_n A_n \subseteq B$.
    \item If $(A_n)$ is increasing, $\varliminf_n A_n = \varlimsup_n A_n = \bigcup_n A_n$.
    \item If $(A_n)$ is decreasing, $\varliminf_n A_n = \varlimsup_n A_n = \bigcap_n A_n$.
\end{enumerate}
\end{proposicion}

\begin{proof}\leavevmode\\
(i) If $x \in \varliminf_n A_n = \bigcup_m \bigcap_{k \ge m} A_k$, there is $m_0$ such that $x \in
A_k$ for every $k \ge m_0$; that is, $x$ is in all the $A_k$ except finitely many. In particular,
$x$ is in infinitely many $A_k$, so that for every $m$ there is $k \ge m$ with $x \in A_k$. This
means $x \in \bigcap_m \bigcup_{k \ge m} A_k = \varlimsup_n A_n$.

(ii) By De Morgan's laws applied twice,
$$ \left(\varliminf_n A_n\right)^c = \left(\bigcup_m \bigcap_{k \ge m} A_k\right)^c =
\bigcap_m \left(\bigcap_{k \ge m} A_k\right)^c = \bigcap_m \bigcup_{k \ge m} A_k^c =
\varlimsup_n A_n^c. $$
The other equality is proved analogously, exchanging $\bigcup$ and $\bigcap$.

(iii) If $x \in A$, then $x \in A_n$ for every $n$, in particular $x \in \bigcap_{k \ge m}
A_k$ for every $m$. Hence $x \in \bigcup_m \bigcap_{k \ge m} A_k = \varliminf_n A_n$, which
gives $A \subseteq \varliminf_n A_n$. Combined with (i), $A \subseteq \varliminf_n A_n \subseteq
\varlimsup_n A_n$.

For the upper bound, $A_n \subseteq B$ for every $n$ is equivalent to $B^c \subseteq A_n^c$ for every
$n$. Applying the inclusion $B^c \subseteq \varliminf_n A_n^c$ (just proved with $B^c$ in
place of $A$ and $A_n^c$ in place of $A_n$) and then (ii),
$$ B^c \subseteq \varliminf_n A_n^c = \left(\varlimsup_n A_n\right)^c, $$
whence, taking complements, $\varlimsup_n A_n \subseteq B$.

(iv) If $(A_n)$ is increasing, $\bigcap_{k \ge m} A_k = A_m$ for every $m$, so that
$$ \varliminf_n A_n = \bigcup_m \bigcap_{k \ge m} A_k = \bigcup_m A_m. $$
On the other hand, $\bigcup_{k \ge m} A_k = \bigcup_k A_k$ for every $m$ (the union does not change when
starting later in an increasing sequence), so that
$$ \varlimsup_n A_n = \bigcap_m \bigcup_{k \ge m} A_k = \bigcup_n A_n. $$

(v) If $(A_n)$ is decreasing, $(A_n^c)$ is increasing. Applying (iv) to $(A_n^c)$ and then (ii),
$$ \left(\varlimsup_n A_n\right)^c = \varliminf_n A_n^c = \bigcup_n A_n^c =
\left(\bigcap_n A_n\right)^c, $$
and analogously for $\varliminf$. Taking complements, $\varliminf_n A_n = \varlimsup_n A_n
= \bigcap_n A_n$.
\end{proof}

\begin{ejemplo}{Computing limits of sets}{ej_limites_conjuntos}
\begin{enumerate}
    \item Let $A_n = [0, 2]$ if $n$ is odd and $A_n = [1, 4]$ if $n$ is even. A point is in
    infinitely many $A_n$ if and only if it belongs to $[0,2] \cup [1,4] = [0,4]$, and it is in all but
    finitely many if and only if it belongs to $[0,2] \cap [1,4] = [1,2]$. Thus
    $$ \varliminf_n A_n = [1, 2], \qquad \varlimsup_n A_n = [0, 4], $$
    and the sequence does not converge as sets.
    \item Let $A_n = [1/n, 1)$. The sequence is increasing, and by Proposition
    \ref{prop:props_limsup_liminf},
    $$ \varliminf_n A_n = \varlimsup_n A_n = \bigcup_n [1/n, 1) = (0, 1). $$
    The sequence converges to the set $(0,1)$; observe that the limit does not contain $0$
    even though each $A_n$ is ``attached'' to it.
\end{enumerate}
\end{ejemplo}

\obs In the space $X = \{0,1\}^{\mathbb{N}}$ of coin tosses with the events
$E_n = \{x : x_n = 1\}$, the limits have a direct probabilistic reading:
$$ \varlimsup_n E_n = \{\text{heads infinitely often}\}, \qquad
\varliminf_n E_n = \{\text{eventually always heads}\}. $$
The inclusion $\varliminf \subseteq \varlimsup$ is, with this reading, transparent: if from
some moment on heads always comes up, in particular heads comes up infinitely often.

\newpage

\section{Non-measurable sets: Vitali's example}

Before moving on to the definition of measure and to Lebesgue's construction, it is convenient to understand
why one cannot expect \emph{every} subset of $\mathbb{R}$ to admit a measure with the
desired properties (translation invariance and $\sigma$-additivity). The
classical example showing that this is impossible is \emph{Vitali's} construction: a subset
$V \subseteq [0, 1]$ that cannot be measurable with respect to any translation-invariant
measure extending the ``length'' $m([0, 1]) = 1$.

\begin{teorema}{Existence of non-measurable sets}{vitali}
There is no function $m: \mathcal{P}(\mathbb{R}) \to [0, \infty]$ that simultaneously satisfies:
\begin{enumerate}
    \item $m$ is $\sigma$-additive: if $(A_n)$ is disjoint, $m(\bigcup_n A_n) = \sum_n m(A_n)$.
    \item $m$ is translation invariant: $m(A + t) = m(A)$ for every $A \subseteq \mathbb{R}$
    and every $t \in \mathbb{R}$.
    \item $m([0, 1]) = 1$.
\end{enumerate}
\end{teorema}

\begin{proof}\leavevmode\\
Assuming the axiom of choice, a $V \subseteq [0, 1]$ is constructed that will lead to a
contradiction under (i)--(iii).

\textit{Construction.} On $[0, 1]$ define the equivalence relation
$$ x \sim y \ \Leftrightarrow \ x - y \in \mathbb{Q}. $$
Each class $[x]$ is dense in $[0, 1]$ (because $\mathbb{Q}$ is dense). By the axiom of choice,
choose one representative for each class and form
$$ V = \{\text{one representative of each class}\} \subseteq [0, 1]. $$

\textit{Translates.} Enumerate $\mathbb{Q} \cap [-1, 1] = \{r_1, r_2, \dots\}$, which is countable.
For each $k$ define
$$ V_k = V + r_k = \{v + r_k : v \in V\}. $$

\textit{The $V_k$ are pairwise disjoint.} If $V_k \cap V_j \ne \emptyset$, there are $v_1, v_2 \in
V$ with $v_1 + r_k = v_2 + r_j$, whence $v_1 - v_2 = r_j - r_k \in \mathbb{Q}$, that is
$v_1 \sim v_2$. Since $V$ contains a single representative per class, $v_1 = v_2$, and then
$r_k = r_j$, that is $k = j$.

\textit{$[0, 1] \subseteq \bigcup_k V_k \subseteq [-1, 2]$.} The inclusion on the right is a
consequence of the fact that $V \subseteq [0, 1]$ and $r_k \in [-1, 1]$ give $V_k \subseteq [-1, 2]$. For the
inclusion on the left, given $x \in [0, 1]$, its class $[x]$ has a representative $v \in V$,
and $x - v \in \mathbb{Q} \cap [-1, 1]$ (because $x, v \in [0, 1]$), so that $x - v = r_k$ for
some $k$, that is $x = v + r_k \in V_k$.

\textit{Contradiction.} Suppose that $V$ is measurable. By (ii), $m(V_k) = m(V)$ for every $k$,
and by (i),
$$ m\!\left(\bigcup_k V_k\right) = \sum_{k=1}^\infty m(V_k) = \sum_{k=1}^\infty m(V). $$
Cases are distinguished according to the value of $m(V) \in [0, \infty]$:
\begin{itemize}
    \item If $m(V) = 0$, the sum gives $0$, and by monotonicity $m([0, 1]) \le m(\bigcup_k V_k) = 0$.
    But $m([0, 1]) = 1$ by (iii), a contradiction.
    \item If $m(V) > 0$, $\sum_k m(V) = +\infty$, whereas by monotonicity $m(\bigcup_k V_k)
    \le m([-1, 2]) = 3$ (the latter by translation of $[0, 1]$ and additivity: $m([-1, 2]) =
    m([-1, 0]) + m([0, 1]) + m([1, 2]) = 3$). Contradiction.
\end{itemize}
In either case a contradiction is reached. Therefore $V$ cannot be measurable with respect to $m$, and
in particular there is no function $m: \mathcal{P}(\mathbb{R}) \to [0, \infty]$ satisfying
(i)--(iii) on all of $\mathcal{P}(\mathbb{R})$.
\end{proof}

\obs A translation-invariant, $\sigma$-additive and
normalized measure cannot be defined on \emph{the whole} $\sigma$-algebra $\mathcal{P}(\mathbb{R})$. The most that
can be achieved (with the axiom of choice) is a measure on a strictly
smaller $\sigma$-algebra: the Borel sets, or somewhat more: the Lebesgue-measurable sets.

\newpage

\section{Measures}

\begin{definicion}{Measure}{def_medida}
Let $(X, \mathcal{A})$ be a pair with $\mathcal{A}$ a $\sigma$-algebra on $X$. A function
$\mu: \mathcal{A} \to [0, \infty]$ is a \emph{measure} if:
\begin{enumerate}
    \item $\mu(\emptyset) = 0$.
    \item ($\sigma$-additivity) If $(A_n) \subset \mathcal{A}$ is a pairwise disjoint family,
    $$ \mu\!\left(\bigcup_{n=1}^\infty A_n\right) = \sum_{n=1}^\infty \mu(A_n). $$
\end{enumerate}
The triple $(X, \mathcal{A}, \mu)$ is called a \emph{measure space}.

If $\mu(X) < \infty$, $\mu$ is said to be \emph{finite}. If $X = \bigcup_n X_n$ with
$\mu(X_n) < \infty$ for every $n$, $\mu$ is said to be \emph{$\sigma$-finite}. If $\mu(X) = 1$,
$\mu$ is a \emph{probability}.
\end{definicion}

\subsection{Examples}

\begin{proposicion}{Dirac delta}{delta_dirac}
Let $X$ be a nonempty set and $x_0 \in X$ a fixed point. The function $\delta_{x_0}:
\mathcal{P}(X) \to [0, \infty]$ defined by
$$ \delta_{x_0}(A) = \begin{cases} 1, & x_0 \in A, \\ 0, & x_0 \notin A, \end{cases} $$
is a measure (called the \emph{Dirac delta} concentrated at $x_0$).
\end{proposicion}

\begin{proof}\leavevmode\\
\textit{$\delta_{x_0}(\emptyset) = 0$.} Since $x_0 \notin \emptyset$, $\delta_{x_0}(\emptyset) = 0$.

\textit{$\sigma$-additivity.} Let $(A_n)_{n \in \mathbb{N}} \subset \mathcal{P}(X)$ be a pairwise
disjoint family, and $A = \bigcup_{n} A_n$. Two cases are distinguished according to whether $x_0$ is in
$A$ or not.

\textit{Case 1: $x_0 \in A$.} There is $n_0 \in \mathbb{N}$ with $x_0 \in A_{n_0}$. Since the
$A_n$ are pairwise disjoint, that $n_0$ is unique: for every $n \ne n_0$, $x_0 \notin A_n$.
Therefore
$$ \sum_{n=1}^{\infty} \delta_{x_0}(A_n) = \delta_{x_0}(A_{n_0}) + \sum_{n \ne n_0}
\delta_{x_0}(A_n) = 1 + 0 = 1 = \delta_{x_0}(A). $$

\textit{Case 2: $x_0 \notin A$.} Then $x_0 \notin A_n$ for every $n$, so that
$\delta_{x_0}(A_n) = 0$ for every $n$, and
$$ \sum_{n=1}^{\infty} \delta_{x_0}(A_n) = 0 = \delta_{x_0}(A). $$
\end{proof}

\obs $\delta_{x_0}$ is a probability ($\delta_{x_0}(X) = 1$). It concentrates all the mass at the
point $x_0$.

\begin{proposicion}{Counting measure}{medida_conteo}
Let $X$ be a nonempty set. The function $\mu: \mathcal{P}(X) \to [0, \infty]$ defined by
$$ \mu(A) = \#A = \begin{cases} \text{number of elements of } A, & \text{if } A \text{ is finite,}
\\ \infty, & \text{if } A \text{ is infinite,} \end{cases} $$
is a measure (called the \emph{counting measure}).
\end{proposicion}

\begin{proof}\leavevmode\\
\textit{$\mu(\emptyset) = 0$.} The empty set has zero elements.

\textit{$\sigma$-additivity.} Let $(A_n)_{n \in \mathbb{N}}$ be pairwise disjoint and
$A = \bigcup_n A_n$. The structure of the proof depends on whether $A$ is finite or infinite.

\textit{Case 1: $A$ finite.} Since each $A_n \subseteq A$, all the $A_n$ are finite. Moreover,
since the $A_n$ are disjoint and $A$ is finite, only finitely many of them can be
nonempty; otherwise $A$ would contain at least one element for each nonempty $A_n$,
contradicting $\#A < \infty$. Renumbering, suppose $A_1, \dots, A_N$ nonempty and $A_n =
\emptyset$ for $n > N$. By elementary counting of finite disjoint unions,
$$ \#A = \#(A_1 \sqcup \cdots \sqcup A_N) = \sum_{n=1}^{N} \#A_n = \sum_{n=1}^{\infty} \#A_n, $$
where the last equality uses that $\#A_n = 0$ for $n > N$.

\textit{Case 2: $A$ infinite.} Then $\mu(A) = \infty$. Two subcases are distinguished:
\begin{itemize}
    \item If there is $n_0$ with $A_{n_0}$ infinite, then $\#A_{n_0} = \infty$, and
    $$ \sum_{n=1}^{\infty} \#A_n \ge \#A_{n_0} = \infty. $$
    \item If all the $A_n$ are finite, since $A$ is infinite there must be infinitely many nonempty
    $A_n$; otherwise $A$ would be a finite union of finite sets, and therefore finite. The series
    of positive terms
    $$ \sum_{n=1}^{\infty} \#A_n $$
    then has infinitely many terms $\ge 1$, and diverges to $\infty$.
\end{itemize}
In both subcases $\sum_n \#A_n = \infty = \#A$.
\end{proof}

\obs The counting measure is finite if and only if $X$ is finite, and it is $\sigma$-finite if and only if
$X$ is at most countable. On $X = \mathbb{N}$ with the counting measure, integrals
with respect to $\mu$ reduce to series.

\begin{ejemplo}{Bernoulli measure}{ej_medida_bernoulli}
On $X = \{0,1\}^{\mathbb{N}}$ with the $\sigma$-algebra generated by the events
$E_n = \{x : x_n = 1\}$ (Example \ref{ej:ej_monedas}), there is a unique probability $P$ that
models independent tosses of a fair coin: on the events that fix $k$
coordinates,
$$ P(\{x : x_{i_1} = \epsilon_1, \dots, x_{i_k} = \epsilon_k\}) = \frac{1}{2^k}. $$
With it the probability of \emph{heads at some point} can be computed. The event is
$\bigcup_n E_n$, which is not a disjoint union; it is made disjoint with
\begin{align*}
A_n = E_n \setminus (E_1 \cup \cdots \cup E_{n-1})
&= \{x : x_1 = \cdots = x_{n-1} = 0, \ x_n = 1\} \\
&= \{\text{the first heads is the } n\text{-th toss}\},
\end{align*}
which fixes $n$ coordinates, so that $P(A_n) = 1/2^n$. By $\sigma$-additivity,
$$ P\!\left(\bigcup_n E_n\right) = \sum_{n=1}^\infty P(A_n) = \sum_{n=1}^\infty \frac{1}{2^n}
= 1. $$
With probability $1$ some heads comes up.
\end{ejemplo}

\obs Of the previous examples, only the Dirac delta and the counting measure have been
\emph{constructed}. For the Bernoulli measure, and for the Lebesgue measure on $\mathbb{R}$
(which assigns to each interval its length), existence has not been proved so far: that is the
goal of the section on generating measures.

\subsection{Basic properties}

\begin{teorema}{Properties of a measure}{prop_medida}
Let $(X, \mathcal{A}, \mu)$ be a measure space. For $A, B, A_n \in \mathcal{A}$:
\begin{enumerate}
    \item If $E \subseteq F$, $\mu(E) \le \mu(F)$.
    \item If $E \subseteq F$ and $\mu(E) < \infty$, $\mu(F \setminus E) = \mu(F) - \mu(E)$.
    \item If $\mu(A \cap B) < \infty$,
    $$ \mu(A \cup B) = \mu(A) + \mu(B) - \mu(A \cap B). $$
    \item For any (not necessarily disjoint) family $(A_n)$,
    $$ \mu\!\left(\bigcup_n A_n\right) \le \sum_n \mu(A_n). $$
    \item If $A_n \nearrow A$ (that is, $A_n \subseteq A_{n+1}$ and $A = \bigcup_n A_n$),
    $$ \mu(A) = \lim_{n \to \infty} \mu(A_n). $$
    \item If $A_n \searrow A$ and $\mu(A_1) < \infty$,
    $$ \mu(A) = \lim_{n \to \infty} \mu(A_n). $$
    \item If $\varlimsup_n A_n = \varliminf_n A_n = A$ (that is, $(A_n)$ converges to $A$ as
    sets) and $\mu\!\left(\bigcup_n A_n\right) < \infty$,
    $$ \mu(A) = \lim_{n \to \infty} \mu(A_n). $$
\end{enumerate}
\end{teorema}

\begin{proof}\leavevmode\\
(i) $F = E \cup (F \setminus E)$ is a disjoint union, so that
$\mu(F) = \mu(E) + \mu(F \setminus E)
\ge \mu(E)$.

(ii) It follows from the previous step by solving for $\mu(F \setminus E)$.

(iii) Decompose $A \cup B$ into three disjoint sets: $A \cup B = (A \setminus B) \sqcup
(B \setminus A) \sqcup (A \cap B)$. By $\sigma$-additivity,
$$ \mu(A \cup B) = \mu(A \setminus B) + \mu(B \setminus A) + \mu(A \cap B). $$
Analogously, $A = (A \setminus B) \sqcup (A \cap B)$ and $B = (B \setminus A) \sqcup
(A \cap B)$, whence
$$ \mu(A) + \mu(B) = \mu(A \setminus B) + \mu(B \setminus A) + 2 \mu(A \cap B). $$
Subtracting $\mu(A \cap B)$ (finite by hypothesis, which avoids indeterminate forms $\infty - \infty$)
yields $\mu(A) + \mu(B) - \mu(A \cap B) = \mu(A \cup B)$.

(iv) Make the family disjoint: $B_1 = A_1$, $B_n = A_n \setminus (A_1 \cup \cdots \cup A_{n-1})$. The $B_n$ are
disjoint, $B_n \subseteq A_n$ (so $\mu(B_n) \le \mu(A_n)$), and $\bigcup B_n = \bigcup A_n$. Hence
$$ \mu\!\Big(\!\bigcup A_n\Big) = \sum \mu(B_n) \le \sum \mu(A_n). $$

(v) If $A_n \nearrow A$, make it disjoint with $B_1 = A_1$, $B_n = A_n \setminus A_{n-1}$.
Then $A_n = \bigsqcup_{k=1}^n B_k$ and $A = \bigsqcup_k B_k$, so that
$$ \mu(A) = \sum_{k=1}^\infty \mu(B_k) = \lim_n \sum_{k=1}^n \mu(B_k) = \lim_n \mu(A_n). $$

(vi) If $A_n \searrow A$, define $C_n = A_1 \setminus A_n$, which increases to $A_1 \setminus A$.
Applying (v) and (ii) (using $\mu(A_1) < \infty$),
$$ \mu(A_1) - \mu(A) = \mu(A_1 \setminus A) = \lim_n \mu(C_n) = \lim_n (\mu(A_1) - \mu(A_n)) =
\mu(A_1) - \lim_n \mu(A_n). $$
Cancelling $\mu(A_1)$, $\mu(A) = \lim \mu(A_n)$.

(vii) Define
$$ B_m = \bigcap_{k \ge m} A_k, \qquad C_m = \bigcup_{k \ge m} A_k. $$
By construction, $B_m \nearrow \varliminf_n A_n = A$ and $C_m \searrow \varlimsup_n A_n = A$.
Moreover, $B_m \subseteq A_m \subseteq C_m$ by monotonicity of intersection and union.

By (v) applied to $(B_m)$,
$$ \lim_m \mu(B_m) = \mu(A). $$
By (vi) applied to $(C_m)$ (which is legitimate because $\mu(C_1) \le \mu(\bigcup_n A_n) < \infty$),
$$ \lim_m \mu(C_m) = \mu(A). $$
By (i), $\mu(B_m) \le \mu(A_m) \le \mu(C_m)$ for every $m$. By the squeeze
theorem, $\lim_m \mu(A_m) = \mu(A)$.
\end{proof}

\subsection{Borel-Cantelli lemmas}

Recall the probabilistic interpretation of $\varlimsup A_n$: a point $\omega$ belongs
to $\varlimsup A_n = \bigcap_m \bigcup_{k \ge m} A_k$ if and only if it is in infinitely many $A_n$. That
is why one usually writes $\{A_n \text{ infinitely often}\}$. The Borel-Cantelli lemmas give
conditions under which this event has measure zero or one.

\begin{teorema}{Borel-Cantelli (I)}{borel_cantelli_I}
Let $(X, \mathcal{A}, \mu)$ be a measure space and $(A_n) \subset \mathcal{A}$. If
$$ \sum_{n=1}^\infty \mu(A_n) < \infty, $$
then $\mu(\varlimsup_n A_n) = 0$.
\end{teorema}

\begin{proof}
By definition, $\varlimsup_n A_n = \bigcap_{m=1}^\infty C_m$, where $C_m = \bigcup_{k \ge m}
A_k$. By sub-$\sigma$-additivity,
$$ \mu(C_m) = \mu\!\left(\bigcup_{k \ge m} A_k\right) \le \sum_{k=m}^\infty \mu(A_k). $$
Since $\sum_n \mu(A_n) < \infty$, its tail $\sum_{k \ge m} \mu(A_k)$ tends to $0$ as $m \to
\infty$. In particular, $\mu(C_1) \le \sum_{k=1}^\infty \mu(A_k) < \infty$, so that the
decreasing sequence $(C_m)$ satisfies the hypotheses of continuity along decreasing
sequences. Then
$$ \mu\!\left(\varlimsup_n A_n\right) = \mu\!\left(\bigcap_m C_m\right) = \lim_{m \to \infty}
\mu(C_m) \le \lim_{m \to \infty} \sum_{k=m}^\infty \mu(A_k) = 0. $$
By monotonicity, $\mu(\varlimsup A_n) = 0$.
\end{proof}

\obs Borel-Cantelli (I) requires neither independence of the $A_n$ nor that $\mu$ be a probability:
it holds for every measure.

\begin{definicion}{Independence of events}{independencia_eventos}
Let $(\Omega, \mathcal{F}, \mathbb{P})$ be a probability space.
\begin{enumerate}
    \item Two events $A, B \in \mathcal{F}$ are \emph{independent} if
    $\mathbb{P}(A \cap B) = \mathbb{P}(A) \mathbb{P}(B)$.
    \item A finite family $A_1, \dots, A_n \in \mathcal{F}$ is \emph{independent} if for every
    subset of indices $\{i_1, \dots, i_k\} \subseteq \{1, \dots, n\}$ (with $k \ge 1$),
    $$ \mathbb{P}(A_{i_1} \cap \cdots \cap A_{i_k}) = \mathbb{P}(A_{i_1}) \cdots
    \mathbb{P}(A_{i_k}). $$
    \item A family $(A_n)_{n \in \mathbb{N}}$ is \emph{independent} if every finite subfamily
    of it is.
\end{enumerate}
\end{definicion}

\obs If the $(A_n)$ are independent, so are the complements $(A_n^c)$ (in any
finite subfamily). This follows by inclusion-exclusion: for example, $\mathbb{P}(A_1^c \cap
A_2^c) = 1 - \mathbb{P}(A_1) - \mathbb{P}(A_2) + \mathbb{P}(A_1) \mathbb{P}(A_2) = (1 -
\mathbb{P}(A_1))(1 - \mathbb{P}(A_2)) = \mathbb{P}(A_1^c) \mathbb{P}(A_2^c)$; the general case by
induction.

\begin{teorema}{Borel-Cantelli (II)}{borel_cantelli_II}
Let $(\Omega, \mathcal{F}, \mathbb{P})$ be a probability space and $(A_n) \subset \mathcal{F}$ an
\emph{independent} sequence. If
$$ \sum_{n=1}^\infty \mathbb{P}(A_n) = \infty, $$
then $\mathbb{P}(\varlimsup_n A_n) = 1$.
\end{teorema}

\begin{proof}
It will suffice to prove $\mathbb{P}((\varlimsup_n A_n)^c) = 0$. By the properties seen earlier,
$$ (\varlimsup_n A_n)^c = \varliminf_n A_n^c. $$
By sub-$\sigma$-additivity, it is enough to see that for every $m$,
$$ \mathbb{P}\!\left(\bigcap_{k \ge m} A_k^c\right) = 0. $$
Fix $m$. By continuity along decreasing sequences (the finite intersection $\bigcap_{k =
m}^M A_k^c$ decreases to $\bigcap_{k \ge m} A_k^c$ and has probability $\le 1 < \infty$),
$$ \mathbb{P}\!\left(\bigcap_{k \ge m} A_k^c\right) = \lim_{M \to \infty} \mathbb{P}\!\left(
\bigcap_{k = m}^M A_k^c\right). $$
Since $(A_n)$ is independent, so is $(A_n^c)$. Applying the
factorization to the finite family $A_m^c, \dots, A_M^c$,
$$ \mathbb{P}\!\left(\bigcap_{k = m}^M A_k^c\right) = \prod_{k = m}^M \mathbb{P}(A_k^c) =
\prod_{k = m}^M (1 - \mathbb{P}(A_k)). $$
The inequality $1 - x \le e^{-x}$, valid for every $x \in \mathbb{R}$, is used:
$$ \prod_{k = m}^M (1 - \mathbb{P}(A_k)) \le \prod_{k = m}^M e^{-\mathbb{P}(A_k)} =
\exp\!\left(-\sum_{k = m}^M \mathbb{P}(A_k)\right). $$
Since $\sum_n \mathbb{P}(A_n) = \infty$, the tail $\sum_{k = m}^M \mathbb{P}(A_k) \to \infty$ as
$M \to \infty$ (for each fixed $m$). Therefore the exponent tends to $-\infty$ and
$$ \lim_{M \to \infty} \mathbb{P}\!\left(\bigcap_{k = m}^M A_k^c\right) \le \lim_{M \to \infty}
\exp\!\left(-\sum_{k = m}^M \mathbb{P}(A_k)\right) = 0, $$
that is, $\mathbb{P}(\bigcap_{k \ge m} A_k^c) = 0$. Since this holds for every $m$,
$$ \mathbb{P}\!\left((\varlimsup A_n)^c\right) = \mathbb{P}\!\left(\bigcup_m \bigcap_{k \ge m}
A_k^c\right) \le \sum_m \mathbb{P}\!\left(\bigcap_{k \ge m} A_k^c\right) = 0. $$
It follows that $\mathbb{P}(\varlimsup A_n) = 1$.
\end{proof}

\begin{ejemplo}{Increasingly biased coins}{ej_monedas_sesgadas}
Independent coins are tossed, where the $n$-th comes up heads with probability $p_n$, and let
$E_n$ = ``the $n$-th comes up heads''.
\begin{enumerate}
    \item If $p_n = 1/n$: since $\sum 1/n = \infty$ and the events are independent,
    Borel-Cantelli (II) gives that with probability $1$ \emph{infinitely many} heads come up, even though
    $p_n \to 0$.
    \item If $p_n = (99/100)^n$: the series $\sum (99/100)^n$ is a convergent geometric series, and
    Borel-Cantelli (I) gives that with probability $1$ only \emph{finitely many} heads come up, even though the
    first coins come up heads almost always.
\end{enumerate}
\end{ejemplo}

\begin{ejemplo}{The monkey and Borges}{ej_mono}
A monkey types at random, with equiprobable and independent keys, on a keyboard of $60$
symbols. Fix a work of $N$ characters (say, a short story by Borges). The claim is: with
probability $1$, the monkey types it, and not once but \emph{infinitely many} times.

Consider the disjoint blocks of $N$ consecutive keystrokes (the $k$-th block occupies the
positions $(k-1)N + 1, \dots, kN$) and let $S_k$ be the event ``the $k$-th block coincides
character by character with the work''. The $S_k$ depend on disjoint keystrokes, so that they are
independent, and each one has probability $p = (1/60)^N$: tiny, but positive and the
same for every $k$. Then
$$ \sum_{k=1}^\infty \mathbb{P}(S_k) = \sum_{k=1}^\infty \frac{1}{60^N} = \infty, $$
and Borel-Cantelli (II) gives $\mathbb{P}(\varlimsup_k S_k) = 1$: infinitely many blocks reproduce the
complete work.
\end{ejemplo}

\obs The example illustrates the typical use of Borel-Cantelli (II): to conclude ``infinitely
often'' the events need not be likely; it suffices that they be independent and that
the series of their probabilities diverge, which a constant positive probability guarantees.

\obs Combining both directions: if the $(A_n)$ are \emph{independent}, then
$$ \mathbb{P}(\varlimsup A_n) = \begin{cases} 0, & \text{if } \sum \mathbb{P}(A_n) < \infty, \\
1, & \text{if } \sum \mathbb{P}(A_n) = \infty. \end{cases} $$
That is, under independence the event $\{A_n \text{ infinitely often}\}$ always has
probability $0$ or $1$. This dichotomy is a particular case of a much more
general phenomenon: Kolmogorov's $0$-$1$ law.

\subsection{Kolmogorov's 0-1 law}

\begin{definicion}{Tail $\sigma$-algebra}{sigma_cola}
Let $(\Omega, \mathcal{F}, \mathbb{P})$ be a probability space and $(A_n)_{n \in \mathbb{N}}
\subset \mathcal{F}$. The \emph{tail $\sigma$-algebra} of the sequence $(A_n)$ is
$$ \mathcal{T} = \bigcap_{n=1}^\infty \sigma(A_n, A_{n+1}, A_{n+2}, \dots), $$
where $\sigma(A_n, A_{n+1}, \dots)$ is the $\sigma$-algebra generated by the sequence from
index $n$ on. Its elements are called \emph{tail events}: they are the ones that do not depend on
any finite number of the $A_k$.
\end{definicion}

\obs Examples of tail events: $\varlimsup_n A_n$ and $\varliminf_n A_n$ (both are in
$\sigma(A_n, A_{n+1}, \dots)$ for every $n$, by their very definition). The event ``the sequence
$\mathbf{1}_{A_n}$ converges'' is also a tail event. On the other hand, $A_1$ by itself is not.

\begin{teorema}{Kolmogorov's 0-1 law}{kolmogorov_01}
Let $(\Omega, \mathcal{F}, \mathbb{P})$ be a probability space and $(A_n)_{n \in \mathbb{N}}
\subset \mathcal{F}$ an \emph{independent} sequence. If $B \in \mathcal{T}$ is a tail
event, then
$$ \mathbb{P}(B) \in \{0, 1\}. $$
\end{teorema}

\begin{proof}
The strategy is to prove that $B$ is independent of itself, which forces $\mathbb{P}(B) =
\mathbb{P}(B \cap B) = \mathbb{P}(B)^2$ and therefore $\mathbb{P}(B) \in \{0, 1\}$.

The following criterion for independence of generated $\sigma$-algebras is used; it is a consequence of
Dynkin's $\pi$-$\lambda$ theorem, and its proof is omitted: \emph{if two families of events
$\mathcal{C}_1, \mathcal{C}_2$ are closed under finite intersections and every pair $C_1 \in
\mathcal{C}_1$, $C_2 \in \mathcal{C}_2$ is independent, then $\sigma(\mathcal{C}_1)$ and
$\sigma(\mathcal{C}_2)$ are independent} (every pair of elements, one from each, is).

\textit{$\mathcal{T}$ is independent of $\sigma(A_1, \dots, A_n)$, for each $n$.} Let
$\mathcal{C}_1$ be the family of finite intersections of sets $A_i$ with $i \le n$, and
$\mathcal{C}_2$ that of finite intersections of sets $A_j$ with $j > n$. Both are
closed under finite intersections, and every pair of elements from one and the other is independent
because the sequence $(A_n)$ is (the indices involved are all distinct). By the
criterion, $\sigma(A_1, \dots, A_n)$ and $\sigma(A_{n+1}, A_{n+2}, \dots)$ are independent; since
$\mathcal{T} \subseteq \sigma(A_{n+1}, A_{n+2}, \dots)$, in particular $\mathcal{T}$ is
independent of $\sigma(A_1, \dots, A_n)$.

\textit{$\mathcal{T}$ is independent of $\sigma(A_1, A_2, \dots)$.} Let $\mathcal{C} =
\bigcup_n \sigma(A_1, \dots, A_n)$. This family is closed under finite intersections (two
of its elements coexist in some $\sigma(A_1, \dots, A_n)$ with $n$ sufficiently large),
each of its elements is independent of each element of $\mathcal{T}$ by the previous step, and
$\sigma(\mathcal{C}) = \sigma(A_1, A_2, \dots)$. Applying the criterion again (with
$\mathcal{C}$ and $\mathcal{T}$, which is closed under intersections because it is a $\sigma$-algebra),
the independence of $\sigma(A_1, A_2, \dots)$ and $\mathcal{T}$ is obtained.

\textit{Conclusion.} Let $B \in \mathcal{T}$. Also $B \in \sigma(A_1, A_2, \dots)$, since
$\mathcal{T} \subseteq \sigma(A_1, A_2, \dots)$. The independence just proved, applied to the
pair $(B, B)$, gives
$$ \mathbb{P}(B) = \mathbb{P}(B \cap B) = \mathbb{P}(B) \cdot \mathbb{P}(B), $$
whence $\mathbb{P}(B)(1 - \mathbb{P}(B)) = 0$ and $\mathbb{P}(B) \in \{0, 1\}$.
\end{proof}

\obs The law does not say \emph{which} of the two values $\mathbb{P}(B)$ takes, and deciding it can be
difficult. For the tail event $\varlimsup_n A_n$, the Borel-Cantelli lemmas are the
effective criterion.

\newpage

\section{Generating measures}

\subsection{Outer measure}

\begin{definicion}{Outer measure}{medida_exterior}
A function $\mu^*: \mathcal{P}(X) \to [0, \infty]$ is an \emph{outer measure} if:
\begin{enumerate}
    \item $\mu^*(\emptyset) = 0$.
    \item Monotonicity: $A \subseteq B \implies \mu^*(A) \le \mu^*(B)$.
    \item Sub-$\sigma$-additivity: $\mu^*(\bigcup_n A_n) \le \sum_n \mu^*(A_n)$ for any
    $(A_n) \subseteq \mathcal{P}(X)$.
\end{enumerate}
\end{definicion}

\obs \emph{Typical construction of an outer measure.} Given a function $\ell:
\mathcal{C} \to [0, \infty]$ defined on a family $\mathcal{C} \subseteq \mathcal{P}(X)$
with $\emptyset \in \mathcal{C}$ and $\ell(\emptyset) = 0$, define
$$ \mu^*(E) = \inf\left\{ \sum_{n=1}^\infty \ell(I_n) : (I_n) \subset \mathcal{C}, \ E \subseteq
\bigcup_n I_n \right\}, $$
with $\mu^*(E) = +\infty$ if $E$ admits no cover. Then $\mu^*$ is an outer measure
.

\subsection{Carathéodory's criterion}

\begin{definicion}{$\mu^*$-measurable sets}{caratheodory_mide}
$A \subseteq X$ is \emph{$\mu^*$-measurable} (in the sense of Carathéodory) if for every
$E \subseteq X$:
$$ \mu^*(E) = \mu^*(E \cap A) + \mu^*(E \cap A^c). $$
\end{definicion}

\begin{teorema}{Carathéodory}{caratheodory}
Let $\mu^*$ be an outer measure on $X$, and let
$$ \mathcal{A}^* = \{ A \subseteq X : A \text{ is } \mu^*\text{-measurable} \}. $$
Then $\mathcal{A}^*$ is a $\sigma$-algebra and $\mu = \mu^*|_{\mathcal{A}^*}$ is a measure.
\end{teorema}

\begin{proof}\leavevmode\\
For every $A \subseteq X$ and every $E \subseteq X$, subadditivity gives $\mu^*(E) = \mu^*((E \cap A) \cup (E \cap A^c)) \le \mu^*(E \cap A) + \mu^*(E \cap A^c)$. Therefore, to see that a set is in $\mathcal{A}^*$ it suffices to prove the opposite inequality.

\textit{Step 1: $\mathcal{A}^*$ is an algebra.} The empty set is in $\mathcal{A}^*$, since $\mu^*(E \cap \emptyset) + \mu^*(E \cap X) = 0 + \mu^*(E) = \mu^*(E)$ for every $E$. Closure under complements follows from the fact that Carathéodory's condition is symmetric in $A$ and $A^c$. For closure under finite unions, let $A, B \in \mathcal{A}^*$ and $E \subseteq X$. Applying the measurability of $A$ to $E$, and then that of $B$ to $E \cap A^c$,
\begin{align*}
\mu^*(E) &= \mu^*(E \cap A) + \mu^*(E \cap A^c) \\
&= \mu^*(E \cap A) + \mu^*(E \cap A^c \cap B) + \mu^*(E \cap A^c \cap B^c).
\end{align*}
Since $E \cap (A \cup B) = (E \cap A) \sqcup (E \cap A^c \cap B)$, subadditivity gives $\mu^*(E \cap (A \cup B)) \le \mu^*(E \cap A) + \mu^*(E \cap A^c \cap B)$. Substituting in the previous equality and using $E \cap (A \cup B)^c = E \cap A^c \cap B^c$,
$$ \mu^*(E) \ge \mu^*(E \cap (A \cup B)) + \mu^*(E \cap (A \cup B)^c), $$
so that $A \cup B \in \mathcal{A}^*$.

\textit{Step 2: $\mathcal{A}^*$ is a $\sigma$-algebra.} Let $(A_n)_{n \in \mathbb{N}} \subseteq \mathcal{A}^*$ be a disjoint sequence and $A = \bigcup_n A_n$. Write $S_n = \bigcup_{k=1}^n A_k$. It is proved by induction on $n$ that, for every $E$,
$$ \mu^*(E \cap S_n) = \sum_{k=1}^n \mu^*(E \cap A_k). $$
The case $n = 1$ holds because $S_1 = A_1$. Assuming the formula for $n$, observe that $S_n \in \mathcal{A}^*$ since it is a finite union, and apply Carathéodory's condition for $A_{n+1}$ to the set $E \cap S_{n+1}$:
$$ \mu^*(E \cap S_{n+1}) = \mu^*(E \cap S_{n+1} \cap A_{n+1}) + \mu^*(E \cap S_{n+1} \cap A_{n+1}^c). $$
Since the $A_k$ are disjoint, $S_{n+1} \cap A_{n+1} = A_{n+1}$ and $S_{n+1} \cap A_{n+1}^c = S_n$, so that
$$ \mu^*(E \cap S_{n+1}) = \mu^*(E \cap A_{n+1}) + \mu^*(E \cap S_n). $$
Applying the inductive hypothesis to $\mu^*(E \cap S_n)$,
$$ \mu^*(E \cap S_{n+1}) = \mu^*(E \cap A_{n+1}) + \sum_{k=1}^n \mu^*(E \cap A_k) = \sum_{k=1}^{n+1} \mu^*(E \cap A_k), $$
which is the formula for $n+1$.

Now, by the measurability of $S_n$,
$$ \mu^*(E) = \mu^*(E \cap S_n) + \mu^*(E \cap S_n^c). $$
In the first term the formula just proved is used; in the second, since $S_n \subseteq A$ gives $S_n^c \supseteq A^c$, monotonicity ensures $\mu^*(E \cap S_n^c) \ge \mu^*(E \cap A^c)$. Therefore
$$ \mu^*(E) \ge \sum_{k=1}^n \mu^*(E \cap A_k) + \mu^*(E \cap A^c), $$
and since this holds for every $n$, letting $n \to \infty$,
$$ \mu^*(E) \ge \sum_{k=1}^\infty \mu^*(E \cap A_k) + \mu^*(E \cap A^c). $$
Finally, $E \cap A = \bigcup_k (E \cap A_k)$, so that sub-$\sigma$-additivity gives $\sum_{k=1}^\infty \mu^*(E \cap A_k) \ge \mu^*(E \cap A)$, and then
$$ \mu^*(E) \ge \mu^*(E \cap A) + \mu^*(E \cap A^c), $$
that is, $A \in \mathcal{A}^*$. For an arbitrary countable union $(B_n) \subseteq \mathcal{A}^*$, make it disjoint: $A_1 = B_1$ and $A_n = B_n \setminus (B_1 \cup \cdots \cup B_{n-1})$ for $n \ge 2$. Since $\mathcal{A}^*$ is an algebra, each $A_n \in \mathcal{A}^*$, the $A_n$ are disjoint and $\bigcup_n B_n = \bigcup_n A_n \in \mathcal{A}^*$ by the above. Hence $\mathcal{A}^*$ is a $\sigma$-algebra.

\textit{Step 3: $\mu = \mu^*|_{\mathcal{A}^*}$ is a measure.} Let $(A_n) \subseteq \mathcal{A}^*$ be disjoint and $A = \bigcup_n A_n$. Evaluate the inequality obtained in Step 2 at $E = A$. Since $A \cap A_k = A_k$ and $A \cap A^c = \emptyset$,
$$ \mu^*(A) \ge \sum_{k=1}^\infty \mu^*(A \cap A_k) + \mu^*(A \cap A^c) = \sum_{k=1}^\infty \mu^*(A_k). $$
The opposite inequality $\mu^*(A) \le \sum_k \mu^*(A_k)$ is sub-$\sigma$-additivity, so that $\mu(A) = \sum_n \mu(A_n)$; that is, $\mu$ is $\sigma$-additive on disjoint sequences of $\mathcal{A}^*$. Since moreover $\mu(\emptyset) = \mu^*(\emptyset) = 0$, $\mu$ is a measure.
\end{proof}

\subsection{Lebesgue measure on \texorpdfstring{$\mathbb{R}$}{R}}

\begin{definicion}{Lebesgue outer measure}{lebesgue_ext}
For $E \subseteq \mathbb{R}$, define
$$ \lambda^*(E) = \inf\!\left\{ \sum_{n=1}^\infty (b_n - a_n) : E \subseteq \bigcup_n (a_n, b_n]
\right\}, $$
\end{definicion}

\begin{teorema}{Lebesgue measure}{lebesgue}
$\lambda^*$ is an outer measure. The associated Carathéodory $\sigma$-algebra, denoted
$\mathcal{L}$, contains $\mathcal{B}(\mathbb{R})$, and the restriction $\lambda =
\lambda^*|_{\mathcal{L}}$ is a measure (called \emph{Lebesgue measure}) that satisfies:
\begin{enumerate}
    \item $\lambda(a, b]) = b - a$.
    \item $\mathcal{B}(\mathbb{R}) \subset \mathcal{L} \subsetneq \mathcal{P}(\mathbb{R})$.
    \item $\lambda$ is translation invariant: $\lambda(A + t) = \lambda(A)$ for every $t \in
    \mathbb{R}$ and $A \in \mathcal{L}$.
    \item $(\mathbb{R}, \mathcal{L}, \lambda)$ is complete: if $\lambda(N) = 0$ and $A \subseteq N$,
    then $A \in \mathcal{L}$ (and by monotonicity $\lambda(A) = 0$).
\end{enumerate}
\end{teorema}

\begin{proof}\leavevmode\\
(i) The inequality $\lambda^*((a, b]) \le b - a$
is obtained by taking $(a, b]$ itself as the cover. The other inequality is proved now. Let
$\{(a_k, b_k]\}_k$ be a countable cover of $(a, b]$. Given $\epsilon > 0$, enlarge each
interval to an \emph{open} one $(a_k - \epsilon_k, b_k + \epsilon_k)$, choosing $\epsilon_k > 0$
so that $\sum_k \epsilon_k < \epsilon$. Thus, $\{(a_k - \epsilon_k, b_k + \epsilon_k)\}$ is an
open cover of the compact set $[a, b]$.

By compactness, there is a finite subcover $\{(a_{k_1} - \epsilon_{k_1}, b_{k_1} +
\epsilon_{k_1}), \dots, (a_{k_N} - \epsilon_{k_N}, b_{k_N} + \epsilon_{k_N})\}$ of $[a, b]$. The sum of the lengths
of the intervals of the subcover is at least $b - a$:
$$ b - a \le \sum_{j=1}^N (b_{k_j} + \epsilon_{k_j} - (a_{k_j} - \epsilon_{k_j})) = \sum_{j=1}^N
(b_{k_j} - a_{k_j}) + 2 \sum_{j=1}^N \epsilon_{k_j} \le \sum_k (b_k - a_k) + 2\epsilon. $$
Letting $\epsilon \to 0$ and then taking the infimum over covers, $b - a \le \lambda^*((a, b])$.

(ii) It suffices to prove that each half-line
$(c, \infty)$ is $\lambda^*$-measurable, because the half-lines generate $\mathcal{B}(\mathbb{R})$.
Fix $c \in \mathbb{R}$ and $E \subseteq \mathbb{R}$. Given a cover $E \subseteq \bigcup_k
(a_k, b_k]$, split each interval into its piece to the right of $c$ and its piece to the left:
$$ (a_k, b_k] = ((a_k, b_k] \cap (c, \infty)) \sqcup ((a_k, b_k] \cap (-\infty, c]). $$
Each piece is a half-open interval (possibly empty) and their lengths add up to exactly
$b_k - a_k$. Applying $\lambda^*$ to the two collections of pieces, which cover
$E \cap (c, \infty)$ and $E \cap (-\infty, c]$ respectively,
$$ \lambda^*(E \cap (c, \infty)) + \lambda^*(E \cap (-\infty, c]) \le \sum_k (b_k - a_k). $$
Taking the infimum over covers on the right,
$$ \lambda^*(E \cap (c, \infty)) + \lambda^*(E \cap (c, \infty)^c) \le \lambda^*(E). $$
The opposite inequality is subadditivity. Hence $(c, \infty)$ satisfies Carathéodory and belongs to
$\mathcal{L}$.

(iii) If $\{(a_k, b_k]\}$ covers $E$, then
$\{(a_k + t, b_k + t]\}$ covers $E + t$ with the same sum of lengths. This gives $\lambda^*(E + t)
\le \lambda^*(E)$, and by symmetry (applying the same to $-t$) equality. Invariance carries
over to $\lambda^*$-measurability: if $A \in \mathcal{L}$, Carathéodory's condition
for $A + t$ follows by applying that of $A$ to the set $E - t$.

(iv) If $\lambda(N) = 0$ and $A \subseteq N$, by monotonicity $\lambda^*(A)
= 0$. For any $E$, $\lambda^*(E \cap A) \le \lambda^*(A) = 0$, so that Carathéodory's
condition for $A$ reduces to $\lambda^*(E) = \lambda^*(E \cap A^c)$, which holds by
subadditivity and monotonicity. Hence $A \in \mathcal{L}$.

(v) For $\mathcal{L} \subsetneq \mathcal{P}(\mathbb{R})$, the Vitali
sets cannot be Lebesgue-measurable because $\lambda$ satisfies the
three properties in the conclusion of the theorem.
\end{proof}

\obs There are Lebesgue-null sets with the cardinality of the continuum (Cantor), as well as
nowhere dense sets of positive measure (fat Cantor).

\obs Every countable set has Lebesgue measure zero: for each point,
$\lambda(\{x\}) \le \lambda((x - \epsilon, x]) = \epsilon$ for every $\epsilon > 0$, whence
$\lambda(\{x\}) = 0$, and by $\sigma$-additivity
$$ \lambda(\mathbb{Q}) = \sum_{q \in \mathbb{Q}} \lambda(\{q\}) = 0. $$
Thus, the function $\mathbf{1}_{\mathbb{Q} \cap [0,1]}$ from the motivation, which was not
Riemann-integrable, is zero except on a set of measure zero: its Lebesgue integral
will be $0$.

\subsection{Lebesgue-Stieltjes}

\begin{definicion}{Lebesgue-Stieltjes outer measure}{lebesgue_stieltjes_ext}
Let $F: \mathbb{R} \to \mathbb{R}$ be nondecreasing and right-continuous. For $E \subseteq
\mathbb{R}$,
$$ \lambda_F^*(E) = \inf\!\left\{ \sum_{n=1}^\infty (F(b_n) - F(a_n)) : E \subseteq \bigcup_n
(a_n, b_n] \right\}. $$
\end{definicion}

\begin{teorema}{Lebesgue-Stieltjes measure}{lebesgue_stieltjes}
Let $F$ be as above. Then $\lambda_F^*$ is an outer measure, and the Carathéodory
$\sigma$-algebra $\mathcal{L}_F$ contains $\mathcal{B}(\mathbb{R})$. The restriction $\lambda_F =
\lambda_F^*|_{\mathcal{L}_F}$ is the \emph{Lebesgue-Stieltjes measure} of $F$, and it satisfies
$$ \lambda_F((a, b]) = F(b) - F(a) \qquad \forall a < b. $$
\end{teorema}

\begin{proof}
The identity $\lambda_F((a, b]) = F(b) - F(a)$ uses, in the compactness step, that
$F$ is right-continuous: to show
$F(b) - F(a) \le \sum_k (F(b_k) - F(a_k))$ when $\{(a_k, b_k]\}$ covers $(a, b]$, the
intervals are widened to $(a_k, b_k + \delta_k]$ with $F(b_k + \delta_k) - F(b_k) < \epsilon/2^k$
 and $[a + \delta, b]$ is compactified as for Lebesgue.
The $\lambda_F^*$-measurability of the half-lines and completeness follow exactly in the same way.
\end{proof}

\obs When $F$ is the
distribution function of a random variable $X$, $\lambda_F$ is the
distribution of $X$, denoted $P_X$, and $P_X(B) = P(X \in B)$ for every Borel set $B$.

\newpage

\section{Measurable functions}

\begin{definicion}{Measurable function}{medible}
Let $(X, \mathcal{A}, \mu)$ be a measure space. A function $f: X \to \mathbb{R}$ is
\emph{measurable} if for every $a \in \mathbb{R}$,
$$ f^{-1}((a, \infty)) = \{ x \in X : f(x) > a \} \in \mathcal{A}. $$
\end{definicion}

\obs In probability, a measurable function is called a \emph{random variable}.

\obs The definition points directly at the Lebesgue integral: to approximate $\int f$
by partitioning the \emph{range} into values $a_0 < a_1 < \cdots < a_n$, one must measure the sets
$$ A_k = \{a_{k-1} < f \le a_k\} = \{f > a_{k-1}\} \setminus \{f > a_k\}. $$
Requiring $\{f > a\} \in \mathcal{A}$ for every $a$ is exactly what guarantees that those
sets have a measure.

\begin{proposicion}{Equivalent characterizations of measurability}{caract_medible}
Let $(X, \mathcal{A})$ be a measurable space and $f: X \to \mathbb{R}$. The following conditions are
equivalent:
\begin{enumerate}
    \item $\{f > a\} \in \mathcal{A}$ for every $a \in \mathbb{R}$.
    \item $\{f \ge a\} \in \mathcal{A}$ for every $a \in \mathbb{R}$.
    \item $\{f < a\} \in \mathcal{A}$ for every $a \in \mathbb{R}$.
    \item $\{f \le a\} \in \mathcal{A}$ for every $a \in \mathbb{R}$.
\end{enumerate}
\end{proposicion}

\begin{proof}\leavevmode\\
(i) $\Leftrightarrow$ (iv), (ii) $\Leftrightarrow$ (iii) and (i) $\Leftrightarrow$ (ii) are proved.

\textit{(i) $\Leftrightarrow$ (iv).} $\{f \le a\} = \{f > a\}^c$, and $\mathcal{A}$ is closed under
complements.

\textit{(ii) $\Leftrightarrow$ (iii).} $\{f < a\} = \{f \ge a\}^c$.

\textit{(i) $\Rightarrow$ (ii).} $\{f \ge a\} = \bigcap_{n=1}^\infty \{f > a - 1/n\}$,
a countable
intersection of elements of $\mathcal{A}$.

\textit{(ii) $\Rightarrow$ (i).} $\{f > a\} = \bigcup_{n=1}^\infty \{f \ge a + 1/n\}$, a countable
union of elements of $\mathcal{A}$.

With this, (i)--(iv) are equivalent.
\end{proof}

\subsection{Examples and operations}

\begin{ejemplo}{Constant functions}{}
If $f \equiv c$, then $\{f > a\}$ equals $X$ if $c > a$ and $\emptyset$ if $c \le a$. Since
$X, \emptyset \in \mathcal{A}$, $f$ is measurable.
\end{ejemplo}

\begin{ejemplo}{Characteristic functions}{}
Let $A \in \mathcal{A}$. The characteristic function $\mathbf{1}_A$ is measurable:
$$ \{\mathbf{1}_A > a\} = \begin{cases} X, & a < 0, \\ A, & 0 \le a < 1, \\ \emptyset,
& a \ge 1, \end{cases} $$
and the three sets are in $\mathcal{A}$.
\end{ejemplo}

\begin{teorema}{Continuous functions are Borel-measurable}{continuas_borel}
Every continuous function $f: \mathbb{R} \to \mathbb{R}$ is Borel-measurable.
\end{teorema}

\begin{proof}
It must be seen that $f^{-1}((a, \infty)) \in \mathcal{B}(\mathbb{R})$ for every $a \in \mathbb{R}$.
Since $(a, \infty)$ is open and $f$ continuous, $f^{-1}((a, \infty))$ is open in $\mathbb{R}$, and
every open set belongs to $\mathcal{B}(\mathbb{R})$.

More generally, the preimage under a continuous function of any open set is open, which gives
$f^{-1}(\mathcal{B}(\mathbb{R})) \subseteq \mathcal{B}(\mathbb{R})$: the set $\{B \subseteq
\mathbb{R} : f^{-1}(B) \in \mathcal{B}(\mathbb{R})\}$ is a $\sigma$-algebra containing the
open sets, so that it contains $\mathcal{B}(\mathbb{R})$.
\end{proof}

\begin{teorema}{Monotone functions are Borel-measurable}{monotonas_borel}
Every monotone function $f: \mathbb{R} \to \mathbb{R}$ is Borel-measurable.
\end{teorema}

\begin{proof}
Suppose $f$ nondecreasing (the nonincreasing case follows by applying what follows to $-f$). With
$a \in \mathbb{R}$ fixed, consider the set
$E_a = f^{-1}((a, \infty)) = \{x \in \mathbb{R} : f(x) > a\}$.

If $x \in E_a$ and $y > x$, monotonicity gives $f(y) \ge f(x) > a$, so that $y \in E_a$.
Therefore $E_a$ is a half-line to the right: denoting $c = \inf E_a$ (with the convention
$\inf \emptyset = +\infty$, and $\inf$ of a set not bounded below $= -\infty$), then
$$ E_a = (c, \infty) \quad \text{or} \quad E_a = [c, \infty), $$
according to whether the infimum is attained or not. In either case, $E_a \in
\mathcal{B}(\mathbb{R})$ (open and closed half-lines are Borel sets).
\end{proof}

\begin{ejemplo}{Simple functions}{}
A function $\varphi: X \to \mathbb{R}$ is \emph{simple} if it is a finite linear combination of
characteristic functions:
$$ \varphi = \sum_{i=1}^n c_i \, \mathbf{1}_{A_i}, \qquad c_i \in \mathbb{R}, \ A_i \in \mathcal{A}.
$$
Equivalently, $\varphi$ takes only finitely many values and the preimage of each one is in
$\mathcal{A}$. Simple functions are measurable because each $\mathbf{1}_{A_i}$ is and the operations of
sum and product by a scalar preserve measurability (Theorem \ref{teo:operaciones_medibles}).
\end{ejemplo}

\begin{ejemplo}{Canonical representation of a simple function}{ej_repr_canonica}
Writing a simple function as a combination of characteristic functions is not unique. With intervals of
$\mathbb{R}$, consider
$$ \varphi = 2 \cdot \mathbf{1}_{[0,2]} + 3 \cdot \mathbf{1}_{[1,5]}. $$
The two intervals overlap on $[1,2]$, where $\varphi$ equals $2 + 3 = 5$. Separating according to
the values that $\varphi$ takes,
$$ \varphi = 2 \cdot \mathbf{1}_{[0,1)} + 5 \cdot \mathbf{1}_{[1,2]} + 3 \cdot \mathbf{1}_{(2,5]},
$$
which is the \emph{canonical representation}: the coefficients are the distinct nonzero values
of $\varphi$, and the sets $\{\varphi = 2\} = [0,1)$, $\{\varphi = 5\} = [1,2]$ and
$\{\varphi = 3\} = (2,5]$ are disjoint. Every simple function admits a unique representation of this
form, which is the one used by the definition of the integral.
\end{ejemplo}

\begin{teorema}{Operations with measurable functions}{operaciones_medibles}
Let $f, g: X \to \mathbb{R}$ be measurable and $c \in \mathbb{R}$. Then $cf$, $f + g$, $f \cdot g$,
$\max(f, g)$, $\min(f, g)$, $|f|$, $f^+ := \max(f, 0)$ and $f^- := \max(-f, 0)$ are measurable.
\end{teorema}

\begin{proof}\leavevmode\\
By Proposition \ref{prop:caract_medible}, to see that a function is measurable it suffices to
verify any of the conditions $\{f > a\}, \{f \ge a\}, \{f < a\}, \{f \le a\} \in
\mathcal{A}$ for every $a \in \mathbb{R}$; in each case the most convenient one is chosen.

\textit{Product by a scalar.}
\begin{itemize}
    \item If $c = 0$, $cf \equiv 0$, which is constant and therefore measurable.
    \item If $c > 0$, $cf(x) > a \Leftrightarrow f(x) > a/c$, so that
    $\{cf > a\} = \{f > a/c\} \in \mathcal{A}$.
    \item If $c < 0$, $cf(x) > a \Leftrightarrow f(x) < a/c$, so that 
    $\{cf > a\} = \{f < a/c\} \in \mathcal{A}$.
\end{itemize}

\textit{Maximum and minimum.} Since
$$ \max(f, g)(x) > a \ \Leftrightarrow \ f(x) > a \ \text{ or } \ g(x) > a, $$
one has $\{\max(f, g) > a\} = \{f > a\} \cup \{g > a\} \in \mathcal{A}$, being a finite union
of measurable sets. Analogously,
$$ \min(f, g)(x) \ge a \ \Leftrightarrow \ f(x) \ge a \ \text{ and } \ g(x) \ge a, $$
so that $\{\min(f, g) \ge a\} = \{f \ge a\} \cap \{g \ge a\} \in \mathcal{A}$.

\textit{Absolute value.} For $a < 0$, $\{|f| > a\} = X \in \mathcal{A}$. For $a \ge 0$,
$$ |f(x)| > a \ \Leftrightarrow \ f(x) > a \ \text{ or } \ f(x) < -a, $$
so that $\{|f| > a\} = \{f > a\} \cup \{f < -a\} \in \mathcal{A}$.

\textit{Positive and negative parts $f^+, f^-$.} The constant function $0$ is measurable, and by the
case of the maximum applied to $f$ and $0$ (or to $-f$ and $0$, using the case of product by a scalar with
$c = -1$),
$$ f^+ = \max(f, 0), \qquad f^- = \max(-f, 0), $$
are measurable. Moreover the identities $f = f^+ - f^-$ and $|f| = f^+ + f^-$ hold.

\textit{Sum.} It is shown that for every $a \in \mathbb{R}$,
$$ \{f + g > a\} = \bigcup_{r \in \mathbb{Q}} \left( \{f > r\} \cap \{g > a - r\} \right). $$
\begin{itemize}
    \item ($\subseteq$) Let $x$ be such that $f(x) + g(x) > a$, that is $f(x) > a - g(x)$. By density of
    $\mathbb{Q}$ in $\mathbb{R}$, there is $r \in \mathbb{Q}$ with
    $$ a - g(x) < r < f(x). $$
    From $r < f(x)$ one gets $x \in \{f > r\}$, and from $a - g(x) < r$ one gets $g(x) > a - r$, that is $x \in
    \{g > a - r\}$.

    \item ($\supseteq$) Let $r \in \mathbb{Q}$ with $f(x) > r$ and $g(x) > a - r$. Adding, $f(x) +
    g(x) > r + (a - r) = a$, that is $x \in \{f + g > a\}$.
\end{itemize}
The right-hand side is a \emph{countable} union of
finite intersections of sets in $\mathcal{A}$, and therefore belongs to $\mathcal{A}$.

\textit{Square.} For $a < 0$, $\{f^2 > a\} = X \in \mathcal{A}$. For $a \ge 0$,
$$ f(x)^2 > a \ \Leftrightarrow \ |f(x)| > \sqrt{a} \ \Leftrightarrow \ f(x) > \sqrt{a} \ \text{ or }
\ f(x) < -\sqrt{a}, $$
and therefore $\{f^2 > a\} = \{f > \sqrt{a}\} \cup \{f < -\sqrt{a}\} \in \mathcal{A}$.

\textit{Product.} The polarization identity is used:
$$ fg = \tfrac{1}{4} \left[ (f + g)^2 - (f - g)^2 \right]. $$
By the previous steps, $f + g$ and $f - g = f + (-1) g$ are measurable.
Their squares $(f + g)^2$ and $(f - g)^2$ are measurable by the previous case. Applying sum and
product by a scalar once more,
$$ fg = \tfrac{1}{4}(f + g)^2 + \left(-\tfrac{1}{4}\right)(f - g)^2 \in \mathcal{A}, $$
that is, $fg$ is measurable.
\end{proof}

\obs The positive and negative parts satisfy, in addition to $f = f^+ - f^-$ and $|f| = f^+ + f^-$,
the inverse formulas
$$ f^+ = \tfrac{1}{2}(|f| + f), \qquad f^- = \tfrac{1}{2}(|f| - f), $$
with $f^+, f^- \ge 0$ of disjoint supports. The decomposition $f = f^+ - f^-$ is what reduces
the integral of a general function to that of two nonnegative ones.

\subsection{Sup, inf, lim sup, lim inf}

\begin{teorema}{Sup and inf of measurable functions}{sup_inf_medibles}
If $(f_n)$ is a sequence of measurable functions, then $\sup_n f_n$, $\inf_n f_n$,
$\varlimsup_n f_n$ and $\varliminf_n f_n$ are measurable.
\end{teorema}

\begin{proof}\leavevmode\\
\textit{Case of the supremum.} For every $a \in \mathbb{R}$,
$$ \sup_n f_n(x) > a \ \Leftrightarrow \ \exists \, n \in \mathbb{N} : f_n(x) > a, $$
whence
$$ \{\sup_n f_n > a\} = \bigcup_{n \in \mathbb{N}} \{f_n > a\}. $$
Each $\{f_n > a\} \in \mathcal{A}$ by measurability of $f_n$, and the union is countable, so that
$\{\sup_n f_n > a\} \in \mathcal{A}$ for every $a$.

\textit{Case of the infimum.} Analogously,
$$ \inf_n f_n(x) < a \ \Leftrightarrow \ \exists \, n \in \mathbb{N} : f_n(x) < a, $$
so that $\{\inf_n f_n < a\} = \bigcup_n \{f_n < a\} \in \mathcal{A}$.

\textit{Limit superior.} By definition,
$$ \varlimsup_n f_n = \inf_{m \in \mathbb{N}} \sup_{n \ge m} f_n. $$
For each $m$, the sequence $(f_n)_{n \ge m}$ is measurable, so that $g_m := \sup_{n \ge m} f_n$
is measurable by the case of the sup. The sequence $(g_m)_m$ is then measurable, and
$\varlimsup_n f_n = \inf_m g_m$ is measurable by the case of the inf.

\textit{Limit inferior.} Analogous, exchanging $\sup$ with $\inf$: $\varliminf_n f_n =
\sup_m \inf_{n \ge m} f_n$, where $\inf_{n \ge m} f_n$ is measurable by the case of the inf and then the
sup over $m$ by the case of the sup.
\end{proof}

\obs The four operations $\sup, \inf, \varlimsup, \varliminf$ may yield functions with values
in $[-\infty, +\infty]$. The definition of measurability extends without problems:
$f: X \to [-\infty,
+\infty]$ is measurable if $\{f > a\} \in \mathcal{A}$ for every $a \in \mathbb{R}$.

\begin{corolario}{Pointwise limit of measurable functions is measurable}{lim_medibles}
Let $(f_n)$ be a sequence of measurable functions.
\begin{enumerate}
    \item If $f_n \to f$ pointwise on all of $X$, $f$ is measurable.
    \item If the series $g := \sum_{k=1}^\infty f_k$ converges pointwise, $g$ is measurable.
\end{enumerate}
\end{corolario}

\begin{proof}\leavevmode\\
\textit{(i)} If $f_n \to f$ pointwise, $f = \varlimsup_n f_n = \varliminf_n f_n$, and by the
previous theorem both are measurable. Hence $f$ is measurable.

\textit{(ii)} The partial sums $g_n = \sum_{k=1}^n f_k$ are finite sums of measurable functions, and by
Theorem \ref{teo:operaciones_medibles} they are measurable. Since $g = \lim_n g_n$ pointwise (by
hypothesis), by (i) $g$ is measurable.
\end{proof}

\subsection{Approximation by simple functions}

\begin{teorema}{Approximation by simple functions}{aprox_simples}
Let $f: X \to \mathbb{R}$ be measurable and nonnegative. There is a sequence $(\varphi_n)$ of nonnegative
simple functions such that $\varphi_n \nearrow f$ pointwise. If moreover $f$ is bounded, the
convergence is uniform.
\end{teorema}

\begin{proof}\leavevmode\\
The idea is to discretize the \emph{values} of $f$. For $n \in \mathbb{N}$ and $0 \le k \le n 2^n - 1$,
write $I_{n,k} = [k/2^n, (k+1)/2^n)$ and define
$$ \varphi_n(x) = \begin{cases}
\dfrac{k}{2^n}, & \text{if } f(x) \in I_{n,k}, \\
n, & \text{if } f(x) \ge n.
\end{cases} $$

The intervals $I_{n,0}, \dots, I_{n, n2^n - 1}$ together with $[n, \infty)$ form a partition of
$[0, \infty)$, so that $\varphi_n$ is well defined and takes at most $n 2^n + 1$ values.
Its level sets are the preimages $\{\varphi_n = k/2^n\} = f^{-1}(I_{n,k})$ and
$\{\varphi_n = n\} = f^{-1}([n, \infty))$, measurable because $f$ is; hence $\varphi_n$ is simple. On
the other hand, $\varphi_n(x) \le f(x)$ in each case ($k/2^n \le f(x)$ when $f(x) \in I_{n,k}$, and
$n \le f(x)$ when $f(x) \ge n$), so that $0 \le \varphi_n \le f$.

It is shown that $\varphi_n \le \varphi_{n+1}$ pointwise. Since $I_{n,k} = I_{n+1, 2k} \sqcup
I_{n+1, 2k+1}$, if $f(x) \in I_{n,k}$ with $k \le n 2^n - 1$ then $f(x)$ belongs to one of those
two intervals, on which $\varphi_{n+1}(x)$ equals $2k/2^{n+1} = k/2^n$ or
$(2k+1)/2^{n+1} > k/2^n$; in both cases $\varphi_{n+1}(x) \ge k/2^n = \varphi_n(x)$. If instead
$f(x) \ge n$, then $\varphi_n(x) = n$ and $\varphi_{n+1}(x) \ge n$: it equals $n+1$ if $f(x) \ge n+1$,
and if $n \le f(x) < n+1$ there is $k$ with $f(x) \in I_{n+1,k}$ and $k/2^{n+1} \ge n$, whence
$\varphi_{n+1}(x) = k/2^{n+1} \ge n$. In all cases $\varphi_{n+1}(x) \ge \varphi_n(x)$.

Now let $x$ be such that $f(x) < \infty$ and take $n > f(x)$. Then $f(x) < n$, so that
$f(x) \in I_{n,k}$ for some $k$ and, by definition of $I_{n,k}$,
$$ 0 \le f(x) - \varphi_n(x) = f(x) - \frac{k}{2^n} < \frac{1}{2^n}. $$
Since $2^{-n} \to 0$, one obtains $\varphi_n(x) \to f(x)$.

Finally, if $f \le M$, for every $n > M$ one has $f(x) < n$ whatever $x$ is, so that the
bound $0 \le f(x) - \varphi_n(x) < 2^{-n}$ holds uniformly in $x$. Since $2^{-n} \to 0$, the
convergence is uniform.
\end{proof}

\newpage

\section{Convergence of measurable functions}

\begin{definicion}{Convergence almost everywhere}{conv_ctp}
$f_n \to f$ \emph{almost everywhere} (a.e.) if there is $N \in \mathcal{A}$ with
$\mu(N) = 0$ such that $f_n(x) \to f(x)$ for every $x \notin N$.
\end{definicion}

\begin{definicion}{Convergence in measure}{conv_medida}
$f_n \to f$ \emph{in measure} if for every $\epsilon > 0$,
$$ \mu(\{x : |f_n(x) - f(x)| \ge \epsilon\}) \xrightarrow{n \to \infty} 0. $$
\end{definicion}

\begin{definicion}{Almost uniform convergence}{conv_casi_unif}
$f_n \to f$ \emph{almost uniformly} if for every $\delta > 0$ there is $N \in \mathcal{A}$ with
$\mu(N) < \delta$ such that $f_n \to f$ uniformly on $N^c$.
\end{definicion}

\begin{corolario}{A.e. limit of measurable functions}{lim_ctp_medible}
If $\mathcal{A}$ is complete with respect to $\mu$, the $f_n$ are measurable and $f_n \to f$ a.e.,
then $f$ is measurable.
\end{corolario}

\begin{proof}
There is $N \in \mathcal{A}$ with $\mu(N) = 0$ such that $f_n(x) \to f(x)$ for every $x \notin N$. Let
$\tilde f = \varlimsup_n f_n$, measurable by Theorem \ref{teo:sup_inf_medibles}; on $N^c$ the
sequence converges, so that $f = \tilde f$ on $N^c$. With $a \in \mathbb{R}$ fixed, decompose
$$ \{f > a\} = \big(\{f > a\} \cap N^c\big) \cup \big(\{f > a\} \cap N\big). $$
The first set coincides with $\{\tilde f > a\} \cap N^c$, an intersection of two measurable sets. The
second satisfies $\{f > a\} \cap N \subseteq N$, and since $\mathcal{A}$ is complete, every subset
of a $\mu$-null set is measurable. Therefore $\{f > a\}$ is a union of two measurable sets, and $f$ is
measurable.
\end{proof}

\obs Without completeness of $\mathcal{A}$, a.e. convergence does not guarantee measurability: if
$\mu(N) = 0$ but $N$ contains non-measurable subsets, redefining $f$ arbitrarily on $N$
can break measurability.

\obs In general, in any measure space the implications
$$ \text{uniform} \implies \text{pointwise} \implies \text{a.e.}, \qquad
   \text{almost uniform} \implies \text{in measure}, $$
hold, and almost uniform convergence also implies a.e. convergence.
But \emph{no} other implication holds automatically. Typical examples:
\begin{itemize}
    \item $f_n = \mathbf{1}_{[n, n+1]}$ on $\mathbb{R}$ converges pointwise
     to 0, but \emph{not} in measure.
    \item The functions $\mathbf{1}_{[0,1]}, \mathbf{1}_{[0,1/2]},
    \mathbf{1}_{[1/2,1]}, \mathbf{1}_{[0,1/4]}, \dots$ converge to 0 in measure but \emph{not}
    pointwise at any point.
    \item $f_n(x) = x^n$ on $([0, 1], \lambda)$ converges pointwise to $\mathbf{1}_{\{1\}}$
    but \emph{not} uniformly: given $0 < \epsilon < 1$ and $x < 1$, the condition
    $x^n < \epsilon$ requires $n > \ln \epsilon / \ln x$, a bound that blows up as $x \to 1^-$, so
    that no $n_0$ works for all $x$ at once (observe moreover that the pointwise
    limit is not continuous even though each $f_n$ is). It does converge \emph{almost uniformly} to
    $\mathbf{1}_{\{1\}}$: given $\delta > 0$ it suffices to exclude $N = (1 - \delta, 1)$, and on
    $[0, 1 - \delta]$ the convergence to $0$ is uniform because $x^n \le (1 - \delta)^n \to 0$.
\end{itemize}

\subsection{Egoroff's theorem}

\begin{teorema}{Egoroff}{egoroff}
Let $(X, \mathcal{A}, \mu)$ be a finite measure space ($\mu(X) < \infty$) and $f_n \to f$ a.e.
Then:
\begin{enumerate}
    \item $f_n \to f$ almost uniformly (Definition \ref{def:conv_casi_unif}).
    \item As a consequence, $f_n \to f$ in measure.
\end{enumerate}
\end{teorema}

\begin{proof}
Fix $\delta > 0$. For each $p, m \in \mathbb{N}$, let
$$ A_m^{(p)} = \{x \in X : |f_m(x) - f(x)| \ge 1/p\}. $$
It is first shown that, for each $p$,
$$ \mu\!\left(\bigcup_{k \ge m} A_k^{(p)}\right) \xrightarrow{m \to \infty} 0. $$
Indeed, if $x \in \varlimsup_m A_m^{(p)} = \bigcap_m \bigcup_{k \ge m} A_k^{(p)}$, then for
each $m$ there is $k \ge m$ with $|f_k(x) - f(x)| \ge 1/p$, whence $f_k(x) \not\to f(x)$.
Therefore $\varlimsup_m A_m^{(p)} \subseteq \{x \in X : f_n(x) \not\to f(x)\}$, which is $\mu$-null by
hypothesis. Since moreover $\big(\bigcup_{k \ge m} A_k^{(p)}\big)_m$ decreases to
$\varlimsup_m A_m^{(p)}$ and its first term has measure $\le \mu(X) < \infty$, continuity of
the measure gives the announced limit.

(i) is proved. By what was just seen, for each $p \in
\mathbb{N}$ one can choose $m_p$ such that $\mu\big(\bigcup_{k \ge m_p} A_k^{(p)}\big) < \delta/2^p$, and
define
$$ N = \bigcup_{p=1}^\infty \bigcup_{k \ge m_p} A_k^{(p)}, \qquad \text{with } \mu(N) \le
\sum_{p=1}^\infty \frac{\delta}{2^p} = \delta. $$
If $x \in N^c$, then $|f_k(x) - f(x)| < 1/p$ for every $p \in \mathbb{N}$ and every $k \ge m_p$.
Thus, given $\epsilon > 0$ and choosing $p$ with $1/p < \epsilon$, one gets $|f_k(x) - f(x)| < \epsilon$
for every $k \ge m_p$ and every $x \in N^c$. Since $m_p$ does not depend on $x$, the convergence is uniform
on $N^c$.

Finally, (ii) follows from (i): given $\epsilon > 0$ and $\eta > 0$, by almost uniform
convergence there are $N$ with $\mu(N) < \eta$ and $m_0 \in \mathbb{N}$ such that $|f_m - f| < \epsilon$ on
$N^c$ for every $m \ge m_0$. Then $\{|f_m - f| \ge \epsilon\} \subseteq N$ for $m \ge m_0$,
whence $\mu(\{|f_m - f| \ge \epsilon\}) \le \mu(N) < \eta$. Since $\eta$ is arbitrary,
$\mu(\{|f_m - f| \ge \epsilon\}) \to 0$; that is, $f_m \to f$ in measure.
\end{proof}

\obs The $m_p$ that defines $N$ depends on $p$ but \emph{not} on $x$.

\newpage

\section{Integration}

The integral is constructed in three steps: simple, nonnegative, general.

\subsection{Step 1: nonnegative simple functions}

\begin{definicion}{Integral of a nonnegative simple function}{int_simple}
Let $\varphi = \sum_{i=1}^n a_i \mathbf{1}_{A_i}$ be simple with $a_i \ge 0$ and $A_i$ disjoint.
Define
$$ \int_X \varphi \, d\mu = \sum_{i=1}^n a_i \, \mu(A_i), $$
with the convention $0 \cdot \infty = 0$.
\end{definicion}

\obs The definition does not depend on the representation (it suffices to verify it for refinements).

\begin{proposicion}{Linearity and monotonicity on simple functions}{lin_mon_simple}
If $\varphi, \psi$ are nonnegative simple functions and $a, b \ge 0$:
\begin{enumerate}
    \item $\int (a \varphi + b \psi) d\mu = a \int \varphi \, d\mu + b \int \psi \, d\mu$.
    \item If $\varphi \le \psi$, $\int \varphi \, d\mu \le \int \psi \, d\mu$.
\end{enumerate}
Moreover, for $A \in \mathcal{A}$, $\lambda(A) := \int_A \varphi \, d\mu := \int_X \varphi
\mathbf{1}_A \, d\mu$ is a measure on $(X, \mathcal{A})$.
\end{proposicion}

\begin{proof}\leavevmode\\
\textit{(i)} Refining both representations to a common partition, one may suppose that
$\varphi = \sum_i \alpha_i \mathbf{1}_{C_i}$ and $\psi = \sum_i \beta_i \mathbf{1}_{C_i}$ use the
same disjoint partition $\{C_i\}$. Then $a \varphi + b \psi = \sum_i (a \alpha_i + b
\beta_i) \mathbf{1}_{C_i}$, and by definition,
$$ \int (a \varphi + b \psi) d\mu = \sum_i (a \alpha_i + b \beta_i) \mu(C_i) = a \sum_i \alpha_i
\mu(C_i) + b \sum_i \beta_i \mu(C_i) = a \int \varphi \, d\mu + b \int \psi \, d\mu. $$

\textit{(ii)} With the same common partition, $\varphi \le \psi$ implies $\alpha_i \le \beta_i$
on each nonempty $C_i$. Summing with $\mu(C_i) \ge 0$ preserves the inequality:
$\int \varphi \, d\mu \le \int \psi \, d\mu$.

\textit{$\lambda$ is a measure.} $\lambda(\emptyset) = 0$ is direct. For $\sigma$-additivity,
let $(A_n)$ be disjoint and $A = \bigsqcup_n A_n$. With $\varphi = \sum_i a_i \mathbf{1}_{C_i}$
in canonical form,
$$ \lambda(A) = \sum_i a_i \mu(C_i \cap A) = \sum_i a_i \sum_n \mu(C_i \cap A_n) =
\sum_n \sum_i a_i \mu(C_i \cap A_n) = \sum_n \lambda(A_n), $$
where the interchange of sums is legitimate since all terms are nonnegative (Tonelli's
theorem for series).
\end{proof}

\subsection{Step 2: nonnegative measurable functions}

\begin{definicion}{Integral of a nonnegative function}{int_no_neg}
Let $f: X \to [0, \infty]$ be measurable. Define
$$ \int_X f \, d\mu = \sup\left\{ \int_X \varphi \, d\mu : \varphi \text{ nonnegative simple,}
\ 0 \le \varphi \le f \right\}. $$
\end{definicion}

\begin{proposicion}{Basic properties for nonnegative functions}{props_int_pos}
For measurable $f, g \ge 0$:
\begin{enumerate}
    \item If $f \le g$, $\int f \, d\mu \le \int g \, d\mu$.
    \item If $A \subseteq B$ are measurable, $\int_A f \, d\mu \le \int_B f \, d\mu$.
    \item If $A \in \mathcal{A}$, $\nu(A) := \int_A f \, d\mu$ defines a measure on
    $(X, \mathcal{A})$.
\end{enumerate}
\end{proposicion}

\begin{proof}\leavevmode\\
\textit{(i)} Every nonnegative simple $\varphi$ with $0 \le \varphi \le f$ also satisfies
$0 \le \varphi \le g$, so that the supremum defining $\int f$ is taken over a subfamily of the one that
defines $\int g$. Therefore $\int f \le \int g$.

\textit{(ii)} $f \mathbf{1}_A \le f \mathbf{1}_B$ because $\mathbf{1}_A \le \mathbf{1}_B$. By (i),
$\int_A f = \int f \mathbf{1}_A \le \int f \mathbf{1}_B = \int_B f$.

\textit{(iii)} One has $\nu(\emptyset) = 0$. For $\sigma$-additivity, let $(A_n)$ be disjoint and
$A = \bigsqcup_n A_n$.

\emph{Lower bound $\nu(A) \ge \sum_n \nu(A_n)$.} By (ii), for every finite $N$,
$\sum_{n=1}^N \nu(A_n) = \nu\!\left(\bigsqcup_{n=1}^N A_n\right) \le \nu(A)$, where the first
equality follows from finite additivity: for a simple $\varphi \le f$, $\int \varphi
\mathbf{1}_{\bigsqcup_{n=1}^N A_n} = \sum_{n=1}^N \int \varphi \mathbf{1}_{A_n}$ by
$\sigma$-additivity of $\mu$ applied to the measure defined by $\varphi$ (Proposition
\ref{prop:lin_mon_simple}). Letting $N \to \infty$, $\sum_{n=1}^\infty \nu(A_n) \le \nu(A)$.

\emph{Upper bound $\nu(A) \le \sum_n \nu(A_n)$.} Let $\varphi$ be nonnegative simple with
$0 \le \varphi \le f \mathbf{1}_A$. Since $\varphi$ vanishes outside $A$, $\varphi = \sum_n \varphi
\mathbf{1}_{A_n}$ (a pointwise sum where at most one term is nonzero by disjointness). By
$\sigma$-additivity of the measure $E \mapsto \int_E \varphi \, d\mu$ (again Proposition
\ref{prop:lin_mon_simple}),
$$ \int_A \varphi \, d\mu = \sum_n \int_{A_n} \varphi \, d\mu \le \sum_n \nu(A_n), $$
because each $\int_{A_n} \varphi = \int \varphi \mathbf{1}_{A_n} \le \int f \mathbf{1}_{A_n} =
\nu(A_n)$. Taking the supremum over $\varphi$, $\nu(A) \le \sum_n \nu(A_n)$.

Combining both inequalities, $\nu(A) = \sum_n \nu(A_n)$, that is, $\nu$ is $\sigma$-additive.
\end{proof}

\subsection{Step 3: general functions}

\begin{definicion}{Integrable function}{integrable}
Let $f: X \to \mathbb{R}$ be measurable. Decompose $f = f^+ - f^-$, where $f^\pm \ge 0$.
$f$ is said to be \emph{integrable} if $\int f^+ d\mu < \infty$ and $\int f^- d\mu < \infty$, and define
$$ \int_X f \, d\mu = \int_X f^+ d\mu - \int_X f^- d\mu. $$
\end{definicion}

\obs If only one of the two parts is finite, the integral is defined with value in $\mathbb{R} \cup
\{\pm\infty\}$. If both are infinite, the integral is not defined.

\begin{proposicion}{Characterization by $|f|$}{caract_integrable}
$f$ is integrable if and only if $|f|$ is integrable, and in that case
$$ \left| \int_X f \, d\mu \right| \le \int_X |f| \, d\mu \qquad (\text{triangle inequality}). $$
If $|f| \le g$ with $g$ integrable, then $f$ is integrable.
\end{proposicion}

\begin{proof}
$|f| = f^+ + f^-$, so that $|f|$ is integrable iff $f^+$ and $f^-$ are. For the triangle inequality,
$|\int f| = |\int f^+ - \int f^-| \le \int f^+ + \int f^- = \int |f|$. The last claim
follows by monotonicity.
\end{proof}

\begin{proposicion}{Linearity and monotonicity}{lin_int}
If $f, g$ are integrable and $a, b \in \mathbb{R}$:
\begin{enumerate}
    \item $af + bg$ is integrable and $\int (af + bg) d\mu = a \int f \, d\mu + b \int g \, d\mu$.
    \item If $f \le g$, $\int f \, d\mu \le \int g \, d\mu$.
    \item $f = 0$ a.e. $\implies \int f \, d\mu = 0$ (and the converse for $f \ge 0$:
    $\int f \, d\mu = 0 \implies f = 0$ a.e.).
\end{enumerate}
\end{proposicion}

\begin{proof}[Proof of (iii)]
The converse is proved: if $f \ge 0$ and $\int f \, d\mu = 0$, then $f = 0$ a.e. For each $m$,
the set $A_m = \{f \ge 1/m\}$ satisfies $f \ge \tfrac{1}{m} \mathbf{1}_{A_m}$, whence, integrating,
$$ \frac{1}{m} \mu(A_m) \le \int_X f \, d\mu = 0. $$
Hence $\mu(A_m) = 0$ for every $m$, and since $\{f > 0\} = \bigcup_m A_m$, it follows that
$\mu(\{f > 0\}) = 0$.
\end{proof}

\obs For signed functions, $\int_X f \, d\mu = 0$ does not imply $f = 0$ a.e.\ (take an
integrable odd $f$ on $(\mathbb{R}, \lambda)$). The correct version requires all the
restrictions: if
$$ \int_A f \, d\mu = 0 \qquad \text{for every } A \in \mathcal{A}, $$
then $f = 0$ a.e. Indeed, with $A = \{f \ge 0\}$ one gets $f \cdot \mathbf{1}_A = f^+$,
whence $\int_X f^+ \, d\mu = 0$ and $f^+ = 0$ a.e.\ by the previous result; with
$A = \{f \le 0\}$ one gets analogously $f^- = 0$ a.e. This remark is what gives
uniqueness in the Radon-Nikodym theorem and in conditional expectation.

\subsection{Classical inequalities}

\begin{teorema}{Markov's inequality}{markov}
Let $f \ge 0$ be measurable. For every $a > 0$,
$$ \mu(\{f \ge a\}) \le \frac{1}{a} \int_X f \, d\mu. $$
\end{teorema}

\begin{proof}
$f \ge a \cdot \mathbf{1}_{\{f \ge a\}}$, and integrating $\int f \, d\mu \ge a \mu(\{f \ge a\})$.
\end{proof}

\begin{corolario}{Chebyshev's inequality}{chebyshev}
If $X \in L^p$ with $1 \le p < \infty$, for every $a > 0$,
$$ \mu(\{|X| \ge a\}) \le \frac{1}{a^p} \int |X|^p \, d\mu = \frac{\|X\|_p^p}{a^p}. $$
\end{corolario}

\begin{proof}
Apply Markov to $f = |X|^p$ with threshold $a^p$.
\end{proof}

\obs In probability: $P(|X - E(X)| \ge a) \le \mathrm{Var}(X) / a^2$ (Chebyshev with $p = 2$
applied to $X - E(X)$). It is the basis of the weak law of large numbers.

\begin{teorema}{Jensen's inequality (integral)}{jensen_int}
Let $\mu$ be a probability ($\mu(X) = 1$), $f \in L^1(\mu)$, and $\varphi: \mathbb{R} \to \mathbb{R}$
convex with $\varphi \circ f$ integrable. Then
$$ \varphi\!\left(\int_X f \, d\mu\right) \le \int_X \varphi \circ f \, d\mu. $$
\end{teorema}

\begin{proof}
Let $c = \int_X f \, d\mu$. Since $\varphi$ is convex, at each point it admits a supporting line:
there is $m \in \mathbb{R}$ such that
$$ \varphi(y) \ge \varphi(c) + m (y - c) \qquad \forall y \in \mathbb{R}. $$
(The $m$ is any value between the one-sided derivatives $\varphi_-'(c)$ and $\varphi_+'(c)$, which
exist by convexity.) Replacing $y = f(x)$ and integrating with respect to $\mu$,
$$ \int_X \varphi \circ f \, d\mu \ge \int_X \varphi(c) \, d\mu + m \int_X (f - c) \, d\mu
= \varphi(c) \mu(X) + m \left(\int f \, d\mu - c \mu(X)\right). $$
Since $\mu(X) = 1$ and $\int f \, d\mu = c$, the second term is zero and one gets
$\int \varphi \circ f \, d\mu \ge \varphi(c) = \varphi(\int f \, d\mu)$.
\end{proof}

\obs Useful particular cases: $\varphi(t) = t^p$ for $p \ge 1$ gives $|\int f|^p \le \int |f|^p$ (if
$\mu$ is a probability); $\varphi(t) = e^t$ gives $\exp(\int f) \le \int e^f$; $\varphi(t) = -\ln t$
(with $t > 0$) gives $\ln \int f \ge \int \ln f$.

\newpage

\section{Convergence theorems}

When does $\lim_n \int f_n = \int \lim_n f_n$ hold? Not always: let $f_n = n \cdot
\mathbf{1}_{(0, 1/n)}$ on $(\mathbb{R}, \lambda)$. Then $f_n \to 0$ pointwise (with
$x > 0$ fixed, for $n > 1/x$ one has $f_n(x) = 0$; at $x \le 0$ all of them vanish), but
$$ \int_{\mathbb{R}} f_n \, d\lambda = n \cdot \frac{1}{n} = 1 \qquad \text{for every } n, $$
so that $\lim_n \int f_n = 1 \ne 0 = \int \lim_n f_n$: the mass escapes through an increasingly
tall and narrow spike. The same phenomenon in the version for series: with $g_1 = f_1$ and
$g_n = f_n - f_{n-1}$ one has $\sum_{n=1}^m g_n = f_m$, and
$$ \sum_{n=1}^\infty \int_{\mathbb{R}} g_n \, d\lambda \ne \int_{\mathbb{R}}
\sum_{n=1}^\infty g_n \, d\lambda. $$
The theorems of this section give conditions (monotonicity, domination) under which the
interchange is legitimate.

\begin{teorema}{Monotone Convergence (Beppo Levi)}{mct}
Let $(f_n)$ be an increasing sequence of nonnegative measurable functions with $f_n \nearrow f$
pointwise. Then
$$ \int_X f \, d\mu = \lim_{n \to \infty} \int_X f_n \, d\mu. $$
\end{teorema}

\begin{proof}
$f = \lim f_n$ is measurable. Since $f_n \le f$ for every $n$,
$$ \lim_n \int f_n \, d\mu \le \int f \, d\mu. $$
For the other inequality, let $\varphi$ be nonnegative simple with $0 \le \varphi \le f$ and $0 < c < 1$.
Define $A_n = \{x : f_n(x) \ge c \varphi(x)\}$. Since $f_n \nearrow f \ge \varphi > c\varphi$
pointwise on $\{\varphi > 0\}$, the $A_n$ increase to $X$ (or, more precisely, their union covers
$\{\varphi > 0\}$, and outside it there is nothing to prove). Then
$$ \int f_n \, d\mu \ge \int_{A_n} f_n \, d\mu \ge c \int_{A_n} \varphi \, d\mu. $$
Since $\nu(A) := \int_A \varphi \, d\mu$ is a measure, by continuity
$\int_{A_n} \varphi \, d\mu \to \int_X \varphi \, d\mu$. Passing to the limit,
$\lim_n \int f_n \, d\mu \ge c \int \varphi \, d\mu$. Letting $c \to 1$ and then taking the sup over $\varphi$,
one gets $\lim_n \int f_n \, d\mu \ge \int f \, d\mu$.
\end{proof}

\begin{teorema}{Fatou's Lemma}{fatou}
If $(f_n)$ is a sequence of nonnegative measurable functions,
$$ \int_X \varliminf_n f_n \, d\mu \le \varliminf_n \int_X f_n \, d\mu. $$
\end{teorema}

\begin{proof}
Define $g_m = \inf_{n \ge m} f_n$. The sequence $(g_m)$ is increasing and $g_m \nearrow
\varliminf_n f_n$. By MCT,
$$ \int_X \varliminf_n f_n \, d\mu = \lim_m \int_X g_m \, d\mu. $$
Since $g_m \le f_n$ for every $n \ge m$, $\int g_m \le \int f_n$, so that
$\int g_m \le \inf_{n \ge m} \int f_n$. Passing to the limit in $m$,
$\lim_m \int g_m \le \varliminf_n \int f_n$.
\end{proof}

\obs In general the inequality is strict. Example: $f_n = \mathbf{1}_{[n, n+1]}$ on $\mathbb{R}$
satisfies $\varliminf f_n = 0$ but $\int f_n = 1$ for every $n$.

\begin{teorema}{Dominated Convergence (Lebesgue)}{dct}
Let $(f_n)$ be a sequence of measurable functions with $f_n \to f$ a.e., and suppose there is
an integrable function $g$ with $|f_n| \le g$ a.e. for every $n$. Then $f$ is integrable and
$$ \int_X f \, d\mu = \lim_{n \to \infty} \int_X f_n \, d\mu. $$
Moreover, $\int |f_n - f| d\mu \to 0$.
\end{teorema}

\begin{proof}
Passing to a set of full measure, one may suppose convergence and domination at every point.
Since $|f| \le g$, $f$ is integrable. Apply Fatou to the two nonnegative sequences
$g + f_n \ge 0$ and $g - f_n \ge 0$:

\textit{For $g + f_n$:}
$$ \int g + \int f \le \varliminf \int (g + f_n) = \int g + \varliminf \int f_n, $$
so that $\int f \le \varliminf \int f_n$.

\textit{For $g - f_n$:}
$$ \int g - \int f \le \varliminf \int (g - f_n) = \int g - \varlimsup \int f_n, $$
so that $\varlimsup \int f_n \le \int f$.

Combining, $\varlimsup \int f_n \le \int f \le \varliminf \int f_n$, which gives that the limit
exists and equals $\int f$. The claim $\int|f_n - f| d\mu \to 0$ follows by applying DCT to
$|f_n - f|$ with dominating function $2g$.
\end{proof}

\begin{corolario}{Bounded convergence}{conv_acotada}
Let $\mu(X) < \infty$, $f_n \to f$ a.e.\ and $|f_n| \le M$ a.e.\ for every $n$, with $M$
constant. Then $f$ is integrable and $\int_X f \, d\mu = \lim_n \int_X f_n \, d\mu$.
\end{corolario}

\begin{proof}
The constant function $g \equiv M$ is integrable, because $\int_X M \, d\mu = M \mu(X) < \infty$,
and it dominates all the $f_n$. Dominated convergence is applied.
\end{proof}

\begin{ejemplo}{A limit computed with dominated convergence}{ej_tcd}
Compute
$$ L = \lim_{n \to \infty} \int_0^\pi (\sin x)^n \sqrt{x} \, dx. $$
The integrand converges pointwise to $0$ on $[0, \pi] \setminus \{\pi/2\}$ (there
$0 \le \sin x < 1$, so that $(\sin x)^n \to 0$), whereas at
$x = \pi/2$ it equals $\sqrt{\pi/2}$ for every $n$: the pointwise limit is $0$ a.e. The convergence
is not uniform (near $\pi/2$ the integrand stays close to $\sqrt{\pi/2}$), so
that the classical interchange criterion does not apply. But
$$ |(\sin x)^n \sqrt{x}| \le \sqrt{x} =: g(x), \qquad \int_0^\pi \sqrt{x} \, dx =
\tfrac{2}{3} \pi^{3/2} < \infty, $$
and by dominated convergence
$$ L = \int_0^\pi \lim_n (\sin x)^n \sqrt{x} \, dx = \int_0^\pi 0 \, dx = 0. $$
\end{ejemplo}

\subsection{Implications between types of convergence}

There are four types of convergence for sequences of measurable functions: \emph{uniform},
\emph{pointwise} (or a.e.), \emph{in measure} and \emph{in $L^p$}. The elementary
implications are shown, together with two fine classical results: Riesz's theorem (an a.e.\ subsequence from
convergence in measure) and the DCT in measure.

\begin{proposicion}{Convergence in $L^p$ implies convergence in measure}{Lp_implica_medida}
If $f_n \to f$ in $L^p$ with $1 \le p < \infty$, then $f_n \to f$ in measure.
\end{proposicion}

\begin{proof}
By Chebyshev applied to $|f_n - f|$,
$$ \mu(\{|f_n - f| \ge \epsilon\}) \le \frac{\|f_n - f\|_p^p}{\epsilon^p} \to 0. $$
\end{proof}

\begin{teorema}{Riesz: a.e. subsequence}{riesz_subsucesion}
If $f_n \to f$ in measure, there is a subsequence $(f_{n_k})$ with $f_{n_k} \to f$ almost everywhere.
\end{teorema}

\begin{proof}
Construct a subsequence $(f_{n_k})$ with
$$ \mu\!\left(\{|f_{n_k} - f| > 2^{-k}\}\right) < 2^{-k}. $$
This is possible because convergence in measure gives, for each $k$, an $n_k$ satisfying that, and
the $n_k$ can be chosen strictly increasing.

Let $A_k = \{|f_{n_k} - f| > 2^{-k}\}$ and $N = \varlimsup_k A_k$. By Borel-Cantelli I (because
$\sum \mu(A_k) < \infty$), $\mu(N) = 0$. If $x \notin N$, $x$ is in at most finitely many $A_k$,
so
that for large $k$, $|f_{n_k}(x) - f(x)| \le 2^{-k} \to 0$. Hence $f_{n_k} \to f$ outside $N$.
\end{proof}

\obs In general, convergence in measure does \emph{not} imply a.e.\ convergence (the ``walking
waves'' mentioned in the section on convergence). But \emph{a subsequence}
does converge a.e.

\begin{proposicion}{A.e.\ convergence implies convergence in measure if $\mu(X) < \infty$}{ctp_implica_medida}
If $\mu(X) < \infty$ and $f_n \to f$ a.e., then $f_n \to f$ in measure.
\end{proposicion}

\begin{proof}
By Egoroff, for every $\delta > 0$ there is $N$ with $\mu(N) < \delta$ such that $f_n \to f$ uniformly
on $N^c$. Given $\epsilon > 0$, on $N^c$ there is $n_0$ such that $|f_n - f| < \epsilon$ for
$n \ge n_0$,
that is, $\{|f_n - f| \ge \epsilon\} \subseteq N$. Hence $\mu(\{|f_n - f| \ge \epsilon\}) \le \delta$,
and $\delta$ was arbitrary.
\end{proof}

\begin{teorema}{Dominated convergence in measure}{dct_en_medida}
Let $(f_n)$ be a measurable sequence with $f_n \to f$ in measure, and suppose there is $g \ge 0$
integrable with $|f_n| \le g$ a.e.\ for every $n$. Then $f \in L^1$, $\int |f_n - f| d\mu
\to 0$, and in particular $\int f_n d\mu \to \int f \, d\mu$.

More generally, if $|f_n|^p \le g$ with $g$ integrable and $1 \le p < \infty$, then
$f_n \to f$ in $L^p$.
\end{teorema}

\begin{proof}
By Riesz's theorem there is a subsequence $f_{n_k} \to f$ a.e., so that $|f| \le g$ a.e.\
and $f \in L^1$.

For convergence in $L^1$, suppose by contradiction that $\int |f_n - f| d\mu \not\to 0$. There is
$\epsilon_0 > 0$ and a subsequence $(f_{n_j})$ with $\int |f_{n_j} - f| d\mu \ge \epsilon_0$. Since
$f_{n_j} \to f$ in measure, a sub-subsequence $(f_{n_{j_l}})$ with $f_{n_{j_l}} \to f$
a.e. is extracted. By the classical DCT, $\int |f_{n_{j_l}} - f| d\mu \to 0$, which contradicts the lower bound
$\epsilon_0$.

The $L^p$ version is analogous, applying DCT to $|f_n - f|^p \le 2^p (|f_n|^p + |f|^p) \le
2^p (g + |f|^p)$.
\end{proof}

\obs Collecting the implications, in finite measure one has
$$ \text{uniform} \Longrightarrow \text{a.e.} \Longrightarrow \text{in measure}, \qquad
L^p \Longrightarrow \text{in measure}. $$
And with integrable domination, ``in measure'' or ``a.e.'' imply ``in $L^p$''. The other
implications are false in general.

\subsection{Density of simple and continuous functions in \texorpdfstring{$L^p$}{Lp}}

\begin{teorema}{Simple functions are dense in $L^p$}{simples_densas}
Let $1 \le p < \infty$. The set of integrable simple functions (that is, simple $\varphi$
with $\|\varphi\|_p < \infty$) is dense in $L^p(\mu)$.
\end{teorema}

\begin{proof}
Let $f \in L^p$. Decompose $f = f^+ - f^-$. For $f^+ \ge 0$, by the theorem of approximation
by simple functions there are $\varphi_n \nearrow f^+$ with $\varphi_n$ nonnegative simple. Since
$0 \le \varphi_n \le f^+$, $\varphi_n \in L^p$. Moreover $|\varphi_n - f^+|^p \le (f^+)^p \in L^1$,
and it converges to 0 pointwise, so that by DCT $\|\varphi_n - f^+\|_p \to 0$. Analogously for $f^-$.
\end{proof}

\begin{teorema}{Continuous functions are dense in $L^p(\mathbb{R}, \lambda)$}{continuas_densas}
Let $1 \le p < \infty$. The continuous functions with compact support $C_c(\mathbb{R})$ are dense
in $L^p(\mathbb{R}, \lambda)$.
\end{teorema}

\begin{proof}
By the previous theorem it suffices to approximate integrable simple functions by continuous ones. By
linearity, the case $f = \mathbf{1}_A$ with $\lambda(A) < \infty$ suffices.

\textit{Reduction to intervals.} By construction of the Lebesgue measure (via outer
measure with covers by intervals), given $\epsilon > 0$ there is a countable union of
open intervals $U = \bigcup_n (a_n, b_n)$ with $A \subseteq U$ and $\lambda(U \setminus A) <
\epsilon^p / 2$. Truncating to a finite union $V = (a_1, b_1) \cup \cdots \cup (a_N, b_N)$,
one may suppose $\lambda(U \setminus V) < \epsilon^p / 2$. Then $\|\mathbf{1}_A -
\mathbf{1}_V\|_p^p \le \lambda(A \triangle V) \le \lambda(U \setminus A) + \lambda(U \setminus V) <
\epsilon^p$.

\textit{Approximation by continuous functions on each interval.} For $\mathbf{1}_{(a, b)}$, define
$g_\eta: \mathbb{R} \to [0, 1]$ continuous with support in $(a - \eta, b + \eta)$, equal to $1$ on
$[a, b]$ and decreasing linearly to $0$ at the ends. Then $|\mathbf{1}_{(a, b)} - g_\eta|$
vanishes outside $[a - \eta, a] \cup [b, b + \eta]$ and is bounded by 1, so that
$$ \|\mathbf{1}_{(a, b)} - g_\eta\|_p^p \le 2 \eta \to 0 \quad \text{as } \eta \to 0. $$

Combining, $\|\mathbf{1}_V - \sum_k g_\eta^{(k)}\|_p$ becomes arbitrarily small, and the finite
sum of continuous functions with compact support is continuous with compact support. Hence for every
$\epsilon > 0$ there is $g \in C_c(\mathbb{R})$ with $\|\mathbf{1}_A - g\|_p < 2 \epsilon$.
\end{proof}

\obs Density allows proving properties in $L^p$ by verifying them first on simple (or
continuous) functions and then passing to the limit. It is very useful for extending operators: if an operator is
linear and continuous on a dense subset, it extends uniquely to the whole space.

\newpage

\section{The \texorpdfstring{$L^p$}{Lp} spaces}

The theory of the $L^p$ spaces as normed spaces (the Young, Hölder and
Minkowski inequalities, completeness (Riesz-Fischer theorem), the inclusions between the $L^p$ and the
$\ell^p$, the parallelogram law, the dual space and the Riesz representation theorem)
is developed, with proofs, in the \emph{Real analysis} notes. To avoid duplicating it,
only the definition and the statements that the rest of these notes uses are fixed here, and
what does not appear there is proved: the orthogonal decomposition of a Hilbert space, on
which conditional expectation and the proof of Radon-Nikodym rely.

Recall that, for $1 \le p < \infty$,
$$ L^p(\mu) = \left\{ f : X \to \mathbb{R} \text{ measurable} : \int_X |f|^p \, d\mu < \infty
\right\}, \qquad \|f\|_p = \left( \int_X |f|^p \, d\mu \right)^{1/p}, $$
identifying two functions that coincide a.e.\ (without that identification, $\|\cdot\|_p$ is
only a seminorm: $\|f\|_p = 0$ is equivalent to $f = 0$ a.e.). The facts that will be used, all
proved in the \emph{Real analysis} notes:

\begin{enumerate}
    \item \emph{Hölder's inequality:} if $\tfrac{1}{p} + \tfrac{1}{q} = 1$ with
    $1 < p, q < \infty$, $f \in L^p(\mu)$ and $g \in L^q(\mu)$, then $fg \in L^1(\mu)$ and
    $\|fg\|_1 \le \|f\|_p \, \|g\|_q$. For $p = q = 2$ it is the Cauchy-Schwarz inequality.
    \item \emph{Minkowski's inequality:} $\|f + g\|_p \le \|f\|_p + \|g\|_p$. With it,
    $(L^p(\mu), \|\cdot\|_p)$ is a normed space.
    \item \emph{Riesz-Fischer theorem:} $L^p(\mu)$ is complete, that is, a Banach
    space. (The proof runs on the tools of these notes: monotone
    convergence, Fatou's Lemma and dominated convergence.)
    \item $L^2(\mu)$ is a \emph{Hilbert space}: its norm comes from the inner product
    $$ \langle f, g \rangle = \int_X f g \, d\mu, $$
    finite by Cauchy-Schwarz.
    \item \emph{Projection onto a closed convex set:} if $H$ is a Hilbert space,
    $S \subseteq H$ is convex, closed and nonempty, and $x \in H$, there is a unique $s_0 \in S$
    that realizes the distance: $\|x - s_0\| = \inf\{\|x - s\| : s \in S\}$.
    \item \emph{Riesz representation:} every continuous linear functional
    $\varphi: H \to \mathbb{R}$ on a Hilbert space is written $\varphi(x) =
    \langle x, z \rangle$ for a unique $z \in H$. In $L^2(\mu)$: every continuous linear
    functional is ``integrating against a function''. This is exactly the form in which
    the proof of Radon-Nikodym will use it.
\end{enumerate}

\obs If $\mathcal{G} \subseteq \mathcal{F}$ is a sub-$\sigma$-algebra,
$L^p(\mathcal{G}) \subseteq L^p(\mathcal{F})$ is a subspace and, being itself complete
(Riesz-Fischer), it is a \emph{closed} subspace. Conditional expectation will use
this fact.

\begin{teorema}{Orthogonal decomposition}{descomp_ortogonal}
Let $S$ be a closed subspace of a Hilbert space $H$, and let
$S^\perp = \{v \in H : \langle v, s \rangle = 0 \text{ for every } s \in S\}$. Then
$H = S \oplus S^\perp$: every $x \in H$ is written uniquely as $x = s_0 + v$ with
$s_0 \in S$ and $v \in S^\perp$.
\end{teorema}

\begin{proof}
A subspace is convex and, by the projection onto a closed convex set (item (v) above),
there is the $s_0 \in S$ that realizes the distance from $x$ to $S$. It is claimed that
$x - s_0 \in S^\perp$. Let $s \in S$ and $t \in \mathbb{R}$; since $s_0 + t s \in S$,
$$ \|x - s_0\|^2 \le \|x - s_0 - t s\|^2 = \|x - s_0\|^2 - 2 t \langle x - s_0, s \rangle +
t^2 \|s\|^2. $$
The function $\phi(t) = t^2 \|s\|^2 - 2 t \langle x - s_0, s \rangle$ is then $\ge 0$ for
every $t$ and vanishes at $t = 0$: its derivative at $0$ must vanish, that is
$\langle x - s_0, s \rangle = 0$. Since $s \in S$ was arbitrary, $x - s_0 \in S^\perp$ and
$x = s_0 + (x - s_0) \in S + S^\perp$.

The decomposition is unique because $S \cap S^\perp = \{0\}$: if $w$ belongs to both,
$\langle w, w \rangle = 0$ and $w = 0$.
\end{proof}

The operator $P: H \to S$, $P(x) = s_0$, is called the \emph{orthogonal projection} onto $S$.

\newpage

\section{Product measures: Fubini-Tonelli}

Given two measures, how does one integrate a function on the product? The theorems of Tonelli and
Fubini allow interchanging the order of integration.

\begin{definicion}{Product measure}{medida_producto}
Let $(X, \mathcal{A}, \mu)$ and $(Y, \mathcal{B}, \nu)$ be $\sigma$-finite measure spaces. On the
product $X \times Y$ define:
\begin{itemize}
    \item The product $\sigma$-algebra $\mathcal{A} \otimes \mathcal{B} = \sigma(\{A \times B :
    A \in \mathcal{A}, B \in \mathcal{B}\})$.
    \item The product measure $\mu \otimes \nu$ is the unique measure on $\mathcal{A} \otimes
    \mathcal{B}$ with $(\mu \otimes \nu)(A \times B) = \mu(A) \nu(B)$ for measurable rectangles.
\end{itemize}
\end{definicion}

\obs In particular, $\lambda \otimes \lambda$ on $\mathbb{R}^2$ extends to the
$n$-dimensional Lebesgue measure, which assigns to a rectangle $[a, b] \times [c, d]$ the area $(b-a)(d-c)$.

\begin{teorema}{Tonelli (nonnegative functions)}{tonelli}
Let $f: X \times Y \to [0, \infty]$ be measurable with respect to $\mathcal{A} \otimes \mathcal{B}$, with
$\mu$ and $\nu$ $\sigma$-finite. Then:
\begin{enumerate}
    \item For each $x \in X$, $y \mapsto f(x, y)$ is $\mathcal{B}$-measurable (and symmetrically
    for $y$).
    \item $x \mapsto \int_Y f(x, y) \, d\nu(y)$ is $\mathcal{A}$-measurable (and symmetrically).
    \item The equality
    $$ \int_{X \times Y} f \, d(\mu \otimes \nu) = \int_X \!\left( \int_Y f(x, y) \, d\nu(y)
    \right) d\mu(x) = \int_Y \!\left( \int_X f(x, y) \, d\mu(x) \right) d\nu(y) $$
    holds.
\end{enumerate}
\end{teorema}

\begin{proof}
\textit{Step 1: measurability of the sections.} Let $\mathcal{M} \subseteq \mathcal{A} \otimes
\mathcal{B}$ be the class of sets $E$ such that the section $E_x = \{y : (x, y) \in E\}$ is
$\mathcal{B}$-measurable for every $x$. If $E = A \times B$ is a measurable rectangle, $E_x = B$ if
$x \in A$, $\emptyset$ otherwise, in both cases measurable. $\mathcal{M}$ is a $\sigma$-algebra
(directly: $(E^c)_x = (E_x)^c$, $(\bigcup E_n)_x = \bigcup (E_n)_x$). Since it contains the
rectangles, $\mathcal{M} \supseteq \sigma(\text{rectangles}) = \mathcal{A} \otimes
\mathcal{B}$. By approximation by simple functions, for measurable $f$ so is $f(x, \cdot)$.

\textit{Step 2: case $f = \mathbf{1}_E$, $E \in \mathcal{A} \otimes \mathcal{B}$.} Define
$g_E(x) = \nu(E_x)$. It must be seen that $g_E$ is $\mathcal{A}$-measurable and $\int_X g_E \, d\mu = (\mu
\otimes \nu)(E)$. For rectangles $E = A \times B$, $g_E(x) = \nu(B) \mathbf{1}_A(x)$, measurable,
and $\int_X g_E \, d\mu = \mu(A) \nu(B) = (\mu \otimes \nu)(E)$.

The class $\mathcal{D}$ of the $E$ where this property holds is a \emph{monotone class}: closed
under increasing unions and decreasing intersections (assuming $\sigma$-finiteness in order to
interchange limits with integrals via MCT) and under differences of comparable sets. By the
monotone class lemma (or by Dynkin's $\pi$-$\lambda$ lemma), $\mathcal{D}$ contains the
$\sigma$-algebra generated by the rectangles, which is $\mathcal{A} \otimes \mathcal{B}$.

\textit{Step 3: case $f$ nonnegative simple.} By linearity over the terms
$c_i \mathbf{1}_{E_i}$.

\textit{Step 4: case $f \ge 0$ general.} Take nonnegative simple $\varphi_n \nearrow f$
(approximation by simple functions). For each $x$, $\varphi_n(x, \cdot) \nearrow f(x, \cdot)$. By MCT on
$Y$,
$$ \int_Y \varphi_n(x, y) \, d\nu(y) \nearrow \int_Y f(x, y) \, d\nu(y) =: F(x). $$
$F$ is a pointwise limit of $\mathcal{A}$-measurable functions (by Step 3), so that $F$ is
$\mathcal{A}$-measurable. Applying MCT again on $X$,
$$ \int_X F \, d\mu = \lim_n \int_X \!\int_Y \varphi_n \, d\nu \, d\mu = \lim_n \int_{X \times Y}
\varphi_n \, d(\mu \otimes \nu) = \int_{X \times Y} f \, d(\mu \otimes \nu), $$
using the simple case. Symmetrically for the other order.
\end{proof}

\begin{teorema}{Fubini (integrable functions)}{fubini}
Let $f: X \times Y \to \mathbb{R}$ be measurable. If $f$ is integrable with respect to $\mu \otimes \nu$
(or if, by Tonelli, $\int_X \int_Y |f| \, d\nu \, d\mu < \infty$), then:
\begin{enumerate}
    \item For $\mu$-a.e.\ $x$, the function $y \mapsto f(x, y)$ is in $L^1(\nu)$.
    \item The function $x \mapsto \int_Y f(x, y) \, d\nu(y)$, defined $\mu$-a.e., is in
    $L^1(\mu)$.
    \item Tonelli's equality holds for $f$, with finite integrals.
\end{enumerate}
\end{teorema}

\begin{proof}
Decompose $f = f^+ - f^-$ with $f^\pm \ge 0$ measurable. By Tonelli applied to $|f| = f^+ + f^-$,
$$ \int_X \!\int_Y |f| \, d\nu \, d\mu = \int_{X \times Y} |f| \, d(\mu \otimes \nu) < \infty. $$

For the left-hand side to be finite, $\int_Y |f(x, y)| \, d\nu(y) < \infty$ for $\mu$-a.e.\
$x$, that is, $f(x, \cdot) \in L^1(\nu)$ $\mu$-a.e. At those points $x$ define
$$ F(x) = \int_Y f(x, y) \, d\nu(y) = \int_Y f^+(x, y) \, d\nu(y) - \int_Y f^-(x, y) \, d\nu(y), $$
which is the difference of two measurable functions that are finite $\mu$-a.e., so that $F$ is
$\mathcal{A}$-measurable and, by Tonelli applied to $f^\pm$,
$$ \int_X F \, d\mu = \int_{X \times Y} f^+ \, d(\mu \otimes \nu) - \int_{X \times Y} f^- \,
d(\mu \otimes \nu) = \int_{X \times Y} f \, d(\mu \otimes \nu), $$
and all the integrals are finite because $f \in L^1(\mu \otimes \nu)$.
\end{proof}

\obs In practice, the workflow is: \textit{first} verify with Tonelli (on $|f|$) that the
double integral is finite, \textit{then} apply Fubini to interchange the order with $f$.
Skipping the check gives wrong results (the classical computation
$\int_0^1 \int_0^1 (x^2 - y^2)/(x^2 + y^2)^2 \, dx \, dy = \pi/4$ vs.\ $\int \int = -\pi/4$ happens
because $|f|$ is not integrable).

\newpage

\section{Complete measure spaces}

\begin{definicion}{Complete space}{esp_completo}
$(X, \mathcal{A}, \mu)$ is \emph{complete} if every subset of a set of measure zero is
measurable (in particular, it has measure zero):
$$ E \in \mathcal{A}, \ \mu(E) = 0, \ F \subseteq E \implies F \in \mathcal{A}. $$
\end{definicion}

\obs Every measure space can be \emph{completed}: given $(X, \mathcal{A}, \mu)$, define
$\overline{\mathcal{A}} = \{A \cup N : A \in \mathcal{A}, N \subseteq E \text{ with } \mu(E) = 0\}$
and extend $\mu$ to $\overline{\mathcal{A}}$. For example, $(\mathbb{R}, \mathcal{B}, \lambda)$
is not complete, but its completion is $(\mathbb{R}, \mathcal{L}, \lambda)$.

\obs In complete spaces, if $f$ is measurable and $g = f$ a.e., then $g$ is also measurable.
This makes it easier to handle ``a.e.'' properties.

\newpage

\section{Signed measures and the Radon-Nikodym theorem}

\begin{definicion}{Signed measure}{medida_signo}
A function $\nu: \mathcal{A} \to [-\infty, +\infty]$ is a \emph{signed measure} if:
\begin{enumerate}
    \item $\nu(\emptyset) = 0$.
    \item $\nu$ takes at most one of the values $\pm\infty$.
    \item $\sigma$-additivity: $\nu(\bigsqcup A_n) = \sum \nu(A_n)$ (with absolute convergence
    when
    the left-hand side is finite).
\end{enumerate}
\end{definicion}

\begin{definicion}{Absolute continuity}{abs_cont}
Let $\mu, \nu$ be measures on $(X, \mathcal{A})$. $\nu$ is said to be \emph{absolutely continuous}
with respect to $\mu$ (written $\nu \ll \mu$) if
$$ \mu(E) = 0 \implies \nu(E) = 0 \qquad \text{for every } E \in \mathcal{A}. $$
\end{definicion}

\begin{definicion}{Mutually singular}{singular}
$\mu$ and $\nu$ are \emph{mutually singular} (written $\mu \perp \nu$) if there are
disjoint $A, B \in \mathcal{A}$ with $A \cup B = X$ and $\mu(B) = 0 = \nu(A)$.
\end{definicion}

\begin{ejemplo}{Absolutely continuous measures via integration}{}
If $f \ge 0$ is measurable and $\mu$ is a measure, $\nu(A) = \int_A f \, d\mu$ is a measure on the
same measurable space, and it is absolutely continuous with respect to $\mu$: if $\mu(A) = 0$, then
$\int_A f \, d\mu = 0$ (the integral over a $\mu$-null set is zero). The
Radon-Nikodym theorem, later on, will say that \emph{every} absolutely continuous measure (under
reasonable hypotheses) has this form.
\end{ejemplo}

\subsection{Lebesgue decomposition + Radon-Nikodym}

The theorem proved next is the combined version: \emph{every} measure $\nu$ decomposes into an
absolutely continuous part and a singular part with respect to $\mu$. When $\nu \ll \mu$, the singular
part disappears and only the derivative remains.

\begin{teorema}{Lebesgue-Radon-Nikodym}{radon_nikodym}
Let $\mu$ and $\nu$ be $\sigma$-finite measures on $(X, \mathcal{A})$. Then:
\begin{enumerate}
    \item There are unique measures $\nu_a, \nu_\perp$ such that
    $$ \nu = \nu_a + \nu_\perp, \qquad \nu_a \ll \mu, \qquad \nu_\perp \perp \mu. $$
    \item There is a measurable function $h: X \to [0, \infty)$, unique $\mu$-a.e., such that
    $$ \nu_a(E) = \int_E h \, d\mu \qquad \forall E \in \mathcal{A}. $$
\end{enumerate}
The function $h$ is called the \emph{Radon-Nikodym derivative} of $\nu_a$ with respect to $\mu$, and it is
denoted $h = \dfrac{d\nu_a}{d\mu}$. In particular, when $\nu \ll \mu$, $\nu_\perp = 0$ and
$\nu(E) = \int_E h \, d\mu$.
\end{teorema}

\begin{proof}[Proof (von Neumann)]
The case of finite measures is done; the passage to $\sigma$-finite ones is by decomposition into finite
parts.

\textit{Auxiliary functional.} On $L^2(\mu + \nu)$ define
$$ \varphi(f) = \int_X f \, d\nu. $$
It is linear and continuous: by Cauchy-Schwarz (with $g \equiv 1$),
$$ |\varphi(f)| \le \int_X |f| \, d\nu \le \nu(X)^{1/2} \, \|f\|_{L^2(\nu)} \le \nu(X)^{1/2} \,
\|f\|_{L^2(\mu+\nu)}. $$
By the Riesz representation theorem applied to $L^2(\mu + \nu)$, there is $g \in
L^2(\mu + \nu)$ such that
\begin{equation}\label{eq:rn1}
\int_X f \, d\nu = \int_X f g \, d(\mu + \nu) \qquad \forall f \in L^2(\mu + \nu).
\end{equation}

\textit{$0 \le g \le 1$ almost everywhere.} Taking $f = \mathbf{1}_A$ in \eqref{eq:rn1},
$$ \nu(A) = \int_A g \, d(\mu + \nu), $$
whence $0 \le \int_A g \, d(\mu + \nu) \le (\mu + \nu)(A)$. This forces $0 \le g \le 1$
$(\mu + \nu)$-a.e.

\textit{Decomposition.} Let $N = \{x : g(x) = 1\}$ and $M = N^c = \{x : g(x) < 1\}$. Define
$$ \nu_\perp(E) = \nu(E \cap N), \qquad \nu_a(E) = \nu(E \cap M). $$
It is proved that $\nu_\perp \perp \mu$ and $\nu_a \ll \mu$:
\begin{itemize}
    \item $\mu(N) = 0$: taking $f = \mathbf{1}_N$ in \eqref{eq:rn1},
    $\nu(N) = \int_N g \, d(\mu+\nu)
    = (\mu+\nu)(N) = \mu(N) + \nu(N)$, whence $\mu(N) = 0$. Thus $\nu_\perp$ and $\mu$ are
    mutually singular (separated by $N$).
    \item If $\mu(E) = 0$, it must hold that $\nu_a(E) = \nu(E \cap M) = 0$. Rearranging \eqref{eq:rn1}
    (for $f = \mathbf{1}_{E \cap M}$),
    $$ \int_{E \cap M} (1 - g) \, d\nu = \int_{E \cap M} g \, d\mu = 0. $$
    On $M$, $1 - g > 0$, so that necessarily $\nu(E \cap M) = 0$.
\end{itemize}

\textit{Derivative.} Iterating \eqref{eq:rn1} with $f_n = \mathbf{1}_E (1 + g + g^2 + \cdots + g^n)$ and
subtracting, one obtains
$$ \int_E (1 - g^{n+1}) \, d\nu = \int_E (g + g^2 + \cdots + g^{n+1}) \, d\mu. $$
On $M$, $g^{n+1} \to 0$ and $g + g^2 + \cdots \to g/(1-g)$. By monotone convergence,
$$ \nu_a(E) = \nu(E \cap M) = \int_{E \cap M} \frac{g}{1 - g} \, d\mu = \int_E h \, d\mu, $$
where $h = \mathbf{1}_M \cdot g/(1 - g)$.

Uniqueness follows in the standard way: if two decompositions coexisted, their difference would be at the same time
absolutely continuous and singular, which forces it to be zero.
\end{proof}

\subsection{Application: the Fundamental Theorem of Calculus}

Let $F: \mathbb{R} \to \mathbb{R}$ be nondecreasing and right-continuous, and $\lambda_F$ its
Lebesgue-Stieltjes measure. If $\lambda_F \ll \lambda$ (which happens, for example, when $F$ is
\emph{absolutely continuous} in the classical sense), by Radon-Nikodym there is
$f \in L^1_{\mathrm{loc}}(\mathbb{R})$ such that
$$ F(b) - F(a) = \lambda_F((a, b]) = \int_{(a, b]} f \, d\lambda = \int_a^b f(x) \, dx. $$
The function $f = d\lambda_F / d\lambda$ plays the role of the derivative of $F$ in the
distributional sense. This is the \emph{general} version of the Fundamental Theorem of Calculus.

\newpage

\section{Conditional Expectation}

Let $(\Omega, \mathcal{F}, P)$ be a probability space. Recall elementary conditional
probability: given events $A, B$ with $P(B) > 0$,
$$ P(A \mid B) = \frac{P(A \cap B)}{P(B)}. $$
This is a ``basic'' version that only allows conditioning on an event. What is sought now
is to condition with respect to an \emph{entire $\sigma$-algebra} $\mathcal{G} \subseteq \mathcal{F}$,
which represents the available information.

\obs Two direct consequences of the elementary definition. If $A$ and $B$ are independent
($P(A \cap B) = P(A) P(B)$), then $P(A \mid B) = P(A)$: knowing that $B$ occurred provides no
information about $A$. And if $\{B_1, B_2, \dots\}$ is a partition of $\Omega$ with
$P(B_i) > 0$, the sets $\{A \cap B_i\}$ partition $A$, so that
$$ P(A) = \sum_i P(A \cap B_i) = \sum_i P(A \mid B_i) \, P(B_i)
\qquad \text{(total probability)}, $$
and solving from $P(B_k \mid A) \, P(A) = P(A \cap B_k) = P(A \mid B_k) \, P(B_k)$ one obtains
\textbf{Bayes' formula}:
$$ P(B_k \mid A) = \frac{P(A \mid B_k) \, P(B_k)}{\sum_i P(A \mid B_i) \, P(B_i)}. $$
Conditional expectation with respect to a $\sigma$-algebra will generalize this weighted average
over the elements of a partition.

\subsection{Motivation}

\begin{ejemplo}{Three coin tosses}{ejemplo_monedas}
A coin is tossed 3 times. $\Omega$ has 8 elements $\{HHH, HHT, \dots, TTT\}$, and $\mathcal{F}
= \mathcal{P}(\Omega)$. Define nested $\sigma$-algebras that ``gradually reveal''
information:
\begin{align*}
\widetilde{\mathcal{F}}_0 &= \{\emptyset, \Omega\} && \text{(nothing is known)}, \\
\widetilde{\mathcal{F}}_1 &= \sigma(A) = \{\emptyset, \Omega, A, A^c\} && \text{where } A
= \{\text{the 1st toss is }H\}, \\
\widetilde{\mathcal{F}}_2 &= \sigma(B_1, B_2, B_3, B_4) && \text{(according to the first 2 tosses)}, \\
\widetilde{\mathcal{F}}_3 &= \mathcal{P}(\Omega) && \text{(everything is known)}.
\end{align*}
Each $\widetilde{\mathcal{F}}_k$ has more information than $\widetilde{\mathcal{F}}_{k-1}$.
\end{ejemplo}

\begin{definicion}{Filtration}{filtracion}
A \emph{filtration} on $(\Omega, \mathcal{F}, P)$ is an increasing family of
$\sigma$-algebras $(\widetilde{\mathcal{F}}_t)_{t \in T} \subseteq \mathcal{F}$:
$$ s \le t \implies \widetilde{\mathcal{F}}_s \subseteq \widetilde{\mathcal{F}}_t. $$
\end{definicion}

\begin{ejemplo}{Conditional expectation in the three tosses}{}
Continuing the previous example, let $X: \Omega \to \mathbb{R}$ with $X(\omega) = $ ``number of
heads in $\omega$''. The expectation is
$$ E(X) = 0 \cdot \tfrac{1}{8} + 1 \cdot \tfrac{3}{8} + 2 \cdot \tfrac{3}{8} + 3 \cdot \tfrac{1}{8}
= \tfrac{3}{2}. $$
If the 1\textsuperscript{st} toss has already been seen (that is, $\widetilde{\mathcal{F}}_1$ is known),
one \emph{expects} on average:
$$ E(X \mid \widetilde{\mathcal{F}}_1)(\omega) = \begin{cases}
2, & \text{if the 1\textsuperscript{st} was $H$} \ (= 1 + E(\text{2 remaining tosses})), \\
1, & \text{if the 1\textsuperscript{st} was $T$}.
\end{cases} $$
Note that $E(X \mid \widetilde{\mathcal{F}}_1)$ is an $\widetilde{\mathcal{F}}_1$-measurable
\emph{random variable} (it takes constant values on $A$ and on $A^c$).
\end{ejemplo}

\obs In the example measurability can be verified by hand:
$\{E(X \mid \widetilde{\mathcal{F}}_1) \ge 2\} = A \in \widetilde{\mathcal{F}}_1$, whereas
$\{X \ge 2\} = \{HHH, HHT, HTH, THH\} \notin \widetilde{\mathcal{F}}_1$. The conditional
expectation is $\widetilde{\mathcal{F}}_1$-measurable; $X$ is not.

\obs (\emph{Definition on partitions.}) When $\mathcal{G} = \sigma(G_1, G_2, \dots)$ with
$\{G_i\}$ a partition of $\Omega$ and $P(G_i) > 0$, there is an explicit formula. A
$\mathcal{G}$-measurable variable is constant on each atom $G_i$, and the conditional expectation turns out
to be the average of $X$ over the atom where $\omega$ falls:
$$ E(X \mid \mathcal{G})(\omega) = \frac{1}{P(G_i)} \int_{G_i} X \, dP \qquad \text{if }
\omega \in G_i. $$
The characteristic property is checked: every $A \in \mathcal{G}$ is a union of atoms,
$A = \bigcup_{i \in M} G_i$, and
$$ \int_A E(X \mid \mathcal{G}) \, dP = \sum_{i \in M} \left( \frac{1}{P(G_i)} \int_{G_i} X \,
dP \right) P(G_i) = \sum_{i \in M} \int_{G_i} X \, dP = \int_A X \, dP. $$
The general definition (which follows) extends this to $\sigma$-algebras without atoms, where the
local average can no longer be written by dividing by $P(G_i)$.

\subsection{General definition}

\begin{definicion}{Conditional expectation}{esp_cond}
Let $X \in L^1(\Omega, \mathcal{F}, P)$ and let $\mathcal{G} \subseteq \mathcal{F}$ be a
$\sigma$-algebra. The \emph{conditional expectation} of $X$ given $\mathcal{G}$, denoted
$E(X \mid \mathcal{G})$, is the unique $\mathcal{G}$-measurable random variable (up to a.e.\ equality)
such that
$$ \int_A E(X \mid \mathcal{G}) \, dP = \int_A X \, dP \qquad \forall A \in \mathcal{G}. $$
\end{definicion}

\begin{teorema}{Existence (via Radon-Nikodym)}{exist_esp_cond}
Under the previous hypotheses, $E(X \mid \mathcal{G})$ exists and is unique $P$-a.e.
\end{teorema}

\begin{proof}
Suppose first $X \ge 0$. Define $\lambda: \mathcal{G} \to [0, \infty]$ by
$$ \lambda(A) = \int_A X \, dP. $$
$\lambda$ is a measure on $(\Omega, \mathcal{G})$. It is absolutely continuous with respect to
$P|_{\mathcal{G}}$: if $P(A) = 0$, then $\int_A X \, dP = 0$. By Radon-Nikodym, there is
$h \in L^1(\mathcal{G})$ with
$$ \int_A X \, dP = \lambda(A) = \int_A h \, dP \qquad \forall A \in \mathcal{G}. $$
That $h$ is $E(X \mid \mathcal{G})$. For general $X$, decompose $X = X^+ - X^-$ and subtract.
\end{proof}

\subsection{Properties}

\begin{proposicion}{Properties of conditional expectation}{props_esp_cond}
Let $X, Y \in L^1$ and $\mathcal{G}, \mathcal{H} \subseteq \mathcal{F}$ be $\sigma$-algebras.
\begin{enumerate}
    \item $E(aX + bY \mid \mathcal{G}) = a E(X \mid \mathcal{G}) + b E(Y \mid \mathcal{G})$.
    \item If $Y$ is $\mathcal{G}$-measurable (and $XY \in L^1$),
    $$ E(XY \mid \mathcal{G}) = Y \cdot E(X \mid \mathcal{G}). $$
    \item If $X$ is independent of $\mathcal{G}$ (that is, $X^{-1}(B) \perp A$ for every $A \in
    \mathcal{G}$ and Borel set $B$),
    $$ E(X \mid \mathcal{G}) = E(X). $$
    \item If $\mathcal{H} \subseteq \mathcal{G}$,
    $$ E(E(X \mid \mathcal{G}) \mid \mathcal{H}) = E(X \mid \mathcal{H}). $$
\end{enumerate}
\end{proposicion}

\begin{proof}\leavevmode\\
\textit{(i)} For every $A \in \mathcal{G}$,
\begin{align*}
\int_A (aE(X \mid \mathcal{G}) + b E(Y \mid \mathcal{G})) \, dP
&= a \int_A E(X \mid \mathcal{G}) \, dP + b \int_A E(Y \mid \mathcal{G}) \, dP \\
&= a \int_A X \, dP + b \int_A Y \, dP = \int_A (aX + bY) \, dP.
\end{align*}
By uniqueness, $E(aX + bY \mid \mathcal{G}) = a E(X \mid \mathcal{G}) + b E(Y \mid \mathcal{G})$.

\textit{(ii)} It is proved by increasing approximation.
\begin{itemize}
    \item For $Y = \mathbf{1}_C$ with $C \in \mathcal{G}$: for every $A \in \mathcal{G}$,
    $A \cap C \in \mathcal{G}$, and
    $$ \int_A \mathbf{1}_C \cdot E(X \mid \mathcal{G}) dP = \int_{A \cap C} E(X \mid \mathcal{G}) dP
    = \int_{A \cap C} X dP = \int_A \mathbf{1}_C X dP. $$
    By uniqueness, $E(\mathbf{1}_C X \mid \mathcal{G}) = \mathbf{1}_C E(X \mid \mathcal{G})$.

    \item For $Y = \sum_i c_i \mathbf{1}_{C_i}$ simple and $\mathcal{G}$-measurable: by linearity from
    (i).

    \item For $Y \ge 0$ $\mathcal{G}$-measurable: take nonnegative simple $Y_n \nearrow Y$. If
    $X \ge 0$, by monotone convergence $\int_A Y_n X dP \to \int_A YX dP$ and $\int_A Y_n
    E(X \mid \mathcal{G}) dP \to \int_A Y E(X \mid \mathcal{G}) dP$. For general $X$ decompose
    $X = X^+ - X^-$.

    \item For general $Y$: $Y = Y^+ - Y^-$ and the previous step is applied to each part.
\end{itemize}

\textit{(iii)} For $A \in \mathcal{G}$, $\mathbf{1}_A$ and $X$ are independent,
so that $E(\mathbf{1}_A X) = P(A) E(X)$:
$$ \int_A X \, dP = E(\mathbf{1}_A X) = E(\mathbf{1}_A) E(X) = P(A) E(X) = \int_A E(X) \, dP. $$
$E(X)$ is constant (in particular $\mathcal{G}$-measurable), so that by uniqueness
$E(X \mid \mathcal{G}) = E(X)$.

\textit{(iv)} Let $Z = E(X \mid \mathcal{G})$ and $W = E(X \mid \mathcal{H})$. For
$A \in \mathcal{H} \subseteq \mathcal{G}$,
$$ \int_A Z \, dP = \int_A X \, dP = \int_A W \, dP. $$
Since $W$ is $\mathcal{H}$-measurable and the equality holds for every $A \in \mathcal{H}$, by uniqueness of
$E(Z \mid \mathcal{H})$ one has $E(Z \mid \mathcal{H}) = W$, that is $E(E(X \mid \mathcal{G})
\mid \mathcal{H}) = E(X \mid \mathcal{H})$.
\end{proof}

\subsection{Interpretation as a projection}

When $X \in L^2$, conditional expectation admits a geometric characterization:

\begin{teorema}{$E(\cdot \mid \mathcal{G})$ is an orthogonal projection}{proy_L2}
For $X \in L^2(\mathcal{F})$, $E(X \mid \mathcal{G})$ is the orthogonal projection of $X$ onto the
closed subspace $L^2(\mathcal{G}) \subseteq L^2(\mathcal{F})$. That is, $E(X \mid \mathcal{G})$
is the unique $\mathcal{G}$-measurable element that minimizes
$$ \|X - Y\|_{L^2}^2 = E((X - Y)^2) \qquad \text{among } Y \in L^2(\mathcal{G}). $$
\end{teorema}

\begin{proof}
Since $L^2(\mathcal{G})$ is a closed subspace of $L^2(\mathcal{F})$ (Hilbert), there is the
projection $\widehat{X}$ characterized by $X - \widehat{X} \perp L^2(\mathcal{G})$, that is
$$ \int_\Omega (X - \widehat{X}) Y \, dP = 0 \qquad \forall Y \in L^2(\mathcal{G}). $$
Taking $Y = \mathbf{1}_A$ with $A \in \mathcal{G}$, $\int_A X \, dP = \int_A \widehat{X} \, dP$, and
since $\widehat{X}$ is $\mathcal{G}$-measurable, by uniqueness $\widehat{X} = E(X \mid \mathcal{G})$.
\end{proof}

\obs This explains the name ``best predictor of $X$ given the information $\mathcal{G}$'':
$E(X \mid \mathcal{G})$ minimizes the mean squared error among all
$\mathcal{G}$-measurable functions.

\begin{ejemplo}{The three tosses as nested spaces}{ej_L2_anidados}
In the example of the three tosses, $L^2(\widetilde{\mathcal{F}}_k)$ is the space of
functions that are constant on the atoms of $\widetilde{\mathcal{F}}_k$; identifying each
variable with the vector of its $8$ values $(X(HHH), X(HHT), X(HTH), X(HTT), X(THH), X(THT),
X(TTH), X(TTT))$,
$$ \underbrace{L^2(\widetilde{\mathcal{F}}_0)}_{\dim 1} \subset
\underbrace{L^2(\widetilde{\mathcal{F}}_1)}_{\dim 2} \subset
\underbrace{L^2(\widetilde{\mathcal{F}}_2)}_{\dim 4} \subset
\underbrace{L^2(\widetilde{\mathcal{F}}_3)}_{\dim 8}. $$
The variable $X$ = ``number of heads'' is the vector $(3, 2, 2, 1, 2, 1, 1, 0)$, and the
successive orthogonal projections average over increasingly coarse atoms:
$$ (3,2,2,1,2,1,1,0) \ \mapsto \ \left(\tfrac{5}{2}, \tfrac{5}{2}, \tfrac{3}{2}, \tfrac{3}{2},
\tfrac{3}{2}, \tfrac{3}{2}, \tfrac{1}{2}, \tfrac{1}{2}\right) \ \mapsto \ (2,2,2,2,1,1,1,1)
\ \mapsto \ \left(\tfrac{3}{2}, \dots, \tfrac{3}{2}\right), $$
which are $E(X \mid \widetilde{\mathcal{F}}_2)$, $E(X \mid \widetilde{\mathcal{F}}_1)$ and
$E(X \mid \widetilde{\mathcal{F}}_0) = E(X) = 3/2$. Composing projections
($L^2(\mathcal{F}) \to L^2(\mathcal{G}) \to L^2(\mathcal{H})$ with $\mathcal{H} \subseteq
\mathcal{G}$) is the geometric reading of the tower property.
\end{ejemplo}

\obs Another characterization, in terms of covariances: for $X \in L^2$,
$E(X \mid \mathcal{G})$ is the unique $\mathcal{G}$-measurable variable with $E(E(X \mid
\mathcal{G})) = E(X)$ such that
$$ \mathrm{Cov}(E(X \mid \mathcal{G}), Y) = \mathrm{Cov}(X, Y) \qquad \text{for every } Y \in
L^2(\mathcal{G}). $$
Indeed, subtracting,
$$ \mathrm{Cov}(X, Y) - \mathrm{Cov}(E(X \mid \mathcal{G}), Y) = \big\langle X - E(X \mid
\mathcal{G}), \, Y - E(Y) \big\rangle = 0, $$
because $X - E(X \mid \mathcal{G})$ is orthogonal to $L^2(\mathcal{G})$ and $Y - E(Y) \in
L^2(\mathcal{G})$. Conditional expectation preserves all the correlation with what is observable.

\subsection{Conditional Jensen inequality}

\begin{teorema}{Conditional Jensen}{jensen_cond}
Let $X \in L^1$ and $\varphi: \mathbb{R} \to \mathbb{R}$ convex with $\varphi(X) \in L^1$. Then
$$ \varphi(E(X \mid \mathcal{G})) \le E(\varphi(X) \mid \mathcal{G}) \quad P\text{-a.e.} $$
\end{teorema}

\begin{proof}
The idea is the same as in the integral Jensen inequality: by convexity there is a countable family of
supporting lines $\varphi(y) = \sup_n (a_n y + b_n)$. For each fixed $n$, by linearity and monotonicity
of conditional expectation,
$$ E(\varphi(X) \mid \mathcal{G}) \ge E(a_n X + b_n \mid \mathcal{G}) = a_n E(X \mid \mathcal{G})
+ b_n. $$
Taking the supremum over $n$, $E(\varphi(X) \mid \mathcal{G}) \ge \varphi(E(X \mid \mathcal{G}))$.
\end{proof}

\obs Particular cases: $\varphi(t) = |t|^p$ with $p \ge 1$ gives
$$ |E(X \mid \mathcal{G})|^p \le E(|X|^p \mid \mathcal{G}), $$
and taking expectations, $\|E(X \mid \mathcal{G})\|_p \le \|X\|_p$. Thus, the conditional expectation
operator is \emph{contractive} on $L^p$.

\newpage

\section{Martingales}

Martingales formalize the idea of a \emph{fair game}: the conditional expectation of the future
value, knowing everything that has happened so far, coincides with the current value.

\subsection{Adapted processes and stopping times}

\begin{definicion}{Adapted process}{adaptado}
Let $(\widetilde{\mathcal{F}}_t)_{t \in T}$ be a filtration. A stochastic process
$(X_t)_{t \in T}$ is \emph{adapted} to the filtration if each $X_t$ is
$\widetilde{\mathcal{F}}_t$-measurable.
\end{definicion}

\obs If no filtration is given, the \emph{natural filtration}
$\widetilde{\mathcal{F}}_t = \sigma(X_s : s \le t)$ can be used, which is the smallest $\sigma$-algebra that
makes $(X_s)_{s \le t}$ measurable. By construction, $(X_t)$ is adapted to its natural filtration.

\begin{definicion}{Stopping time}{stopping}
Let $(\widetilde{\mathcal{F}}_t)_{t \in T}$ be a filtration. A \emph{stopping time} is a random
variable $\tau: \Omega \to T$ such that for every $t$,
$$ \{\tau \le t\} \in \widetilde{\mathcal{F}}_t. $$
In discrete time $T = \mathbb{N}$ it is equivalent to requiring $\{\tau = n\} \in \widetilde{\mathcal{F}}_n$.
\end{definicion}

\obs Intuitively, $\tau$ is an instant \emph{decidable with the information up to that moment}.
No knowledge of the future is required.

\begin{ejemplo}{Common stopping times}{}
The following are stopping times: playing a fixed number of times ($\tau \equiv N$ constant), ``the first time
the event $A$ appears'' ($\tau = \inf\{n : \omega_n \in A\}$), or ``the $k$-th visit to
$A$''. ``The instant before the maximum'' is \emph{not} one, because knowing that the next step
will be the maximum requires knowing the future. Moreover, $\max(\tau_1, \tau_2)$ and $\min(\tau_1,
\tau_2)$ are stopping times: $\{\max(\tau_1, \tau_2) \le t\} = \{\tau_1 \le t\} \cap \{\tau_2 \le
t\}$ and $\{\min(\tau_1, \tau_2) \le t\} = \{\tau_1 \le t\} \cup \{\tau_2 \le t\}$ are in
$\widetilde{\mathcal{F}}_t$.
\end{ejemplo}

\begin{teorema}{Wald's Lemma}{wald}
Let $(X_i)$ be i.i.d.\ and integrable, $S_n = \sum_{i=1}^n X_i$, and $\tau$ a stopping time
with respect to $\widetilde{\mathcal{F}}_n = \sigma(X_1, \dots, X_n)$ with $E(\tau) < \infty$.
Then
$$ E(S_\tau) = E(X_1) \, E(\tau). $$
\end{teorema}

\begin{proof}
Suppose first $X_i \ge 0$. Write
$$ S_\tau = \sum_{n=1}^\infty X_n \, \mathbf{1}_{\{\tau \ge n\}}. $$
The event $\{\tau \ge n\} = \{\tau \le n - 1\}^c \in \widetilde{\mathcal{F}}_{n-1}$ is
independent of $X_n$, so that
$E(X_n \mathbf{1}_{\{\tau \ge n\}}) = E(X_n) \, P(\tau \ge n) = E(X_1) \, P(\tau \ge n)$. By
monotone convergence (the summands are nonnegative),
$$ E(S_\tau) = \sum_{n=1}^\infty E(X_n \mathbf{1}_{\{\tau \ge n\}}) = E(X_1)
\sum_{n=1}^\infty P(\tau \ge n) = E(X_1) \, E(\tau), $$
using the identity $E(\tau) = \sum_{n \ge 1} P(\tau \ge n)$, valid for variables with values
in $\mathbb{N}$ (it is seen by writing $\tau = \sum_{n \ge 1} \mathbf{1}_{\{\tau \ge n\}}$ and
interchanging expectation and sum by monotone convergence). For general $X_i$ the above is applied
to the parts $X_i^+$ and $X_i^-$ and subtracted; both expectations are finite since
$E|X_1| < \infty$ and $E(\tau) < \infty$.
\end{proof}

\begin{ejemplo}{Compound dice}{ej_wald}
A die is rolled and $N$ comes up; then $N$ dice are rolled and added: $S = X_1 + \cdots + X_N$. Here
$N$ is independent of the $X_i$ (a particular case of stopping time), and
$$ E(S) = E(X_1) \, E(N) = \frac{7}{2} \cdot \frac{7}{2} = \frac{49}{4} = 12.25, $$
with no need to condition by hand on the $6$ values of $N$.
\end{ejemplo}

\subsection{Definition of a martingale}

\begin{definicion}{Martingale}{martingala}
A process $(X_n)_{n \in \mathbb{N}}$ adapted to $(\widetilde{\mathcal{F}}_n)$ is a
\emph{martingale} if each $X_n$ is integrable and
$$ E(X_n \mid \widetilde{\mathcal{F}}_{n-1}) = X_{n-1} \qquad \forall n \ge 1. $$
\end{definicion}

\obs Equivalently, $E(X_n \mid \widetilde{\mathcal{F}}_m) = X_m$ for every $m \le n$ (by
the tower property).

\obs As a consequence, $E(X_n) = E(X_0)$ for every $n$. All the variables of a martingale
have the same expectation.

\begin{definicion}{Submartingale and supermartingale}{submartingala}
An adapted integrable process $(X_n)$ is a \emph{submartingale} if
$$ E(X_n \mid \widetilde{\mathcal{F}}_{n-1}) \ge X_{n-1} \qquad \forall n \ge 1, $$
and a \emph{supermartingale} if the opposite inequality holds,
$E(X_n \mid \widetilde{\mathcal{F}}_{n-1}) \le X_{n-1}$.
\end{definicion}

\obs In the language of games: a martingale is a fair game, a submartingale one
favorable to the player and a supermartingale an unfavorable one. A process is a martingale if and
only if it is both a sub- and a supermartingale, and the expectations $E(X_n)$ turn out to be increasing in the
sub case and decreasing in the super case. If $(X_n)$ is a martingale and $\varphi$ is convex with
$\varphi(X_n)$ integrable, then $(\varphi(X_n))$ is a submartingale: by conditional Jensen
(Theorem \ref{teo:jensen_cond}),
$E(\varphi(X_n) \mid \widetilde{\mathcal{F}}_{n-1}) \ge
\varphi(E(X_n \mid \widetilde{\mathcal{F}}_{n-1})) = \varphi(X_{n-1})$. In particular
$(|X_n|)$ and $(|X_n|^p)$, with $p \ge 1$, are submartingales.

\begin{ejemplo}{Constants and averages}{}
If $X \in L^1$ and one defines $X_n = E(X \mid \widetilde{\mathcal{F}}_n)$, then $(X_n)$ is a
martingale with respect to $(\widetilde{\mathcal{F}}_n)$. By the tower property,
$E(X_n \mid \widetilde{\mathcal{F}}_{n-1}) = E(E(X \mid \widetilde{\mathcal{F}}_n) \mid
\widetilde{\mathcal{F}}_{n-1}) = E(X \mid \widetilde{\mathcal{F}}_{n-1}) = X_{n-1}$.
\end{ejemplo}

\begin{ejemplo}{Symmetric random walk}{}
Let $(X_n)_{n \ge 1}$ be i.i.d.\ with $E(X_n) = 0$ and $S_n = \sum_{k=1}^n X_k$ (with $S_0 = 0$).
With respect to the natural filtration,
$$ E(S_n \mid \widetilde{\mathcal{F}}_{n-1}) = E(S_{n-1} + X_n \mid \widetilde{\mathcal{F}}_{n-1})
= S_{n-1} + E(X_n) = S_{n-1}, $$
using that $S_{n-1}$ is $\widetilde{\mathcal{F}}_{n-1}$-measurable and $X_n$ is independent. Thus
$(S_n)$ is a martingale.
\end{ejemplo}

\begin{ejemplo}{$S_n^2 - n$}{}
With $X_k$ Rademacher ($\pm 1$ with probability $1/2$), $(S_n^2 - n)$ is a martingale. Computation:
\begin{align*}
E(S_n^2 \mid \widetilde{\mathcal{F}}_{n-1}) &= E((S_{n-1} + X_n)^2 \mid
\widetilde{\mathcal{F}}_{n-1})
= S_{n-1}^2 + 2 S_{n-1} E(X_n) + E(X_n^2) \\
&= S_{n-1}^2 + 0 + 1.
\end{align*}
Subtracting $n$ from both sides, $E(S_n^2 - n \mid \widetilde{\mathcal{F}}_{n-1}) = S_{n-1}^2 - (n-1)$.
\end{ejemplo}

\subsection{Predictable strategy: martingale transform}

\begin{definicion}{Predictable process}{previsible}
$(H_n)_{n \ge 1}$ is \emph{predictable} (with respect to $(\widetilde{\mathcal{F}}_n)$) if each $H_n$ is
$\widetilde{\mathcal{F}}_{n-1}$-measurable. Intuitively, $H_n$ is decided \emph{before} step $n$.
\end{definicion}

\begin{proposicion}{Martingale transform}{transformada}
If $(X_n)$ is a martingale and $(H_n)$ is predictable and bounded, then
$$ Y_n = \sum_{k=1}^n H_k (X_k - X_{k-1}) $$
is a martingale.
\end{proposicion}

\begin{proof}
$Y_n - Y_{n-1} = H_n (X_n - X_{n-1})$. Since $H_n$ is $\widetilde{\mathcal{F}}_{n-1}$-measurable,
$$ E(Y_n - Y_{n-1} \mid \widetilde{\mathcal{F}}_{n-1}) = H_n E(X_n - X_{n-1} \mid
\widetilde{\mathcal{F}}_{n-1}) = H_n \cdot 0 = 0. $$
\end{proof}

\obs The interpretation: in a fair game, no predictable strategy (one that only uses the
past) can turn the game into a favorable one. This justifies mathematically that ``there is no
infallible system'' in a fair game.

\obs A famous strategy: \emph{double or nothing} (the classical martingale, where the name comes from).
$H_n = 2^{n-1}$ if $X_{k-1} = -k$ for every $k < n$, and $H_n = 0$ if the player has already won. The accumulated
gain is 1 when the player finally wins, but ruin arrives earlier with significant probability.

\newpage

\section{Doob's theorems}

\subsection{Optional Stopping}

\begin{teorema}{Optional Stopping (bounded version)}{optional_stopping}
Let $(X_n)$ be a martingale and $\tau$ a stopping time with $\tau \le M$ bounded. Then
$$ E(X_\tau) = E(X_0). $$
\end{teorema}

\begin{proof}
Define $X_n^\tau = X_{\tau \wedge n}$, the process ``stopped'' at $\tau$. It is verified that
$(X_n^\tau)$ is a martingale (transform by the predictable strategy
$H_n = \mathbf{1}_{\tau \ge n}$,
which is $\widetilde{\mathcal{F}}_{n-1}$-measurable because $\{\tau \ge n\} = \{\tau \le n-1\}^c$).
Thus $E(X_\tau) = E(X_M^\tau) = E(X_0^\tau) = E(X_0)$.
\end{proof}

\begin{teorema}{Optional Stopping (uniformly integrable martingale)}{optional_stopping_UI}
If $(X_n)$ is a uniformly integrable martingale (there is $X_\infty$ with $X_n \to X_\infty$ in
$L^1$) and $\tau$ is a stopping time that is finite a.e., then
$$ E(X_\tau) = E(X_0). $$
\end{teorema}

\begin{proof}
Let $\tau_n = \tau \wedge n$, which is a bounded stopping time. By the bounded version,
$E(X_{\tau_n}) = E(X_0)$ for every $n$.

Since $\tau < \infty$ a.e., $X_{\tau_n} \to X_\tau$ a.e.\ (eventually $\tau_n = \tau$). It remains
to see that $\int X_{\tau_n} dP \to \int X_\tau dP$. Write
$$ X_{\tau_n} = X_\tau \mathbf{1}_{\tau \le n} + X_n \mathbf{1}_{\tau > n}. $$
Since $(X_n)$ is uniformly integrable, in particular $\{X_n \mathbf{1}_{\tau > n}\}$ is UI (their
absolute values are dominated by those of $(X_n)$), and since $P(\tau > n) \to 0$, that term
tends to zero in $L^1$. The first term $X_\tau \mathbf{1}_{\tau \le n}$ converges to $X_\tau$ in
$L^1$ by dominated convergence (with dominating function $|X_\tau| \in L^1$, which also follows from UI:
by Fatou and UI, $E|X_\tau| \le \varliminf E|X_{\tau_n}| < \infty$).

Letting $n \to \infty$, $E(X_0) = E(X_{\tau_n}) \to E(X_\tau)$.
\end{proof}

\subsection{Doob's inequality}

\begin{teorema}{Doob's maximal inequality}{doob_ineq}
Let $(X_n)$ be a martingale with $X_n \ge 0$, and $c > 0$. Then
$$ P\left(\max_{k \le n} X_k \ge c\right) \le \frac{E(X_n)}{c}. $$
\end{teorema}

\begin{proof}
Define the stopping time
$$ \tau = \min\{m : \max_{k \le m} X_k \ge c\} \wedge n. $$
By construction $\tau \le n$ and $\{\tau < n\} = \{\max_k X_k \ge c\}$ (if the maximum reaches $c$
before $n$). By Optional Stopping,
$$ E(X_n) = E(X_\tau) = E(X_\tau \mathbf{1}_{\tau < n}) + E(X_\tau \mathbf{1}_{\tau = n})
\ge c \cdot P(\tau < n), $$
because on $\{\tau < n\}$ one has $X_\tau \ge c$ and $X_\tau \mathbf{1}_{\tau = n} \ge 0$. Solving,
$P(\max X_k \ge c) = P(\tau < n) \le E(X_n)/c$.
\end{proof}

\begin{corolario}{Gaussian tail of the supremum of Brownian motion}{cola_browniano}
Let $(W_t)$ be a Brownian motion. For all $a, T > 0$,
$$ P\!\left(\sup_{t \le T} W_t \ge a\right) \le e^{-a^2/(2T)}. $$
\end{corolario}

\begin{proof}
Fix $\lambda > 0$. If $Z \sim \mathcal{N}(0, t)$, completing the square one computes
$E(e^{\lambda Z}) = e^{\lambda^2 t/2}$, so that the process
$$ M_t = e^{\lambda W_t - \lambda^2 t / 2} $$
is a positive martingale: for $s < t$, since $W_t - W_s \sim \mathcal{N}(0, t - s)$ is
independent of the past,
$$ E(M_t \mid \mathcal{F}_s) = M_s \, e^{-\lambda^2 (t - s)/2} \, E\big(e^{\lambda (W_t -
W_s)}\big) = M_s \, e^{-\lambda^2 (t-s)/2} \, e^{\lambda^2 (t-s)/2} = M_s. $$
Since $e^x$ is increasing and $e^{-\lambda^2 t/2} \ge e^{-\lambda^2 T/2}$ for $t \le T$, the
inclusion of events
$$ \left\{ \sup_{t \le T} W_t \ge a \right\} \subseteq \left\{ \sup_{t \le T} M_t \ge
e^{\lambda a - \lambda^2 T/2} \right\} $$
holds, and Doob's maximal inequality applied to the positive martingale $(M_t)$ (the
continuous version is obtained from the discrete one by taking suprema over dyadic times and using the
continuity of the paths), with $E(M_T) = E(M_0) = 1$, gives
$$ P\!\left(\sup_{t \le T} W_t \ge a\right) \le \frac{E(M_T)}{e^{\lambda a - \lambda^2 T/2}} =
e^{\frac{\lambda^2 T}{2} - \lambda a}. $$
The right-hand side is minimized at $\lambda = a/T$ (setting the derivative $\lambda T - a$ of the
exponent to zero), and substituting yields $e^{a^2/(2T) - a^2/T} = e^{-a^2/(2T)}$.
\end{proof}

\subsection{Convergence of martingales}

\begin{teorema}{Martingale convergence (Doob)}{conv_martingalas}
Let $(X_n)$ be a martingale with $\sup_n E|X_n| < \infty$ (bounded in $L^1$). Then there is
$X \in L^1$ with $X_n \to X$ almost everywhere.
\end{teorema}

\begin{proof}
\emph{Upcrossings} technique. Given $a < b$, define the number of complete upcrossings
$U_n^{[a,b]}(\omega)$ as the number of times that $(X_k(\omega))_{k=0}^n$ crosses from below
$a$ to above $b$. More precisely, let $\tau_0 = 0$ and, recursively,
$$ \tau_{2j+1} = \min\{k > \tau_{2j} : X_k < a\}, \qquad \tau_{2j+2} = \min\{k > \tau_{2j+1} : X_k
> b\}, $$
with the convention $\min \emptyset = \infty$. Then $U_n^{[a,b]} = \max\{j : \tau_{2j} \le n\}$.

\textit{Doob's upcrossing inequality.} For any martingale $(X_n)$,
$$ (b - a) \, E(U_n^{[a,b]}) \le E((X_n - a)^+). $$
Idea of the proof: the \emph{strategy} $H_k = \mathbf{1}_{\tau_{2j+1} < k \le \tau_{2j+2} \text{
for some } j}$ (buy after falling below $a$, sell after rising above $b$) is
predictable and bounded. The accumulated gain $Y_n = \sum_{k=1}^n H_k(X_k - X_{k-1})$ is
nonnegative because each complete upcrossing from $a$ to $b$ contributes at least $b - a$; more
formally, $Y_n \ge (b - a) U_n^{[a,b]} - (X_n - a)^+$ (the remaining term absorbs a
possible incomplete upcrossing at the end). Since $Y_n$ is a martingale (transform by a predictable process),
$E(Y_n) = 0$, and one gets $(b - a) E(U_n) \le E((X_n - a)^+)$.

\textit{Conclusion.} If $\sup_n E|X_n| < \infty$, then $E((X_n - a)^+) \le |a| + E|X_n|$
is
uniformly bounded, so that
$$ E(U_\infty^{[a,b]}) = \lim_n E(U_n^{[a,b]}) \le \frac{|a| + \sup_n E|X_n|}{b - a} < \infty. $$
In particular, $U_\infty^{[a,b]} < \infty$ almost everywhere.

Let $\Lambda_{a,b} = \{\omega : \varliminf X_n(\omega) < a < b < \varlimsup X_n(\omega)\}$.
On $\Lambda_{a,b}$, the sequence $X_n(\omega)$ crosses the interval $[a, b]$ infinitely many times, so
that $U_\infty^{[a,b]}(\omega) = \infty$. Hence $P(\Lambda_{a,b}) = 0$.

Taking a countable union over rationals $a < b$,
$$ P\!\left(\bigcup_{a < b \in \mathbb{Q}} \Lambda_{a,b}\right) = 0. $$
Outside that null set, $\varliminf X_n = \varlimsup X_n$, that is, the limit
$X = \lim X_n$ exists in $\mathbb{R} \cup \{\pm\infty\}$. By Fatou's Lemma applied to $|X_n|$,
$E|X| \le \varliminf E|X_n| \le \sup E|X_n| < \infty$, so that $X \in L^1$ (in particular,
finite a.e.).
\end{proof}

\obs The convergence in the theorem is a.e., \emph{not} in $L^1$. Counterexample: let $(Z_k)$ be
i.i.d.\ with $P(Z_k = 0) = P(Z_k = 2) = 1/2$ (so that $E(Z_k) = 1$) and $M_n = \prod_{k=1}^n Z_k$.
Then $(M_n)$ is a martingale with respect to $\widetilde{\mathcal{F}}_n = \sigma(Z_1, \dots,
Z_n)$:
$$ E(M_{n+1} \mid \widetilde{\mathcal{F}}_n) = M_n \, E(Z_{n+1} \mid
\widetilde{\mathcal{F}}_n) = M_n \, E(Z_{n+1}) = M_n, $$
and $E(M_n) = 1$ for every $n$, with $\sup_n E|M_n| = 1 < \infty$. By the theorem, $M_n$ converges
a.e., and the limit is $0$: as soon as a $Z_k = 0$ comes up (which happens with probability $1$),
the product vanishes forever. But then
$$ E|M_n - 0| = E(M_n) = 1 \not\to 0: $$
there is no convergence in $L^1$, and the expectation is lost in the limit
($1 = \lim_n E(M_n) \ne E(\lim_n M_n) = 0$).

\begin{teorema}{$L^p$ convergence, $p > 1$}{conv_Lp_martingalas}
If $\sup_n E|X_n|^p < \infty$ with $p > 1$, then $X_n \to X$ also in $L^p$, and
$X_n = E(X \mid \widetilde{\mathcal{F}}_n)$.
\end{teorema}

\begin{proof}
Since $L^p \subseteq L^1$ in finite measure ($P$ is a probability), $\sup E|X_n| < \infty$ and by the
previous theorem $X_n \to X$ a.e.\ with $X \in L^1$.

\textit{Uniform bound in $L^p$ via Doob's inequality.} Applying Doob's maximal inequality to
the positive martingale $|X_n|^p$ (which is a submartingale by conditional Jensen with
$\varphi(t) = |t|^p$), one obtains the \emph{$L^p$ maximal inequality}:
$$ E\!\left(\sup_{k \le n} |X_k|^p\right) \le \left(\frac{p}{p - 1}\right)^p \sup_{k \le n}
E|X_k|^p. $$
Letting $n \to \infty$ and by monotone convergence, $E(\sup_n |X_n|^p) \le (p/(p-1))^p \sup_n
E|X_n|^p < \infty$. Call $M = \sup_n |X_n|$, which satisfies $M \in L^p$.

\textit{Convergence in $L^p$.} One has $|X_n - X|^p \le 2^p (|X_n|^p + |X|^p) \le 2^p (M^p + M^p)
= 2^{p+1} M^p$ (because $|X| \le M$ in the limit). Since $X_n \to X$ a.e.\ and $|X_n - X|^p \le
2^{p+1} M^p \in L^1$, by dominated convergence $E|X_n - X|^p \to 0$.

\textit{Identity $X_n = E(X \mid \widetilde{\mathcal{F}}_n)$.} For $A \in
\widetilde{\mathcal{F}}_n$ and $m \ge n$, by the martingale property $\int_A X_n dP = \int_A X_m
dP$. By convergence in $L^1$, $\int_A X_m dP \to \int_A X dP$. Hence $\int_A X_n dP = \int_A X
dP$, and since $X_n$ is $\widetilde{\mathcal{F}}_n$-measurable, by uniqueness $X_n = E(X \mid
\widetilde{\mathcal{F}}_n)$.
\end{proof}

\newpage

\section{Brownian Motion}

The passage from discrete to continuous time leads to \emph{Brownian motion}, the limit of
rescaled random walks. Before the axiomatic definition, the random walk is studied in
some detail: its explicit computations are what motivate the properties of the limit
process.

\subsection{The simple random walk}

Let $(X_n)_{n \ge 1}$ be i.i.d.\ with $P(X_n = 1) = P(X_n = -1) = 1/2$, and $S_n = X_1 + \cdots +
X_n$ (with $S_0 = 0$): a particle starts at the origin and at each step moves one unit to the
right or to the left according to a fair coin. The basic computations: $E(X_i) = 0$ and
$\mathrm{Var}(X_i) = E(X_i^2) = 1$, whence
$$ E(S_n) = 0, \qquad \mathrm{Var}(S_n) = E(S_n^2) = \sum_{i,j} E(X_i X_j) = \sum_{i=1}^n
E(X_i^2) = n, $$
using that for $i \ne j$, $E(X_i X_j) = E(X_i) E(X_j) = 0$ by independence. Moreover, for
$m < n$, $S_n - S_m = \sum_{i = m+1}^n X_i$ is another random walk, independent of
$S_1, \dots, S_m$, with $\mathrm{Var}(S_n - S_m) = n - m$.

\begin{ejemplo}{Distribution of $S_n$ and returns to the origin}{ej_retorno_origen}
After $n$ steps, $S_n = a$ requires $k$ steps $+1$ and $n - k$ steps $-1$ with $a = k - (n - k) =
2k - n$; in particular $a$ and $n$ have the same parity, and $k = (n + a)/2$. Of the $2^n$
equiprobable paths, the favorable ones are the choices of where to place the $+1$ steps:
$$ P(S_n = a) = \frac{1}{2^n} \binom{n}{\frac{n+a}{2}}. $$
In particular, the probability of being at the origin after $2m$ steps is $P(S_{2m} = 0) =
\binom{2m}{m} / 2^{2m}$, which with Stirling's formula $n! \approx n^n e^{-n} \sqrt{2 \pi n}$
is estimated as
$$ P(S_{2m} = 0) = \frac{(2m)!}{(m!)^2 \, 2^{2m}} \approx
\frac{(2m)^{2m} e^{-2m} \sqrt{4 \pi m}}{m^{2m} e^{-m} \cdot m^{m} e^{-m} \cdot (2 \pi m) \cdot
2^{2m}} \cdot \frac{m^{m}}{m^{m}} = \frac{\sqrt{4 \pi m}}{2 \pi m} = \frac{1}{\sqrt{\pi m}}. $$
It decays, but very slowly: if $V = \sum_{m \ge 1} \mathbf{1}_{\{S_{2m} = 0\}}$ is the number
of returns to the origin,
$$ E(V) = \sum_{m=1}^\infty P(S_{2m} = 0) \approx \sum_{m=1}^\infty \frac{1}{\sqrt{\pi m}} =
\infty. $$
(Borel-Cantelli (II) is not applicable here: the events $\{S_{2m} = 0\}$ are not
independent.)
\end{ejemplo}

\obs $E(V) = \infty$ does not by itself imply $V = \infty$ a.e.: a variable can have
infinite expectation and be finite with probability $1$. This is the \emph{St.\ Petersburg
paradox}: a coin is tossed until the first tails, and if $k$ heads came up before, one wins
$2^k$ pesos; the expected gain is
$$ E(G) = \sum_{k=0}^\infty 2^k \cdot \frac{1}{2^{k+1}} = \sum_{k=0}^\infty \frac{1}{2} =
\infty, $$
but $G$ is finite with probability $1$. For the random walk the correct argument is a different one. Let
$\rho$ be the probability of ever returning to the origin. Each time the walk returns, the
process starts afresh (the later steps are independent and have the same
distribution), so that $P(V \ge k) = \rho^k$. If $\rho < 1$,
$$ E(V) = \sum_{k \ge 1} P(V \ge k) = \sum_{k \ge 1} \rho^k = \frac{\rho}{1 - \rho} < \infty, $$
contradicting $E(V) = \infty$. Therefore $\rho = 1$: the random walk returns to the origin (and,
iterating, visits it \emph{infinitely many times}) with probability $1$.

\begin{ejemplo}{The Galton board and the Gaussian bell curve}{ej_galton}
On a Galton board, small balls fall through $n$ levels of pegs; each ball performs a random
walk of $n$ steps, and the column over position $a$ accumulates a fraction of balls
proportional to $\binom{n}{(n+a)/2}$. Measuring each column in units of the central column,
the height over $a$ is
$$ H_a = \binom{n}{\frac{n+a}{2}} \Big/ \binom{n}{\frac{n}{2}} =
\frac{(n/2)! \, (n/2)!}{\left(\frac{n+a}{2}\right)! \left(\frac{n-a}{2}\right)!}. $$
Write $\rho = n/2$ and $c = a/2$; cancelling factorials and taking logarithms,
$$ H_a = \frac{(\rho - c + 1) \cdots \rho}{(\rho + 1) \cdots (\rho + c)}, \quad
\ln H_a = \sum_{k = \rho - c + 1}^{\rho} \ln k - \sum_{k = \rho + 1}^{\rho + c} \ln k
\approx \int_{\rho - c}^{\rho} \ln x \, dx - \int_{\rho}^{\rho + c} \ln x \, dx. $$
With the antiderivative $x \ln x - x$,
$$ \ln H_a \approx 2 \rho \ln \rho - (\rho - c) \ln(\rho - c) - (\rho + c) \ln(\rho + c). $$
The Taylor polynomial of order $2$ of $f(x) = x \ln x$ at $\rho$ (with $f'(x) = \ln x + 1$ and
$f''(x) = 1/x$) gives $f(\rho \pm c) \approx \rho \ln \rho \pm (\ln \rho + 1) c +
\frac{c^2}{2\rho}$; on substituting, the terms of order $0$ and $1$ cancel and one gets
$$ \ln H_a \approx -\frac{c^2}{\rho} = -\frac{a^2/4}{n/2} = -\frac{a^2}{2n}
= -\frac{1}{2} \left( \frac{a}{\sqrt{n}} \right)^{\!2}, \quad \text{that is} \quad
H_a \approx e^{-x^2/2} \quad \text{with } x = \frac{a}{\sqrt{n}}. $$
Rescaling the horizontal axis by $\sqrt{n}$, the profiles fit the Gaussian bell curve:
this is the \emph{de Moivre-Laplace} theorem,
$$ \frac{S_n}{\sqrt{n}} \xrightarrow{d} \mathcal{N}(0, 1), \qquad
P\!\left( \alpha \le \frac{S_n}{\sqrt{n}} \le \beta \right) \approx \frac{1}{\sqrt{2 \pi}}
\int_\alpha^\beta e^{-x^2/2} \, dx. $$
\end{ejemplo}

\subsection{From the random walk to Brownian motion}

Heuristically, Brownian motion is a random walk with time steps of length
$\Delta t = 1/N$ and space steps of length $\epsilon$, in the limit $N \to \infty$: one
defines $W_{n/N} = \epsilon S_n$ and normalizes by requiring $\mathrm{Var}(W_1) = 1$. Since
$N$ steps are taken to reach $t = 1$,
$$ 1 = E(W_1^2) = \epsilon^2 E(S_N^2) = \epsilon^2 N \implies \epsilon = \frac{1}{\sqrt{N}}. $$
With that choice, for $t = n/N$,
$$ W_t = \frac{S_n}{\sqrt{N}} = \frac{S_n}{\sqrt{n}} \cdot \sqrt{\frac{n}{N}} =
\frac{S_n}{\sqrt{n}} \cdot \sqrt{t}, $$
and since $S_n / \sqrt{n} \sim \mathcal{N}(0, 1)$ for large $n$ (de Moivre-Laplace), one expects
$W_t \sim \mathcal{N}(0, t)$. From $\mathrm{Var}(S_n - S_m) = n - m$ and from the independence of the
stretches of the random walk one also expects $W_t - W_s \sim \mathcal{N}(0, t - s)$ and independent
increments. These properties, together with the continuity of the paths, are
exactly the definition.

\subsection{Definition and properties}

\begin{definicion}{Brownian Motion}{browniano}
A \emph{Brownian motion} (Wiener process) is a stochastic process
$(W_t)_{t \ge 0}$ that satisfies:
\begin{enumerate}
    \item $W_0 = 0$ almost everywhere.
    \item For $0 \le s < t$, $W_t - W_s \sim \mathcal{N}(0, t - s)$.
    \item Independent increments: for $0 \le t_0 < t_1 < \cdots < t_n$, the variables
    $W_{t_1} - W_{t_0}, W_{t_2} - W_{t_1}, \dots, W_{t_n} - W_{t_{n-1}}$ are independent.
    \item $t \mapsto W_t$ is continuous almost everywhere.
\end{enumerate}
\end{definicion}

\obs Existence (Wiener, 1923): such a process exists. The construction uses, for example,
Kolmogorov's extension or a random series expansion.

\obs The paths $t \mapsto W_t(\omega)$ are continuous but, almost surely, they are not
differentiable at any point. The heuristics for both facts lies in the scale of the
increments: since $\mathrm{Var}(W_{t+h} - W_t) = h$, the typical increment is
$|W_{t+h} - W_t| \approx \sqrt{h}$, which tends to $0$ (continuity) but whose difference
quotient $\sqrt{h}/h = 1/\sqrt{h} \to \infty$ (non-differentiability). Moreover: a derivative at
$t_0$ could be estimated with past increments ($r < t_0$) and would determine the nearby future
increments, but $W_s - W_{t_0}$ is \emph{independent} of $\{W_r : r \le t_0\}$: the
process cannot remember a slope. This is what makes classical calculus fail and
forces the construction of the Itô integral.

\obs Transformations that preserve Brownian motion: if $(W_t)$ is a Brownian motion, so are
\begin{gather*}
(c \, W_{t/c^2})_{t \ge 0} \quad \text{(self-similarity)}, \qquad
(W_{t+s} - W_s)_{t \ge 0} \quad \text{(restart)}, \\
(t \, W_{1/t})_{t \ge 0} \quad \text{(time inversion)}.
\end{gather*}
With drift and scale, $(\mu t + \sigma W_t)$ is called \emph{Brownian motion with drift} and
$(C e^{\mu t + \sigma W_t})$ \emph{geometric Brownian motion}.

\subsection{Quadratic variation}

\begin{teorema}{Quadratic variation of Brownian motion}{var_cuad}
Let $a = t_0 < t_1 < \cdots < t_m = b$ be a partition of $[a, b]$ with mesh
$\delta = \max_k (t_{k+1} - t_k)$. Define
$$ Q_m = \sum_{k=1}^m (W_{t_k} - W_{t_{k-1}})^2. $$
Then $Q_m \xrightarrow{L^2} b - a$ as $\delta \to 0$.
\end{teorema}

\begin{proof}
Call $\Delta_k = W_{t_k} - W_{t_{k-1}}$ and $\Delta t_k = t_k - t_{k-1}$. By definition of
Brownian motion, $\Delta_k \sim \mathcal{N}(0, \Delta t_k)$ and the $\Delta_k$ are independent.

\textit{Expectation.} $E(\Delta_k^2) = \mathrm{Var}(\Delta_k) = \Delta t_k$, so that
$$ E(Q_m) = \sum_{k=1}^m E(\Delta_k^2) = \sum_{k=1}^m \Delta t_k = b - a. $$

\textit{Variance.} By independence of the $\Delta_k^2$,
$$ \mathrm{Var}(Q_m) = \sum_{k=1}^m \mathrm{Var}(\Delta_k^2). $$
Compute $\mathrm{Var}(\Delta_k^2)$: if $Z \sim \mathcal{N}(0, \sigma^2)$, then
$Z^2 = \sigma^2
\xi$ with $\xi \sim \chi^2_1$. Since $E(\xi^2) = 3$ (fourth moment of the standard normal) and
$E(\xi) = 1$, $\mathrm{Var}(\xi) = 2$. Hence $\mathrm{Var}(Z^2) = \sigma^4 \cdot 2 = 2 \sigma^4$, and
in this case $\mathrm{Var}(\Delta_k^2) = 2 (\Delta t_k)^2$.

Substituting,
$$ \mathrm{Var}(Q_m) = 2 \sum_{k=1}^m (\Delta t_k)^2 \le 2 \delta \sum_{k=1}^m \Delta t_k =
2 \delta (b - a) \xrightarrow{\delta \to 0} 0. $$

\textit{Conclusion.} $E((Q_m - (b - a))^2) = \mathrm{Var}(Q_m) \to 0$, that is $Q_m \to b - a$
in $L^2$.
\end{proof}

\obs Heuristically, this is written as $(dW)^2 = dt$. In the limit one also has
$$ \sum (\Delta t)^2 \to 0, \qquad \sum \Delta t \, \Delta W \to 0, \qquad \sum (\Delta W)^2 \to b
- a, $$
which are the rules used when constructing the Itô integral and the corresponding lemma.

\newpage

\section{Itô Integral}

One would like to define $\int_a^b X_t \, dW_t$ following Riemann: approximate by sums
$\sum X_{\xi_k} (W_{t_{k+1}} - W_{t_k})$ with $\xi_k \in [t_k, t_{k+1}]$. But unlike the
classical case, the choice of $\xi_k$ \emph{does} matter, because the paths of Brownian motion
do not have bounded variation. The choice that gives an integral with good martingale properties
is the \emph{left endpoint}: $\xi_k = t_k$.

\begin{definicion}{Itô Integral}{ito_integral}
Let $(X_t)_{t \in [a, b]}$ be an adapted process with $X_\cdot \in L^2([a, b] \times \Omega)$
(sufficiently integrable).
For a partition $a = t_0 < t_1 < \cdots < t_m = b$ define the \emph{left-endpoint Riemann
sum}
$$ R_m = \sum_{k=1}^m X_{t_{k-1}} (W_{t_k} - W_{t_{k-1}}). $$
The Itô integral is defined as
$$ \int_a^b X \, dW = \lim_{\delta \to 0} R_m, $$
in $L^2(\Omega)$. The limit exists and does not depend on the sequence of partitions.
\end{definicion}

\obs The choice of the left endpoint is essential. Taking $\xi_k = t_k$ (right endpoint) or
$\xi_k = (t_{k-1} + t_k)/2$ (midpoint, Stratonovich) yields \emph{different} integrals.
Itô's is the only one that keeps the martingale property. The difference is
computed directly for $\int_a^b W \, dW$: if $B_m$ uses the right endpoint and $A_m$ the
left one,
$$ B_m - A_m = \sum_k \left( W_{t_k} - W_{t_{k-1}} \right) \left( W_{t_k} - W_{t_{k-1}} \right)
= Q_m \xrightarrow{L^2} b - a \ne 0 $$
by the quadratic variation: both limits exist, but they differ by exactly $b - a$.

\subsection{Properties of the Itô integral}

\begin{proposicion}{Properties}{props_ito}
Let $X, Y$ be adapted processes that are Itô integrable, and $\alpha, \beta \in \mathbb{R}$:
\begin{enumerate}
    \item $\int_a^b (\alpha X + \beta Y) dW = \alpha \int_a^b X \, dW + \beta \int_a^b Y \, dW$.
    \item $\displaystyle E\!\left[\int_a^b X \, dW\right] = 0$.
    \item $\displaystyle E\!\left[\left(\int_a^b X \, dW\right)^2 \right] = E\!\left[\int_a^b X^2
    \, dt\right]$.
    \item $\displaystyle E\!\left[\int_a^b X \, dW \cdot \int_a^b Y \, dW\right] =
    E\!\left[\int_a^b XY \, dt\right]$.
\end{enumerate}
\end{proposicion}

\begin{proof}
(ii) and (iii) are proved at the level of the left-endpoint Riemann sums (the
general case follows by passing to the limit in $L^2$). Write $R = \sum_k X_{t_{k-1}} \Delta W_k$ with
$\Delta W_k = W_{t_k} - W_{t_{k-1}}$ and $\Delta t_k = t_k - t_{k-1}$.

\textit{(ii)} Since $X_{t_{k-1}}$ is $\mathcal{F}_{t_{k-1}}$-measurable (the process is adapted) and
$\Delta W_k$ is independent of $\mathcal{F}_{t_{k-1}}$ with mean $0$,
$$ E(X_{t_{k-1}} \Delta W_k) = E(X_{t_{k-1}}) \, E(\Delta W_k) = 0, $$
and summing, $E(R) = 0$.

\textit{(iii)} Expanding the square,
$$ E(R^2) = \sum_k E\big(X_{t_{k-1}}^2 (\Delta W_k)^2\big) + 2 \sum_{j < k}
E\big(X_{t_{j-1}} \Delta W_j \, X_{t_{k-1}} \Delta W_k\big). $$
In the cross terms with $j < k$, the factor $X_{t_{j-1}} \Delta W_j X_{t_{k-1}}$ is
$\mathcal{F}_{t_{k-1}}$-measurable, and $\Delta W_k$ is independent of it and centered: each
cross term vanishes. In the diagonal ones, $(\Delta W_k)^2$ is independent of
$X_{t_{k-1}}^2$ and $E((\Delta W_k)^2) = \Delta t_k$, whence
$$ E(R^2) = \sum_k E(X_{t_{k-1}}^2) \, \Delta t_k = E\!\left( \sum_k X_{t_{k-1}}^2 \, \Delta
t_k \right) \xrightarrow{\delta \to 0} E\!\left( \int_a^b X^2 \, dt \right). $$
Property (i) is the linearity of the Riemann sums, and (iv) follows from (iii) by
polarization, $XY = \tfrac{1}{4} \left[ (X + Y)^2 - (X - Y)^2 \right]$.
\end{proof}

\obs The Itô isometry is the formal change $(dW)^2 \mapsto dt$ mentioned earlier: the square
of the stochastic integral translates into the classical integral of $X^2$ with respect to time.

\subsection{Itô's Lemma}

The change of variables in stochastic calculus is not the classical one: an additional
term appears due to $(dW)^2 = dt$.

\obs (\emph{Orders of magnitude}.) The reason why the second-order term survives is seen by
comparing, over a partition of $[a, b]$ with mesh $\delta$, the size of the four
differentials that appear when expanding by Taylor:
\begin{align*}
dt &: & \sum_k (t_k - t_{k-1}) &= b - a, \\
(dW)^2 &: & \sum_k (W_{t_k} - W_{t_{k-1}})^2 &\xrightarrow{L^2} b - a
&& \text{(quadratic variation)}, \\
(dt)^2 &: & \sum_k (t_k - t_{k-1})^2 &\le \delta \, (b - a) \longrightarrow 0, \\
dt \, dW &: & \sum_k |W_{t_k} - W_{t_{k-1}}| \, (t_k - t_{k-1}) &\approx \sqrt{\delta} \,
(b - a) \longrightarrow 0,
\end{align*}
where the last line used that $E((W_{t_k} - W_{t_{k-1}})^2) = t_k - t_{k-1}$, so that
the typical increment is $|W_{t_k} - W_{t_{k-1}}| \approx \sqrt{t_k - t_{k-1}} \le \sqrt{\delta}$.
That is: $dt$ and $(dW)^2$ are of the same order and contribute, whereas $(dt)^2$ and $dt \, dW$
vanish in the limit. In classical calculus (where $(dt)^2$ appears instead of $(dW)^2$) the
second-order term disappears and the change of variables keeps a single integral.

\begin{teorema}{Itô's Lemma}{lema_ito}
Let $g \in C^2(\mathbb{R})$. If $Z_t = g(W_t)$,
$$ dZ_t = g'(W_t) \, dW_t + \tfrac{1}{2} g''(W_t) \, dt. $$
In integral form:
$$ g(W_b) - g(W_a) = \int_a^b g'(W_s) \, dW_s + \frac{1}{2} \int_a^b g''(W_s) \, ds. $$
\end{teorema}

\begin{proof}
Fix a partition $a = t_0 < t_1 < \cdots < t_m = b$ with mesh $\delta = \max_k (t_k -
t_{k-1})$. Write $\Delta W_k = W_{t_k} - W_{t_{k-1}}$ and $\Delta t_k = t_k - t_{k-1}$.
By a telescoping sum,
$$ g(W_b) - g(W_a) = \sum_{k=1}^m [g(W_{t_k}) - g(W_{t_{k-1}})]. $$

\textit{Taylor of order 2.} Since $g \in C^2$, for each $k$ there is $\xi_k$ between $W_{t_{k-1}}$ and
$W_{t_k}$ with
$$ g(W_{t_k}) - g(W_{t_{k-1}}) = g'(W_{t_{k-1}}) \Delta W_k + \tfrac{1}{2} g''(\xi_k) (\Delta
W_k)^2. $$
Summing over $k$, the sum splits into two terms:
$$ g(W_b) - g(W_a) = \underbrace{\sum_{k=1}^m g'(W_{t_{k-1}}) \Delta W_k}_{I_m} +
\underbrace{\tfrac{1}{2} \sum_{k=1}^m g''(\xi_k) (\Delta W_k)^2}_{J_m}. $$

\textit{Term $I_m$.} By the construction of the Itô integral as the $L^2$ limit of
left-endpoint Riemann sums,
$$ I_m \xrightarrow{L^2(\Omega)}{} \int_a^b g'(W_s) \, dW_s. $$

\textit{Term $J_m$.} It is split in two:
$$ J_m = \tfrac{1}{2} \sum_k g''(W_{t_{k-1}}) (\Delta W_k)^2 + \tfrac{1}{2} \sum_k [g''(\xi_k) -
g''(W_{t_{k-1}})] (\Delta W_k)^2 =: J_m^{(1)} + J_m^{(2)}. $$

For $J_m^{(1)}$, the quadratic variation is used: given $\epsilon > 0$, fine partitions make
$\sum_k (\Delta W_k)^2 \to b - a$ in $L^2$. Moreover, for a continuous function
$h$, the convergence
$$ \sum_k h(W_{t_{k-1}}) (\Delta W_k)^2 \xrightarrow{L^2}{} \int_a^b h(W_s) \, ds $$
holds. (This result is the analogue of the fundamental theorem for quadratic variation.)
Applying it with $h = g''$,
$$ J_m^{(1)} \xrightarrow{L^2}{} \tfrac{1}{2} \int_a^b g''(W_s) \, ds. $$

For $J_m^{(2)}$: since $g''$ is continuous and the $\xi_k$ lie between $W_{t_{k-1}}$ and $W_{t_k}$ with
$|W_{t_k} - W_{t_{k-1}}| \to 0$ by uniform continuity of Brownian motion on compact sets,
$\max_k |g''(\xi_k) - g''(W_{t_{k-1}})| \to 0$ almost everywhere. Since $\sum (\Delta W_k)^2$ is
bounded in $L^1$ (its expectation is $b - a$),
$$ |J_m^{(2)}| \le \tfrac{1}{2} \max_k |g''(\xi_k) - g''(W_{t_{k-1}})| \cdot \sum_k (\Delta W_k)^2
\xrightarrow{}{} 0. $$

\textit{Conclusion.} Letting $\delta \to 0$ in
$g(W_b) - g(W_a) = I_m + J_m^{(1)} + J_m^{(2)}$, one gets
$$ g(W_b) - g(W_a) = \int_a^b g'(W_s) \, dW_s + \tfrac{1}{2} \int_a^b g''(W_s) \, ds. $$
\end{proof}

\textbf{Version with $t$.} If $g \in C^{1, 2}(\mathbb{R}_{\ge 0} \times \mathbb{R})$ and
$Z_t = g(t, W_t)$,
$$ dZ_t = \frac{\partial g}{\partial t} dt + \frac{\partial g}{\partial x} dW_t + \frac{1}{2}
\frac{\partial^2 g}{\partial x^2} dt. $$

\subsection{Examples}

\begin{ejemplo}{Square of Brownian motion}{cuadrado_W}
Applying Itô to $g(x) = x^2$, $g'(x) = 2x$, $g''(x) = 2$. One gets
$$ W_b^2 - W_a^2 = \int_a^b 2 W_s \, dW_s + \int_a^b 1 \, ds, $$
whence $\int_a^b W \, dW = \tfrac{W_b^2 - W_a^2}{2} - \tfrac{b - a}{2}$. Note the term
$-(b-a)/2$, absent in classical calculus.
\end{ejemplo}

\begin{ejemplo}{Geometric Brownian motion}{browniano_geo}
Solve the stochastic differential equation
$$ dX_t = \mu X_t \, dt + \sigma X_t \, dW_t, \qquad X_0 \text{ given}. $$
Take $Y_t = \ln X_t$ and apply Itô with $g(x) = \ln x$, $g'(x) = 1/x$, $g''(x) = -1/x^2$:
\begin{align*}
d(\ln X_t) &= \frac{1}{X_t} dX_t - \frac{1}{2 X_t^2} (dX_t)^2 \\
&= \frac{1}{X_t} (\mu X_t \, dt + \sigma X_t \, dW_t) - \frac{1}{2 X_t^2} \sigma^2 X_t^2 \, dt \\
&= \left(\mu - \tfrac{\sigma^2}{2}\right) dt + \sigma \, dW_t.
\end{align*}
Integrating, $\ln X_t = \ln X_0 + (\mu - \tfrac{\sigma^2}{2}) t + \sigma W_t$, so that
$$ X_t = X_0 \exp\!\left(\sigma W_t + \left(\mu - \tfrac{\sigma^2}{2}\right) t\right). $$
This is the classical model of the price of an asset in finance (Black-Scholes model).
\end{ejemplo}

\newpage

\subsection{General stochastic processes}

\begin{definicion}{Itô process}{proceso_ito}
An \emph{Itô process} is a process of the form
$$ X_t = X_0 + \int_0^t f(X_s, s) \, ds + \int_0^t g(X_s, s) \, dW_s, $$
written in differential form as
$$ dX_t = f(X_t, t) \, dt + g(X_t, t) \, dW_t. $$
$f$ is called the \emph{drift} and $g$ the \emph{diffusion coefficient}.
\end{definicion}

\begin{teorema}{Generalized Itô}{ito_general}
If $X_t$ is an Itô process as above and $h \in C^{1,2}(\mathbb{R}_{\ge 0} \times \mathbb{R})$,
$M_t = h(t, X_t)$ satisfies
$$ dM_t = \left(\frac{\partial h}{\partial t} + f \frac{\partial h}{\partial x} + \frac{g^2}{2}
\frac{\partial^2 h}{\partial x^2}\right) dt + g \frac{\partial h}{\partial x} \, dW_t. $$
\end{teorema}

\begin{proof}
Same strategy as Itô's lemma. Partition $[0, T]$ with $0 = t_0 < t_1 < \cdots < t_m =
T$ and apply Taylor of order 2 to $h$:
$$ h(t_k, X_{t_k}) - h(t_{k-1}, X_{t_{k-1}}) = h_t \Delta t_k + h_x \Delta X_k + \tfrac{1}{2}
[h_{tt} (\Delta t_k)^2 + 2 h_{tx} \Delta t_k \Delta X_k + h_{xx} (\Delta X_k)^2] + \cdots, $$
with derivatives evaluated at $(t_{k-1}, X_{t_{k-1}})$. Since $\Delta X_k = f \Delta t_k + g \Delta W_k
+ o(\sqrt{\Delta t_k})$,
$$ (\Delta X_k)^2 = g^2 (\Delta W_k)^2 + 2 f g \Delta t_k \Delta W_k + f^2 (\Delta t_k)^2 +
\cdots. $$
The asymptotic rules of quadratic variation give
$$ \sum (\Delta t_k)^2 \to 0, \qquad \sum \Delta t_k \Delta W_k \to 0, \qquad \sum (\Delta W_k)^2
\to \int dt, $$
so that only the term $h_{xx} g^2 (\Delta W)^2 / 2 \to \tfrac{g^2}{2} h_{xx} \, dt$ survives.
Adding the first-order terms, $h_t \, dt + h_x (f \, dt + g \, dW) = (h_t + f h_x) dt + g h_x
\, dW$. Adding the surviving quadratic term of $h_{xx}$, the stated formula
is obtained.
\end{proof}

\newpage

\section{Girsanov's Theorem}

What happens to the distribution of a process if the probability measure is changed?
Girsanov says exactly how a Brownian motion transforms under an absolutely continuous
change of measure.

\subsection{Motivation: change of mean of a normal}

Let $Z \sim \mathcal{N}(0, 1)$ under $P$. It is shown that there is a measure $Q$, equivalent to $P$ (that
is, $P \ll Q \ll P$), such that $Z \sim \mathcal{N}(\mu, 1)$ under $Q$. For this, first compute
$$ E_P(e^{\mu Z}) = \frac{1}{\sqrt{2\pi}} \int_{-\infty}^\infty e^{\mu t} e^{-t^2/2} \, dt =
\frac{1}{\sqrt{2\pi}} \int e^{-(t-\mu)^2/2} dt \cdot e^{\mu^2/2} = e^{\mu^2/2}. $$
Therefore $E_P(e^{\mu Z - \mu^2/2}) = 1$.

\textbf{Change of measure.} For $A \in \mathcal{F}$ define
$$ Q(A) = \int_A e^{\mu Z - \mu^2/2} \, dP. $$
$Q$ is a probability measure (positive, $\sigma$-additive, $Q(\Omega) = 1$). Under $Q$,
the distribution of $Z$ is computed:
\begin{align*}
F_Q(a) = Q(\{Z \le a\}) &= \int_{\{Z \le a\}} e^{\mu Z - \mu^2/2} \, dP \\
&= \frac{1}{\sqrt{2\pi}} \int_{-\infty}^\infty \mathbf{1}_{\{t \le a\}} e^{\mu t - \mu^2/2}
e^{-t^2/2} dt \\
&= \frac{1}{\sqrt{2\pi}} \int_{-\infty}^a e^{-(t - \mu)^2/2} dt,
\end{align*}
that is, $Z \sim \mathcal{N}(\mu, 1)$ under $Q$.

\subsection{Continuous version}

Before the statement, observe why the candidate density is a martingale. Let
$$ Y_t = \int_0^t X_s \, dW_s - \frac{1}{2} \int_0^t X_s^2 \, ds, \qquad M_t = e^{Y_t}, $$
that is $dY_t = X_t \, dW_t - \tfrac{1}{2} X_t^2 \, dt$, with $(dY_t)^2 = X_t^2 \, dt$ by the
rules of quadratic variation. By Itô's lemma with $g(s) = e^s$,
$$ dM_t = e^{Y_t} \, dY_t + \tfrac{1}{2} e^{Y_t} (dY_t)^2 = e^{Y_t} \left( X_t \, dW_t -
\tfrac{1}{2} X_t^2 \, dt + \tfrac{1}{2} X_t^2 \, dt \right) = M_t X_t \, dW_t. $$
The drift term cancels exactly (that is what the $-\tfrac{1}{2} \int_0^t X_s^2
\, ds$ in the definition is for), and one gets $M_t = 1 + \int_0^t M_s X_s \, dW_s$: an Itô integral,
which is a martingale with mean $E(M_t) = 1$ provided the integrand is sufficiently
integrable. The Novikov condition in the statement guarantees exactly that.

\begin{teorema}{Girsanov}{girsanov}
Let $(W_t)_{t \in [0, T]}$ be a Brownian motion under $P$ and let $(X_t)$ be a sufficiently integrable
adapted process. Define
$$ M_t = \exp\!\left( \int_0^t X_s \, dW_s - \frac{1}{2} \int_0^t X_s^2 \, ds \right). $$
Under reasonable conditions (\emph{Novikov's condition}: $E[e^{\tfrac{1}{2} \int_0^T X_s^2 ds}]
< \infty$), $M_t$ is a martingale and $E(M_T) = 1$. Define the measure
$$ Q(A) = \int_A M_T \, dP \qquad (A \in \mathcal{F}_T). $$
Then $Q$ is equivalent to $P$, and the process
$$ \widetilde{W}_t = W_t - \int_0^t X_s \, ds $$
is a Brownian motion under $Q$.
\end{teorema}

\obs In the particular case $X_t \equiv \mu/\sigma$ (constant), taking $W_t' = \sigma W_t$ with
drift $\mu$, Girsanov says that $\widetilde{W}_t = W_t - \mu t/\sigma$ is a Brownian motion under
$Q$ with $dQ/dP = \exp(\tfrac{\mu}{\sigma} W_T - \tfrac{\mu^2}{2\sigma^2} T)$. This justifies the
change of the geometric Black-Scholes model from the real-world measure to the risk-neutral
measure, on which the pricing of derivatives rests.

\obs Girsanov is the continuous version of Radon-Nikodym applied to processes. The derivative
$dQ/dP = M_T$ is called the \emph{Doléans-Dade exponential} of $\int X \, dW$.

\newpage

\section{Application: the Black-Scholes model}

The tools constructed so far (Brownian motion, Itô, Girsanov) have an emblematic
application: the Black-Scholes formula for pricing financial options. It is a case where
the whole apparatus is put to work.

\subsection{Model of the asset}

Consider an \emph{underlying asset} with price $S_t$ that evolves according to a geometric
Brownian motion:
$$ dS_t = \mu S_t \, dt + \sigma S_t \, dW_t, \qquad S_0 \text{ given}, $$
where $\mu \in \mathbb{R}$ is the \emph{expected rate of return} and $\sigma > 0$ the
\emph{volatility}. As seen when solving the SDE with Itô,
$$ S_t = S_0 \exp\!\left(\sigma W_t + \left(\mu - \tfrac{\sigma^2}{2}\right) t\right). $$
Alongside there is a \emph{risk-free asset} (bond) that grows at a constant rate $r$:
$$ B_t = e^{rt}, \qquad dB_t = r B_t \, dt. $$

\textbf{Derivative.} A \emph{derivative} is a contract that pays a function $\Phi(S_T)$ at a
future time $T$. For example, a European \emph{call option} with \emph{strike} $K$ pays
$$ \Phi(S_T) = (S_T - K)^+ = \max(S_T - K, 0). $$
The problem: what is the fair price $C_0$ to pay today for this contract?

\subsection{No-arbitrage argument and risk-neutral measure}

\textbf{Strategy.} If one can build, with $C_0$ dollars today, a portfolio that at $T$
replicates exactly the payoff $\Phi(S_T)$, then $C_0$ is the only price compatible with the
\emph{absence of arbitrage} (no possibility of a sure gain without investment).

\textbf{Reduction to the discounted asset.} Let $\widetilde{S}_t = S_t / B_t = e^{-rt} S_t$, the price
of the asset measured in bonds. By Itô,
$$ d\widetilde{S}_t = (\mu - r) \widetilde{S}_t \, dt + \sigma \widetilde{S}_t \, dW_t. $$
If the drift $(\mu - r) \widetilde{S}_t \, dt$ could be made to disappear, $\widetilde{S}$
would be a
martingale. This is exactly what Girsanov provides: taking
$$ X_t \equiv \frac{\mu - r}{\sigma}, $$
the measure $Q$ with $dQ/dP = \exp(-\tfrac{\mu - r}{\sigma} W_T - \tfrac{1}{2}(\tfrac{\mu-r}{\sigma})^2
T)$
makes
$$ \widetilde{W}_t = W_t + \tfrac{\mu - r}{\sigma} t $$
a Brownian motion under $Q$. Rewriting,
$$ d\widetilde{S}_t = \sigma \widetilde{S}_t \, d\widetilde{W}_t \qquad \text{(under $Q$)}, $$
that is, $\widetilde{S}$ is a martingale under $Q$.

\begin{definicion}{Risk-neutral measure}{Q_neutral}
The \emph{risk-neutral measure} is the unique measure $Q$ equivalent to $P$ such that
$\widetilde{S}_t = e^{-rt} S_t$ is a martingale under $Q$. Under $Q$, the price of the asset evolves
with drift $r$ (not $\mu$):
$$ dS_t = r S_t \, dt + \sigma S_t \, d\widetilde{W}_t, \qquad S_t = S_0 \exp\!\left(\sigma
\widetilde{W}_t + (r - \tfrac{\sigma^2}{2}) t\right). $$
\end{definicion}

\obs The name ``risk-neutral'' is because under $Q$ \emph{every asset} (both $S$ and $B$)
grows, on average, at the same rate $r$: there is no risk premium. The investors' attitude
toward risk, which affects the real-world drift $\mu$, disappears under $Q$.

\subsection{Pricing formula}

\begin{teorema}{Black-Scholes price (general version)}{bs_general}
Under no arbitrage and market completeness, the fair price at $t = 0$ of a derivative with payoff
$\Phi(S_T)$ is
$$ C_0 = e^{-rT} E_Q(\Phi(S_T)). $$
\end{teorema}

\begin{proof}
\textit{Self-financing portfolio.} A portfolio is described by adapted processes $(\alpha_t, \beta_t)$
that indicate how many units of the asset and of the bond are held at each moment. Its value
is
$$ V_t = \alpha_t S_t + \beta_t B_t. $$
The portfolio is \emph{self-financing} if its variation comes only from the price variations:
$$ dV_t = \alpha_t \, dS_t + \beta_t \, dB_t. $$
In terms of the discounted value $\widetilde{V}_t = V_t / B_t$, by Itô applied to
$V_t \cdot e^{-rt}$,
$$ d\widetilde{V}_t = \alpha_t \, d\widetilde{S}_t. $$

\textit{Replication and martingale under $Q$.} Suppose there is a self-financing portfolio that
\emph{replicates} the payoff: $V_T = \Phi(S_T)$. Under $Q$, $\widetilde{S}_t$ is a martingale (by
Girsanov), and then
$$ \widetilde{V}_t - \widetilde{V}_0 = \int_0^t \alpha_s \, d\widetilde{S}_s $$
is an Itô integral with respect to a martingale, therefore $\widetilde{V}_t$ is a martingale
under $Q$ (assuming the necessary integrability). Taking expectations,
$$ \widetilde{V}_0 = E_Q(\widetilde{V}_T) = E_Q(e^{-rT} \Phi(S_T)). $$
And $V_0 = \widetilde{V}_0$ because $B_0 = 1$, which gives $C_0 = V_0 = e^{-rT} E_Q(\Phi(S_T))$.

\textit{Existence of the replicating portfolio.} In the Black-Scholes model (one asset, one bond),
the market is \emph{complete}: by a martingale representation theorem (the discounted
martingale $\widetilde{C}_t = E_Q(e^{-rT} \Phi(S_T) \mid \widetilde{\mathcal{F}}_t)$ is written as an
Itô integral with respect to $\widetilde{W}$, which gives $\alpha_t$), such a portfolio always exists.
The \emph{no-arbitrage} argument implies that its initial price is unique.
\end{proof}

\subsection{Explicit formula for the European call}

Take $\Phi(s) = (s - K)^+$. Under $Q$, $S_T = S_0 \exp(\sigma \widetilde{W}_T + (r -
\sigma^2/2) T)$, with $\widetilde{W}_T \sim \mathcal{N}(0, T)$, that is
$\ln S_T \sim \mathcal{N}(\ln S_0 + (r - \sigma^2/2) T, \ \sigma^2 T)$.

Define
$$ d_1 = \frac{\ln(S_0/K) + (r + \sigma^2/2) T}{\sigma \sqrt{T}}, \qquad
d_2 = d_1 - \sigma \sqrt{T}, $$
and $N(\cdot)$ the distribution function of a standard normal.

\begin{teorema}{Black-Scholes formula}{bs_formula}
The price of a European call option with strike $K$ and maturity $T$ is
$$ C_0 = S_0 \, N(d_1) - K e^{-rT} \, N(d_2). $$
\end{teorema}

\begin{proof}[Proof (direct computation)]
Compute $E_Q((S_T - K)^+)$. Under $Q$ write
$$ S_T = S_0 \, e^{\sigma \sqrt{T} Z + (r - \sigma^2/2) T}, \qquad Z \sim \mathcal{N}(0, 1). $$
The condition $S_T > K$ is equivalent to
$$ Z > \frac{\ln(K/S_0) - (r - \sigma^2/2) T}{\sigma \sqrt{T}} = -d_2. $$
Therefore,
\begin{align*}
E_Q(S_T \mathbf{1}_{S_T > K}) &= S_0 e^{(r - \sigma^2/2) T} \int_{-d_2}^\infty e^{\sigma \sqrt{T} z}
\frac{e^{-z^2/2}}{\sqrt{2\pi}} dz \\
&= S_0 e^{rT} \int_{-d_2}^\infty \frac{e^{-(z - \sigma\sqrt{T})^2/2}}{\sqrt{2\pi}} dz
\qquad \text{(completing the square)} \\
&= S_0 e^{rT} \cdot P(Z' > -d_2 - \sigma\sqrt{T}) \quad \text{with } Z' \sim \mathcal{N}(0,1) \\
&= S_0 e^{rT} \, N(d_2 + \sigma\sqrt{T}) = S_0 e^{rT} \, N(d_1).
\end{align*}
And $E_Q(K \mathbf{1}_{S_T > K}) = K \cdot P(Z > -d_2) = K \, N(d_2)$. Substituting in
$C_0 = e^{-rT} E_Q((S_T - K)^+) = e^{-rT}[E_Q(S_T \mathbf{1}_{S_T > K}) - K E_Q(\mathbf{1}_{S_T >
K})]$,
$$ C_0 = e^{-rT}[S_0 e^{rT} N(d_1) - K N(d_2)] = S_0 N(d_1) - K e^{-rT} N(d_2). $$
\end{proof}

\subsection{Financial interpretation}

\begin{itemize}
    \item $N(d_2)$ is the probability under $Q$ that the option \emph{ends in-the-money}, that
    is, $S_T > K$. The second term $K e^{-rT} N(d_2)$ is the discounted value of paying $K$
    in that case.
    \item $S_0 N(d_1)$ is the \emph{delta hedge}: the amount of underlying (in present value)
    that must be held in the portfolio to replicate the payoff of the call.
    \item The formula \emph{does not depend} on $\mu$, the real drift of the asset. It depends only on
    $\sigma$, $r$, $S_0$, $K$, $T$. This was a historical surprise: the fair price does not require
    estimating how the stock grows, only how much it fluctuates.
\end{itemize}

\obs The model has strong assumptions (constant volatility, no jumps, no transaction
costs, possibility of continuous trading). In practice systematic deviations are observed
(volatility smile, jumps) that motivate more sophisticated models (Heston, models with
jumps, stochastic volatility). But Black-Scholes is still the baseline reference.

\end{document}
